\pdfinclusioncopyfonts=1
\RequirePackage{latex/atlaslatexpath}
\documentclass[NOTE, atlasdraft=true, texlive=2020, UKenglish]{atlasdoc}


\usepackage[subcaption=true]{atlaspackage}
\usepackage{atlasbiblatex}


\usepackage{atlascontribute}

\usepackage{atlasphysics}


\usepackage[output=true]{atlastodo}

\addbibresource{bib/ATLAS.bib}
\addbibresource{bib/ATLAS-useful.bib}
\addbibresource{bib/CMS.bib}
\addbibresource{bib/ConfNotes.bib}
\addbibresource{bib/PubNotes.bib}
\addbibresource{bib/mydocument.bib}

\graphicspath{{logos/}{figures/}}

\usepackage{mydocument-defs}
\usepackage{graphicx}
\usepackage{multirow}
\usepackage{multicol}
\usepackage{xfrac}
\usepackage{wrapfig}
\usepackage[mathscr]{euscript}
\usepackage{csquotes}
\usepackage{float} \usepackage{enumitem}






\newcommand{\NL}{\newline\noindent}
\newcommand{\NLC}{\NL$\circ\,$}

\DeclareRobustCommand{\rchi}{{\mathpalette\irchi\relax}}
\newcommand{\irchi}[2]{\raisebox{\depth}{$#1\chi$}} 

\newcommand{\INT}{\textcolor{violet}{- - - - - INT NOTE - - - - -}}
\newcommand{\PAP}{\textcolor{violet}{- - - - - PAPER - - - - -}}

\newcommand{\TCR}[1]{\textcolor{red}{#1}}
\newcommand{\TCO}[1]{\textcolor{orange}{#1}}
\newcommand{\TCB}[1]{\textcolor{blue}{#1}}
\newcommand{\TCV}[1]{\textcolor{violet}{#1}}
\newcommand{\TCP}[1]{\textcolor{Crimson}{#1}}
\newcommand{\TCC}[1]{\textcolor{cyan}{#1}}
\newcommand{\TCG}[1]{\textcolor{green}{#1}}

\newcommand{\RA}{$\rightarrow$}
\newcommand{\CM}{\TCG{\checkmark}}





\newcommand{\AZDC}{$\langle \text{E}_\text{ZDC}^\text{Pb-side} \rangle$}
\newcommand{\AZDCp}{$\langle \text{E}_\text{ZDC}^\text{p-side} \rangle$}
\newcommand{\EZDC}{$\text{E}_\text{ZDC}^\text{Pb-side}$}
\newcommand{\EZDCp}{$\text{E}_\text{ZDC}^\text{p-side}$}

\newcommand{\EAZDC}{\langle \text{E}_\text{ZDC}^\text{Pb-side} \rangle}
\newcommand{\EAZDCp}{\langle \text{E}_\text{ZDC}^\text{p-side} \rangle}
\newcommand{\EEZDC}{\text{E}_\text{ZDC}^\text{Pb-side}}
\newcommand{\EEZDCp}{\mathrm{E}_\text{ZDC}^\text{p-side}}

\newcommand{\EATPB}{$\langle \text{FCal} \sum \text{E}_{\text{T}}^{\text{Pb-side} } \rangle$}
\newcommand{\EATP}{$\langle \text{FCal} \sum \text{E}_{\text{T}}^{\text{p-side} }\rangle$}
\newcommand{\ETPB}{$\text{FCal} \sum \text{E}_{\text{T}}^{\text{Pb-side} }$}
\newcommand{\ETP}{$\text{FCal} \sum \text{E}_{\text{T}}^{\text{p-side} }$}
\newcommand{\PT}{$p_{\text{T}}$}
\newcommand{\PTE}{p_{\text{T}}}
\newcommand{\PTONE}{$p_{\text{T}_\text{1}}$}
\newcommand{\PTTWO}{$p_{\text{T}_\text{2}}$}
\newcommand{\PTAE}{p_{\text{T,Avg}}}
\newcommand{\PTT}{$p_{\text{T}}^{\text{truth}}$}
\newcommand{\PTTE}{p_{\text{T}}^{\text{truth}}}
\newcommand{\PTR}{$p_{\text{T}}^{\text{reco}}$}
\newcommand{\PTRE}{p_{\text{T}}^{\text{reco}}}
\newcommand{\PTA}{$p_{\text{T,Avg}}$}

\newcommand{\ETA}{$\eta$}
\newcommand{\RAP}{$y$}
\newcommand{\YBOOSTE}{y_{\text{b}}}
\newcommand{\YSTARE}{y^{*}}
\newcommand{\YSTAR}{$y^{*}$}
\newcommand{\YBOOST}{$y_{\text{b}}$}
\newcommand{\AVGYS}{$\langle$\YSTAR$\rangle$}
\newcommand{\AVGYB}{$\langle$\YBOOST$\rangle$}


\newcommand{\RCP}{$R_\text{CP}$}
\newcommand{\RCPZDC}{$R_\text{CP}^\text{ZDC}$}
\newcommand{\RCPE}{R_\text{CP}}

\newcommand{\xp}{$x_{p}$}
\newcommand{\xPb}{$x_{\text{Pb}}$}
\newcommand{\xf}{$x_{\text{F}}$}

\newcommand{\TAB}{\ensuremath{T_{\text{AB}}}}
\newcommand{\TPPB}{$\langle T_{\text{AB}} \rangle$}
\newcommand{\TAA}{$\langle T_{\text{AA}} \rangle$}
\newcommand{\TPPBE}{\langle T_{\text{AB}} \rangle}
\newcommand{\NCOLL}{$\langle N_{\text{coll}} \rangle$}



\newcommand{\NPART}{$\langle N_{\TN{part}} \rangle$}
\newcommand{\SNN}{\ensuremath{\sqrt{s_{_\text{NN}}}}}

\newcommand{\JES}{$\langle \PTRE / \PTTE \rangle$}



\newcommand{\MBT}{\texttt{HLT\_mb\_sptrk}}
\newcommand{\MBTLO}{\texttt{HLT\_mb\_sptrk\_L1MBTS\_1}}


\newcommand\ddfrac[2]{\frac{\displaystyle #1}{\displaystyle #2}}
\newcommand{\pythia}{\mbox{\textsc{Pythia}8}}



\newcommand{\vn}{\mbox{$v_{n}$}}
\newcommand{\vtwo}{\mbox{$v_{\mathrm{2}}$}}
\newcommand{\vthree}{\mbox{$v_{\mathrm{3}}$}}
\newcommand{\vfour}{\mbox{$v_{\mathrm{4}}$}}


\newcommand{\pxp}{$pp$}
\newcommand{\pxA}{$p$+A}
\newcommand{\AxA}{A+A}
\newcommand{\dxAu}{$d$+Au}
\newcommand{\pxPb}{$p$+Pb}
\newcommand{\Pbxp}{Pb+$p$}
\newcommand{\PbxPb}{Pb+Pb}

\newcommand{\DS}{dataset }
\newcommand{\TN}[1]{ \text{#1} }

\newcommand{\NB}{$\textnormal{nb}^{-1}$}
























 







\newif\ifinternalnote 
\internalnotefalse 


\renewcommand{\thesection}{\arabic{section}}
\renewcommand{\theequation}{\arabic{equation}}
\setcounter{secnumdepth}{3}






\AtlasTitle{Characterization of nuclear break-up as a function of hard-scattering kinematics using dijets in $p$+Pb collisions at $\sqrt{S_\mathrm{NN}}$=8.16 TeV with the ATLAS detector}

\AtlasVersion{0.8}

\AtlasAbstract{This note presents the details of the first measurement of forward neutron topologies in proton-lead collisions carried out by ATLAS at a center-of-mass energy of 8.16 TeV. A total luminosity of about 60 nb$^{-1}$, collected in 2016, is used for the measurement. This analysis uses dijets as proxies to the kinematic of the initial state and correlates it with the energy deposited in the Zero Degree Calorimeter (ZDC). The energy deposited in the ATLAS Forward Calorimeter (FCal) is also analyzed as a function of the event kinematics. All the details of the measurement are discussed in this document. 



}

\author[a]{Matthew Hoppesch}
\author[a]{Riccardo Longo}
\author[a]{Aric Tate}
\author[a]{Anne Sickles}

\author[b]{Peter Steinberg}
\author[c]{Brian Cole}
\affil[a]{University of Illinois at Urbana-Champaign}
\affil[b]{Brookhaven National Laboratory}
\affil[c]{Columbia University}







\AtlasNote{ANA-HION-2023-15}
 \hypersetup{pdftitle={ATLAS document},pdfauthor={Matthew Hoppesch, Aric Tate, Riccardo Longo}}

\begin{document}

\maketitle

\tableofcontents



\clearpage{}\section{Note Technical Information} 
\subsection{List of Contributions}
\begin{center}
\begin{tabular}{ l r }
\hline
Matthew Hoppesch \hspace{3cm} & Main analyzer and editor of the note \\
Riccardo Longo & Co-Main analyzer and editor of the note \\
Aric Tate \hspace{3cm} & Co-Main analyzer and editor of the note \\
Anne Sickles & Analysis support and editor of the note \\
Brian Cole & Analysis support and editor of the note \\
Peter Steinberg & Analysis support and editor of the note \\
\hline
\end{tabular}
\end{center}

\subsection{Changelog}
\textbf{v0.1 (2$^{nd}$ July 2024)}
\begin{itemize}
\item First version of the new internal note uploaded to CDS to request Editorial Board. 
\end{itemize}

\textbf{v0.2 (26$^{th}$ July 2024)}
\begin{itemize}
    \item Pile-up subtraction via template fit added to the note and applied to update the results.   
    \item Described 2D unfolding procedure, with inclusion of first \EZDC\ results and comparison with results without unfolding.
    \item Included outline of systematic studies for 2D unfolding, that will be included in the next iteration. 
\end{itemize}

\textbf{v0.3 (8$^{th}$ August 2024)}
\begin{itemize}
    \item Added results with full systematic studies on 2D unfolding. 
\end{itemize}
\textbf{v0.4 (13$^{th}$ August 2024)}
\begin{itemize}
    \item Addressed all comments by the editorial board - mostly related to better explanation of certain passages and additional material provided in the appendix. 
\end{itemize}
\textbf{v0.5 (22$^{nd}$ August 2024)}
\begin{itemize}
    \item Added convergence plots as per request of the HIJet group during sub-group approval review. 
    \item Added ratios to comparison of different \xp\ bins of ZDC and FCal distributions to better render the difference between them. 
    \item Added systematic uncertainty on fit parameters of FCal-ZDC correlation, obtained by varying the fit range as suggested by the EdB. 
    \item Added systematic propagation on the ratio between FCal and ZDC means. 
    \item Added comparison between different \xp\ bins of FCal vs ZDC correlations, with fit. 
\end{itemize}
\textbf{v0.6 (5$^{th}$ September 2024)}
\begin{itemize}
    \item Added eta-pt maps identifying the trigger used as requested by HI Group Reader during group approval meeting.
    \item Added response matrix integrated over bins in \EZDC\ as requested by HI Group Reader during group approval meeting. 
    \item Expanded plot on convergence of unfolding iterations to the first iteration as requested by HI Group Reader during group approval meeting.
    \item Changed systematic uncertainties on \EZDC\ and \ETPB\ to properly represent the uncertainty on self normalized distributions. 
    \item Shifted raw distributions of \AZDC\ and \EATPB\ wrt the unfolded distributions as requested by HI Group Reader during group approval meeting.
    \item Changed several results plots to reflect style using in CONF. 
    \item Modified text in places where typos were pointed out or where further clarification was needed as requested by HI Group Reader during group approval meeting.
\end{itemize}
\textbf{v0.7 (18$^{th}$ September 2024)}
\begin{itemize}
    \item Added correlation plot between \EZDC\ and \ETPB\ in both dijet events and minimum bias events.  These plots are binned in the same binning as the CONF. 
    \item Added normalized distributions of \EZDC\ and \ETPB\ in minimum bias events.
\end{itemize}
\textbf{v0.8 (2$^{nd}$ February 2025)}
First note version after the CONF - and toward the final paper. We try to reference each new section here - such that you can find easily the new additons. 
\begin{itemize}
    \item Added a description of the new cut to reject events affected by Out Of Time Pile-Up in the FCal by using the ZDC - see Section~\ref{sec:oot_zdc}.  
    \item Added section on the study of the proton going energy deposited in the FCal - see Section~\ref{sec:pgoing_fcal}.
    \item Added minimum bias event distributions as per request of theorist who have manifested interest in our preliminary results - see Section~\ref{sec:mb_dist}.
\end{itemize}



%
\clearpage{}
\clearpage{}\section{Introduction}
\label{sec:intro}
Relativistic proton-nucleus (\pxA) collisions provide unique experimental access to a wide array of physics effects. In the recent past, proton-lead (\pxPb) data collected at the Large Hadron Collider (LHC) have contributed to advance the understanding of hot and cold nuclear matter, thanks to several measurements carried out by LHC experiments. 
Initial state effects introduced by the nuclear environment, resulting in modifications of parton distribution functions (PDFs), have been addressed by measurements of the W \cite{CMS_W_2015565} and Z \cite{ATLAS:2015mwq} bosons by CMS and ATLAS, respectively. The measurement of high transverse momentum (\PT{}) probes, such as charged hadrons or jets, has also contributed to the advancement of nuclear PDFs (nPDFs) parametrizations ~\cite{Rabah::nPDF2021,Eskola2019,Eskola2020,PAUKKUNEN2017241, Helenius:2022rmf}. CMS measured dijets in \pxPb\ collisions at 5.02 TeV \cite{CMS-HIN-13-001}, while ATLAS carried out analyses of the dijet production in ultra-peripheral collisions (UPCs) \cite{ATLAS:2022cbd}. 
Recently, ATLAS also published results for the centrality dependence of the dijet production in \pxPb\ collisions at 8.16 TeV \cite{ATLAS:2023zfx}, showing a clear event-activity bias driven by the kinematic of the initial state. These results shed further light on color fluctuation effects \cite{Mueller:1982, Brodsky:1982, PhysRevLett.47.297} in \pxA\ collisions, providing new input to their modeling after a first measurement of inclusive jet production in \pxPb\ at 5.02 TeV carried out using Run 1 data \cite{ATLAS:2014cpa}. 

The analysis of dijets detected in \pxPb\ collisions at \SNN, = 8.16 TeV (high-luminosity sample, $\sim$165~nb$^{-1}$ considering both beam orientations) using the full acceptance of the ATLAS calorimeter can further help in understanding the nuclear environment and its dynamics. Recently, ATLAS published a characterization of the event activity-classified dijet production on the kinematic of the initial state, showing how the proton configuration at the time of the \pxPb\ collision is indisputably driving the production of forward energy released by the nuclear debris \cite{ATLAS:2023zfx}.  An ongoing analysis of dijet production cross-section over the full trigger and calorimeter acceptance during the 2016 \pxPb\ run will provide unprecedented input to the constraining of nPDFs, while a recent calibration of the Zero Degree Calorimeter performance during the 2016 \pxPb\ run \cite{Hoppesch:2900504} has enabled a number of new studies using this unique dataset, which is likely to be the last high-luminosity \pxPb\ sample until the HL-LHC (2029++). 

The \xp\ dependence of the energy production in \pxA\ collisions at different rapidities is one of the most intriguing manifestations of color fluctuations in the size and interaction strength of the proton interacting with the nucleus. The relation between the energy deposited in the ATLAS Forward Calorimeter (FCal) and \xp\ was investigated in \pxPb\ \cite{ATLAS:2014cpa} and \pxp\ \cite{HION-2015-02}. The study presented in this Note reports a comprehensive characterization of the dependence of the energy production on \xp, including not only the Pb-facing side of the FCal, but both sides of the calorimeter, as well as the ZDC. This characterization will serve both the modeling of color fluctuation effects in the event activity characterizing \pxPb\ collision and of nuclear debris originating from the evaporation of the ion target. 

The nuclear breakup after the \pxA\ collision and the associated forward neutron production represent two of the least known aspects toward a comprehensive understanding of \pxA\ reactions. A few forward neutron/energy measurements, carried out in Minimum Bias (MB) events with the ZDC facing the nuclear debris, have been published by ALICE \cite{ALICE:2021poe} and CMS \cite{Suranyi:2021ssd}. On the other hand, there is a lack of data to model the nuclear breakup and the production of forward neutrons in \pxPb\ collisions characterized by the presence of a hard scattering. This note reports the first analysis of forward neutral energy detected in the ZDC in  \pxPb\ reactions, examining both MB collisions as well as events characterized by the presence of a dijet pair originating from a hard-scattering. This second class of event is particularly interesting to correlate the production of forward energy with the kinematics of the initial state, and can help the advancement of recent models that could connect the hard-scattering with the dynamics at play during the nuclear breakup. 
The idea that forward neutron topologies could be correlated to the number of wounded nucleons in photon-nucleus interactions was recently proposed for UPCs in \cite{Alvioli:2024cmd} as a way to constrain the parameterization of nuclear shadowing in nPDFs. The paper focuses on real photons and explores the importance of the framework used to describe hadronic fluctuations of the real photon, and their subsequent effect on the interaction with the nucleus and the number of wounded nucleons. The authors are currently working towards including this modeling also in \pxPb\ reactions. The use of dijets serves, also in this case, as a key proxy to access the initial state kinematics and, therefore, the proton configuration at the moment of the \pxPb\ collision. 
Therefore, the results included in this note could also be used to study the interplay between the kinematics of the initial state, color fluctuations, number of wounded nucleons and nuclear breakup. 



















\subsection{General Analysis Strategy}
\label{sec:GAS}
The analysis presented in this note involves the first investigation of forward neutron topologies in \pxPb\ collisions using the ATLAS detector and is enabled by the recent calibration of the Zero Degree Calorimeter (ZDC) in the 2016 8.16 \pxPb\ dataset \cite{Hoppesch:2900504}. The energy deposited in the ATLAS FCal is also analyzed, to provide a comprehensive overview of the energy production in \pxPb\ collisions. 

This analysis builds upon a previous ATLAS measurement that analyzed the centrality dependence of the dijet production via dijet per-event yields (see the corresponding Internal Note ~\cite{Longo:2860813}). 
In this recent paper, the three kinematic variables defined were the average transverse momentum, \PTA, boost of the dijet system, \YBOOST, and half rapidity separation, \YSTAR, defined as:
\begin{equation}
\label{eq:kinevar_definition} 
\PTAE = \frac{p_{\TN{T,1}}+p_{\TN{T,2}}}{2}, \, \quad y_{\TN{b}} = \frac{ y_1^{\TN{CM}}+y_2^{\TN{CM}} }{2}, \, \quad y^{*} = \frac{| y_1^\TN{CM}-y_2^\TN{CM} | }{2} .
\end{equation}
Please note that the superscript CM indicates that the variables are computed in the nucleon-nucleon center-of-mass reference frame, and the subscripts `1' and `2' refer to the leading and subleading jet, respectively. 

In the case of \pxPb, the CM value for the rapidity can be obtained by adding/subtracting $\Delta y = \pm 0.465$, as discussed in detail in Section~\ref{sec:data}.
This value arises from the fact that
 the proton beam has an energy of 6.5 TeV, while the Pb beam has an energy of 2.56 per nucleon in the laboratory frame.



In this work, the kinematics of the detected leading (\PT{$_{,1}$} and $y^\TN{CM}_{1}$) and sub-leading (\PT{$_{,2}$} and $y^\TN{CM}_{2}$) jets are used, together with the center of mass energy, to access the momentum fraction carried by the parton participating in the hard scattering which originate from the colliding nucleons. 

The chosen kinematic variables allow for full access to the parton level kinematics:
\begin{equation}
\label{eq:xp}
x_p = 
\frac{ p_{\TN{T},1}e^{y^\TN{CM}_{1}} 
+ 
p_{\TN{T},2}e^{y^\TN{CM}_{2}}}
{\sqrt{s}}
\simeq
\frac{2p_{\TN{T},\TN{Avg}}}
{\sqrt{s}}
e^{y_b}
\cosh(y^{*
}),
\end{equation}
\begin{equation}
\label{eq:xPb}
x_{\TN{Pb}} 
= 
\frac{ p_{\TN{T},1}e^{-y^\TN{CM}_{1}} 
+ 
p_{\TN{T},2}e^{-y^\TN{CM}_{2}}}
{\sqrt{s}} 
\simeq
\frac{2p_{\TN{T},\TN{Avg}}}
{\sqrt{s}}
e^{-y_b}
\cosh(y^{*}),
\end{equation}
\begin{equation}
\label{eq:xF}
x_\mathrm{F} = x_p-x_{\TN{Pb}}
\end{equation}



where \xp\,and \xPb\,are the momentum fractions carried by the partons participating in the hard scattering, which originate from the colliding proton and the lead nucleus, respectively, and \xf\ is Feynman's variable. 
Please note that the $\simeq$ symbol in these formulas is meant to remind the reader that the kinematics reconstructed using dijet variables represent an approximation of the one reconstructed using the two leading and sub-leading jet kinematics. The implications of this approximation were studied in a previous analysis by analyzing dijet events generated with Pythia 8, see the Appendix of \cite{Longo:2860813}, as well as Section~\ref{sec:xp_res} of this note.
This analysis mainly utilizes the left half of Equations \ref{eq:xp} and \ref{eq:xPb}. 
The \xf\ variable is built starting from \xp\ and \xPb\ and therefore its level of approximation follows from the choice mentioned above. The definition of the parton kinematics as represented on the right side of Equations \ref{eq:xp} and \ref{eq:xPb} may occasionally be employed in the note to provide input directly comparable to the recently published dijet results. 





\subsection{Measured Observables}
\label{sec:meas_obs}

This analysis aims at measuring correlations between neutron topologies detected in the ZDC and initial state kinematics. Particular focus is given to the ZDC receiving the neutrons from the Pb nucleus evaporation where the energy deposited in the calorimeter, \EZDC, can be measured either in TeV or in terms of number of neutrons using the ZDC calibration described in \cite{Hoppesch:2900504}. This variable is measured in dijet events as a function of \xp, reconstructed according to Equation \ref{eq:xp}, left side. In addition, a broader analysis of the energy production in the same class of events is carried out by looking at the energy deposited in the $p$-going side of the ZDC, \EZDCp, and in the $p$ and Pb-going arms of the FCal, \ETP\ and \ETPB, respectively. The analysis of the energy production in the forward and backward regions is carried out to better understand the dynamics of dijet events initiated by partons with different \xp, and further inform the modeling of hard-scatterings in \pxA. 

In addition to this core analysis, the ratio between \EZDC\ in central and peripheral dijet events
\begin{equation}
    R_\text{CP}^\text{ZDC} = \frac{\langle \EEZDC \rangle_{0-10\%} }{\langle \EEZDC \rangle_{60-90\%}}
\end{equation}

using the same definition of centrality adopted in \cite{ATLAS:2023zfx} (see Section~\ref{sec:centrality} for more details), is also reported as a function of the initial state kinematics. The ultimate goal of this part of the analysis is to provide further understanding of collisions involving small proton spatial configurations by looking, for the first time, at the neutral energy deposited in the ZDC. 

It should be noted that the 2016 \pxPb\ dataset at \SNN = 8.16 TeV includes data collected with two opposite beam orientations. The Zero Degree Calorimeter (ZDC) configuration was different between the two arms, with the LHCf detector installed instead of the ZDC Electromagnetic (EM) section on Arm C. The consequences of this asymmetry, and its selecting the data for this analysis, are described in Section \ref{sec:data}. At the time of EB formation, after an analysis of both beam orientations and of the implications of the LHCf presence in place of the ZDC EM module, the analysis team decided to focus this analysis only on the data where the Pb debris are traveling toward the full ZDC configuration on Arm A, denoted in the following as \pxPb. Distributions obtained from the analysis of the other beam orientation, \Pbxp, are reported in the Supporting Material for completeness, see Sections ~\ref{auxsec:both_periods} and \ref{auxsec:both_periods_RCP}. 


\newpage
\subsection{Binning}
\label{sec:binning}
The analysis is carried out as a function of the initial state kinematic variables, \xp\ and \xf.  For these variables, the binning was chosen as follows: 
\begin{itemize}
    \item \textbf{\xp}: Logarithmic binning, defined as: \{ 0.00063, 0.00131, 0.00191, 0.00275, 0.00398, 0.00575, 0.00832, 0.01202, 0.01738, 0.02512, 0.03631, 0.05248, 0.07586, 0.10965, 0.15849, 0.22909, 0.33113, 0.47863, 1. \}.
    \item \textbf{\xf}: Linear binning, defined as: \{ -1., -0.9, -0.8, -0.7, -0.6, -0.5, -0.4, -0.3, -0.2, -0.1, 0 , 0.1, 0.2, 0.3, 0.4, 0.5, 0.6, 0.7, 0.8, 0.9, 1.0 \}
\end{itemize}
A linear binning was used to analyze the energy in the Pb and $p$-going sides of the ZDC:
\begin{itemize}
\item \textbf{\EZDC\ (in TeV)}: \{ 0.0, 1.0, 20.2, 39.4, 58.6, 77.8, 97.0, 116.2, 135.4, 154.6, 173.8, 193.0, 212.2, 231.4, 250.6, 269.8, 289.0, 308.2, 327.4, 346.6, 365.8, 385.0 \}. 
\item \textbf{\EZDCp\ (in TeV)}: \{ 1, 2, 3, 4, 5, 6, 7, 8, 9, 10, 11, 12, 13, 14, 15, 16, 17, 18, 19, 20 \}. 
\end{itemize}
A logarithmic binning is used for both the energies accumulated in the FCal arms:
\begin{itemize}
\item \textbf{\ETPB\ and \ETP\ (in TeV)}: \{ 0.0010, 0.0014, 0.0020, 0.0028, 0.0040, 0.0056, 0.0079, 0.0112, 0.0158, 0.0224, 0.0316, 0.0447, 0.0631, 0.0891, 0.1259, 0.1778, 0.2512, 0.3548, 0.5012, 0.7079, 1.0000 \} 
\end{itemize}\clearpage{}
\clearpage{}\section{Experimental Setup}
\label{sec:setup}
The ATLAS detector \cite{ATLAS_Collaboration_2008} at the LHC covers nearly the entire solid angle around the collision point.\footnote{ATLAS uses a right-handed coordinate system with its origin at the nominal interaction point (IP) in the center of the detector and the $z$-axis along the beam pipe. The $x$-axis points from the IP to the center of the LHC ring, and the $y$-axis points upwards. Cylindrical coordinates ($\rho$,$\phi$) are used in the transverse plane, $\phi$ being the azimuthal angle around the $z$-axis. The psuedorapidity is defined in terms of the polar angle $\theta$ as $\eta$ = -$\ln$~tan($\theta$/2). Angular distance is measured in units of $\Delta R\equiv \sqrt{ (\Delta \eta)^2 + (\Delta \phi)^2 } $.} It consists of an inner tracking detector surrounded by a thin super-conducting solenoid, electromagnetic and hadronic calorimeters, and a muon spectrometer incorporating three large super-conducting toroidal magnets. The inner-detector system (ID) is immersed in a 2 T axial magnetic field and provides charged particle tracking in the range $|\eta|<2.5$.

The high-granularity silicon pixel detector covers the vertex region and typically provides four measurements per track, the first hit being normally in the innermost layer. It is followed by the silicon microstrip tracker which usually provides four two-dimensional measurement points per track. These silicon detectors are complemented by the transition radiation tracker, which enables radially extended track reconstruction up to $|\eta|=2.0$. The transition radiation tracker also provides electron identification information based on the fraction of hits (typically 30 in total) above a higher energy deposit threshold corresponding to transition radiation.

 The calorimeter system covers the psuedorapidity range $|\eta|<4.9$. Within the region $|\eta|<3.2$, electromagnetic calorimetry is provided by barrel and endcap high-granularity lead/liquid-argon (LAr) electromagnetic calorimeters, with an additional thin LAr presampler covering $|\eta|<1.8$, to correct for energy loss in material upstream of the calorimeters. Hadronic calorimetry is provided by the steel/scintillating-tile calorimeter, segmented into three barrel structures within $|\eta|<1.7$, and two copper/LAr hadronic endcap calorimeters. The solid angle coverage is completed with forward copper/LAr and tungsten/LAr calorimeter modules optimized for electromagnetic and hadronic measurements respectively.

  
 The muon spectrometer (MS) comprises separate trigger and high-precision tracking chambers measuring the deflection of muons in a magnetic field generated by superconducting air-core toroids. The field integral of the toroids ranges between 2.0 and 6.0~T~m across most of the detector. A set of precision chambers covers the region $|\eta|<2.7$ with three layers of monitored drift tubes, complemented by cathode strip chambers in the forward region, where the background is highest. The muon trigger system covers the range $|\eta|<2.4$ with resistive plate chambers in the barrel, and thin gap chambers in the endcap regions. 
The ATLAS trigger system is used to record events of interest, and is discussed in Section~\ref{sec:trig_evsel}.

The ATLAS Zero Degree Calorimeter (ZDC) consists of two detectors orientated symmetrically around the interaction point (IP) on the beamline axis that measure neutral particles emitted at very small rapidity separation from the incident nuclei. The ZDC covers pseudorapidity  $|\eta|>8.3$.  Each side (or arm) consists of four modules each with a material budget of 1.14 hadronic interaction lengths.  The modules are made of layers of tungsten plates with quartz rods interspersed between, see Figure~\ref{fig:zdc_setup}.

\begin{figure}[ht]
    \centering
    \includegraphics[width=0.95\linewidth]{sections/Setup/fig/ZDCInt.001.png}
    \caption{(Left): 3D view of the ATLAS ZDC Calorimeter, as setup for the 2016 \pxPb\ run in Arm A. Please note that the BRAN detector depicted corresponds to the Run 3 design (same dimension, but increased material budget). (Center): Side view of an ATLAS ZDC module on Arm A. On this side of the detector, Electromagnetic (EM) and Hadronic (HAD) modules have essentially the same design, except for HAD1 that is instrumented with additional transverse pixels (not anymore functional) that are readout in a dedicated back-end about 6 cm deep. (Right): Top view of the ZDC module active area. }
    \label{fig:zdc_setup}
\end{figure}\clearpage{}
\clearpage{}\section{Data Sets}
\label{sec:data}
\subsection{2016 \texorpdfstring{\pxPb}{p+Pb} Data}
The \pxPb\ data at \SNN~=~8.16~\TeV\ used for this analysis were delivered by the LHC to ATLAS in November and December of 2016. In \pxPb\ collisions at 8.16 \TeV\ the colliding beams have different energies per nucleon (6.5~\TeV\ protons colliding with (Z/A) $\times$ 6.5 \TeV\ $\sim$ 2.6~\TeV\ per nucleon in a lead nucleus), resulting in a collision system with \SNN~=~8.16~\TeV\ per nucleon pair. In these conditions, the center of
mass frame is shifted in rapidity, \RAP, by $\Delta$\RAP\ = 0.465 in the direction of the proton beam. Data were collected in two different periods, characterized by opposite beam orientations. The runs in these periods are organized  as follows:
\begin{itemize}
\item \textbf{Period B:} 55.9 \NB\ of total luminosity collected. Beam 1: protons. Beam 2: Pb ions. Protons travels towards negative rapidities, ``A$\rightarrow$C'' direction, corresponding to clockwise motion in the accelerator (according to ATLAS conventions). In the following, this will be denoted as \pxPb\ configuration. 
\item \textbf{Period C:} 104.2 \NB\ of total luminosity collected. Beam 1: Pb ions. Beam 2: protons. Protons travels towards positive rapidities, ``C$\rightarrow$A'' direction, corresponding to counterclockwise motion in the accelerator (according to ATLAS conventions). In the following, this will be denoted as \Pbxp\ configuration. 
\end{itemize}
All the data sets used in this analysis, listed in Table~\ref{tab:datasets}, were processed using ATHENA release 21. The trigger strategy is discussed in Section~\ref{sec:trig_evsel}.
The physics results for this analysis are based only on Period B (\pxPb\ orientation), for reasons discussed in Section~\ref{sec:missing_zdc_mod}. The event selection is discussed in Section~\ref{sec:event_selec}.
\begin{table}[ht]
\centering
\begin{tabular}{ | c | c | c |}
\hline
\rule{0pt}{2.25ex}  
\textbf{2016 \SNN\ = 8.16 \TeV\ Data Samples} & \textbf{Period} & \textcolor{black}{ \textbf{Luminosity [nb$^{-1}$] } } \\
\hline
\hline
data16\_hip8TeV.00313063.physics\_Main.merge.AOD.r11416\_p3868 & B & 0.021 \\
\hline
data16\_hip8TeV.00313067.physics\_Main.merge.AOD.r11416\_p3868 & B & 0.486 \\ 
\hline
data16\_hip8TeV.00313100.physics\_Main.merge.AOD.r11416\_p3868 & B & 8.787 \\ 
\hline
data16\_hip8TeV.00313107.physics\_Main.merge.AOD.r11416\_p3868 & B & 11.053 \\
\hline
data16\_hip8TeV.00313136.physics\_Main.merge.AOD.r11416\_p3868 & B & 9.602 \\
\hline
data16\_hip8TeV.00313187.physics\_Main.merge.AOD.r11416\_p3868 & B & 3.030 \\
\hline
data16\_hip8TeV.00313259.physics\_Main.merge.AOD.r11416\_p3868 & B & 4.659 \\
\hline
data16\_hip8TeV.00313285.physics\_Main.merge.AOD.r11416\_p3868 & B & 4.272 \\
\hline
data16\_hip8TeV.00313295.physics\_Main.merge.AOD.r11416\_p3868 & B & 9.923 \\ 
\hline
data16\_hip8TeV.00313333.physics\_Main.merge.AOD.r11416\_p3868 & B & 3.791 \\
\hline
data16\_hip8TeV.00313435.physics\_Main.merge.AOD.r11416\_p3868 & B & 0.301 \\
\hline
\hline
data16\_hip8TeV.00313574.physics\_Main.merge.AOD.r11416\_p3868 & C & 1.187 \\
\hline
data16\_hip8TeV.00313575.physics\_Main.merge.AOD.r11416\_p3868 & C & 7.168 \\ 
\hline
data16\_hip8TeV.00313603.physics\_Main.merge.AOD.r11416\_p3868 & C & 8.283 \\ 
\hline
data16\_hip8TeV.00313629.physics\_Main.merge.AOD.r11416\_p3868 & C & 6.347 \\
\hline
data16\_hip8TeV.00313630.physics\_Main.merge.AOD.r11416\_p3868 & C & 6.734 \\
\hline
data16\_hip8TeV.00313688.physics\_Main.merge.AOD.r11416\_p3868 & C & 7.475 \\
\hline
data16\_hip8TeV.00313695.physics\_Main.merge.AOD.r11416\_p3868 & C & 4.203 \\ 
\hline
data16\_hip8TeV.00313833.physics\_Main.merge.AOD.r11416\_p3868 & C & 4.769 \\ 
\hline
data16\_hip8TeV.00313935.physics\_Main.merge.AOD.r11416\_p3868 & C & 10.368 \\ 
\hline
data16\_hip8TeV.00313984.physics\_Main.merge.AOD.r11416\_p3868 & C & 2.232 \\ 
\hline
data16\_hip8TeV.00314014.physics\_Main.merge.AOD.r11416\_p3868 & C & 5.954 \\ 
\hline
data16\_hip8TeV.00314077.physics\_Main.merge.AOD.r11416\_p3868 & C & 9.642 \\ 
\hline
data16\_hip8TeV.00314105.physics\_Main.merge.AOD.r11416\_p3868 & C & 6.157 \\ 
\hline
data16\_hip8TeV.00314112.physics\_Main.merge.AOD.r11416\_p3868 & C & 9.947 \\ 
\hline
data16\_hip8TeV.00314157.physics\_Main.merge.AOD.r11416\_p3868 & C & 9.078 \\ 
\hline
data16\_hip8TeV.00314170.physics\_Main.merge.AOD.r11416\_p3868 & C & 4.665 \\ 
\hline
\end{tabular}
\caption{\textcolor{black}{Data samples from \SNN\,= 8.16 \TeV\,\pxPb\,collisions collected during the 2016 Heavy Ion Run. The luminosity collected for each run is reported, together with the corresponding data taking period. All the data were processed using ATHENA Release 21. In this analysis, only data from period B are currently used to produce the final results, see Section~\ref{sec:missing_zdc_mod} for more details on this choice. Period C data are analyzed but currently only included in the supporting material.}}
\label{tab:datasets}
\end{table}



\subsection{Monte Carlo}
\label{sec:mc_samples}

In order to evaluate the performance of jet reconstruction and the jet response within ATLAS, dijet Monte Carlo samples generated using Pythia8~\cite{Sjostrand:2014zea} with the A14 tune and the NNPDF23LO PDFs~\cite{BALL2013244} were produced. Dijet events are filtered using the JZ$*$R04 filters, which select events containing an \antikt\ R~=~0.4 jet in specific truth \PT{} ranges, refered to using the prefix `JZ' and an integer ranging from 1-5 in this analysis. 
The detector response was simulated using the full ATLAS detector simulation based on the Geant4 toolkit~\cite{AGOSTINELLI2003250}. For \pxPb\ collisions, Pythia8 truth events are overlaid with real minimum-bias \pxPb\ data, with separate MC samples for Period B and Period C. For each event, the Pythia generation and subsequent simulations were run with conditions matching the collected data, allowing for the simulation to account for underlying event effects in the jet reconstruction. Please note that, in this MC sample, the event activity is originated mainly from the overlay contribution, and therefore has no truth associated. This aspect is relevant for the discussion on the unfolding strategy that will take place in Section~\ref{ssec:unf_strat}. 

Tables~\ref{tab:pPbMC}~and~\ref{tab:pPbMCFF} list the \PT{} range, Pythia8 cross-section and filter efficiency, and the number of events for the inclusive and forward-filtered \pxPb\ and  \Pbxp\ MC samples, respectively\footnote{Please note that \Pbxp\ samples are reported since in this note we have included also studies using \Pbxp\ data. Still, it is important to note that \Pbxp\ data are not part of the final CONF analysis.}
For forward-filtered samples a specific filter, characterized by its own efficiency, is applied in order to preferentially select events in which a jet exists at forward $\eta$, where the cross-section is lower compared to central jets. This allows for smaller MC productions to provide increased statistics for events characterized by dijets in the forward region, i.e. low $x_{\TN{Pb}}$ and high $x_p$, allowing for a satisfactory MC population over the full phase space considered. Please note that, for the forward-filtered samples generated for this analysis, only jets in the proton-going direction are considered forward (e.g. at least one jet with $\eta$ > 3, if the proton travels from negative to positive rapidities, is required to pass the filter) given the asymmetry of the collision system (instead of considering generally $|\eta| > 3.0$ as done for standard forward filtered \pxp\ or \PbxPb\ samples). 





\begin{table}[!ht]
\centering
\begin{tabular}{ | c | l | c | c | c |}
           \hline
           \rule{0pt}{2.25ex}  
           JZ   &  \hspace{2mm} $\PTE^{\TN{truth}}$ [GeV] & $\sigma$ [nb] & $\epsilon$ &  $N_{\TN{evt}}$ \\[0.25ex]
           \hline
           \hline
           \rule{0pt}{2.25ex} 
            1    & 20 - 60      &   7.3070x10$^7$  &  5.8782x10$^{-3}$ & 4M \\[0.1ex]
            \hline
            \rule{0pt}{2.25ex} 
            2    & 60 - 160     &  1.2916x10$^6$   &  5.6911x10$^{-3}$ & 4M \\[0.1ex]
            \hline
            \rule{0pt}{2.25ex} 
            3    & 160 - 400    &  1.1775x10$^4$   &  7.4572x10$^{-3}$ & 4M \\[0.1ex]
            \hline
            \rule{0pt}{2.25ex} 
            4    & 400 - 800    &  8.9941x10$^1$   &  8.1320x10$^{-3}$ & 4M \\[0.1ex]
            \hline
            \rule{0pt}{2.25ex} 
            5    & 800+         &  1.1951x10$^0$   &  7.1577x10$^{-3}$ & 4M \\[0.1ex]
            \hline               
        \end{tabular}
        \vspace{2mm}
        \caption{Summary of inclusive \SNN\ = 8.16 \TeV\ \pxPb\,and \Pbxp\ Pythia8 MC samples.}
        \label{tab:pPbMC}
\end{table}

\begin{table}[h]
\centering
\begin{tabular}{ | c | l | c | c | c | }
           \hline
           \rule{0pt}{2.25ex}  
           JZ   &  \hspace{2mm} $\PTE^{\TN{truth}}$ [GeV] & $\sigma$ [nb] & $\epsilon$ & $N_{\TN{evt}}$ \\[0.25ex]
           \hline
           \hline
           \rule{0pt}{2.25ex} 
            1    & 20 - 60      &   7.3070x10$^7$  &  9.0671x10$^{-4}$ & 4M \\[0.1ex]
            \hline
            \rule{0pt}{2.25ex} 
            2    & 60 - 160     &  1.2916x10$^6$   &  5.3628x10$^{-4}$ & 2M \\[0.1ex]
            \hline               
        \end{tabular}
        \vspace{2mm}
        \caption{Summary of forward-filtered \SNN\ = 8.16 \TeV\ \pxPb\ and \Pbxp\ Pythia8 MC samples. } 
        \label{tab:pPbMCFF}
\end{table}

The ATLAS data-set names for the MC samples used in this analysis are listed below.
\vspace{1.25cm}
\begin{quote}
\SNN\ = 8.16 \TeV\ \Pbxp\ Inclusive\\
\noindent\rule{12cm}{0.4pt}\\
{\scriptsize
    \texttt{mc16\_pPb8TeV.420011.Pythia8EvtGen\_A14NNPDF23LO\_jetjet\_JZ1R04.merge.AOD.e6519\_d1604$*$}\\
    \texttt{mc16\_pPb8TeV.420012.Pythia8EvtGen\_A14NNPDF23LO\_jetjet\_JZ2R04.merge.AOD.e6395\_d1604$*$}\\
    \texttt{mc16\_pPb8TeV.420013.Pythia8EvtGen\_A14NNPDF23LO\_jetjet\_JZ3R04.merge.AOD.e6519\_d1604$*$}\\
    \texttt{mc16\_pPb8TeV.420014.Pythia8EvtGen\_A14NNPDF23LO\_jetjet\_JZ4R04.merge.AOD.e6519\_d1604$*$}\\
    \texttt{mc16\_pPb8TeV.420015.Pythia8EvtGen\_A14NNPDF23LO\_jetjet\_JZ5R04.merge.AOD.e6519\_d1604$*$}\\
}
\end{quote}
\begin{quote}
\SNN\,= 8.16 \TeV\ \Pbxp\, Forward-filtered\\
\noindent\rule{12cm}{0.4pt}\\
{\scriptsize                    
    \texttt{mc15\_pPb8TeV.420048.Pythia8EvtGen\_A14NNPDF23LO\_jetjet\_JZ1R04\_MaxEta\_3p0.merge.AOD.e8310\_d1604$*$}\\ \texttt{mc16\_pPb8TeV.420049.Pythia8EvtGen\_A14NNPDF23LO\_jetjet\_JZ2R04\_MaxEta\_3p0.merge.AOD.e8310\_d1604$*$}\\
}
\end{quote}

\begin{quote}
\SNN\,= 8.16 \TeV\ \pxPb\ Inclusive\\
\noindent\rule{12cm}{0.4pt}\\
{\scriptsize
    \texttt{mc16\_pPb8TeV.420011.Pythia8EvtGen\_A14NNPDF23LO\_jetjet\_JZ1R04.merge.AOD.e6518\_d1603$*$}\\
    \texttt{mc16\_pPb8TeV.420012.Pythia8EvtGen\_A14NNPDF23LO\_jetjet\_JZ2R04.merge.AOD.e6394\_d1603$*$}\\
    \texttt{mc16\_pPb8TeV.420013.Pythia8EvtGen\_A14NNPDF23LO\_jetjet\_JZ3R04.merge.AOD.e6518\_d1603$*$}\\
    \texttt{mc16\_pPb8TeV.420014.Pythia8EvtGen\_A14NNPDF23LO\_jetjet\_JZ4R04.merge.AOD.e6518\_d1603$*$}\\
    \texttt{mc16\_pPb8TeV.420015.Pythia8EvtGen\_A14NNPDF23LO\_jetjet\_JZ5R04.merge.AOD.e6518\_d1603$*$}\\
}
\end{quote}
\begin{quote}
\SNN\,= 8.16 \TeV\ \pxPb\ Forward-filtered\\
\noindent\rule{12cm}{0.4pt}\\
{\scriptsize
    \texttt{mc16\_pPb8TeV.420018.Pythia8EvtGen\_A14NNPDF23LO\_jetjet\_JZ1R04\_MaxEta\_m3p0.merge.AOD.e8288\_d1603$*$}\\
    \texttt{mc16\_pPb8TeV.420019.Pythia8EvtGen\_A14NNPDF23LO\_jetjet\_JZ2R04\_MaxEta\_m3p0.merge.AOD.e8288\_d1603$*$}\\
}
\end{quote}


\clearpage

1D distributions of leading truth jet \PT{} and \ETA\ are shown in Figure~\ref{fig:incff} for the \pxPb\ inclusive MC sample (grey), and the combined inclusive and forward-filtered MC samples (red). Particular attention should be paid to $p_{\TN{T},1}^{\TN{truth}}~=~160$~GeV in the left figure and $\eta_{1}^{\TN{truth}}$~=~-3.0 in the right figure. This $p_{\TN{T},1}^{\TN{truth}}$ value corresponds to the transition between JZ2 and JZ3, where there is no forward-filtered MC sample. Therefore, no matter the value of $\eta_{1}^{\TN{truth}}$, any event with $p_{\TN{T},1}^{\TN{truth}}~>~160$~GeV will only be used if it comes from the inclusive MC sample. In the region where $\eta_{1}^{\TN{truth}}$~<~-3.0 and $p_{\TN{T},1}^{\TN{truth}}~<~160$~GeV, only events from the forward-filtered MC sample are used. The bottom panels in Figure~\ref{fig:incff} show that this strategy leads to a smooth transition, where the inclusive sample has been divided by the combined inclusive and forward-filtered sample, using the strategy outlined above. 



\begin{figure}[ht]
    \centering
    \includegraphics[width=0.95\linewidth,trim={0cm 4cm 0cm 8cm},clip]{sections/Datasets/fig/mc_combination.pdf}
    \caption{1D distributions of leading truth jet \PT{} (left) and \ETA\,(right) for both the \pxPb\,inclusive (INC) MC sample, denoted by grey squares, and the combined inclusive and forward-filtered (INC+FF) MC samples, denoted by red triangles. The bottom panels show the ratio of the inclusive sample to the combined inclusive and forward-filtered samples. }
    \label{fig:incff}
\end{figure}



\subsection{Issues related to the simulation of spectator neutrons in Athena}
\label{sec:zdc_mc_event_gen_issues}
The simulation of spectator neutrons in \AxA\ and \pxA\ collisions represents one of the biggest outstanding challenges in the MC description of HI collisions. Overall, a limited amount of data has been published at LHC energies, which would enable the modeling of the spectators in these collisions. In addition, to this date the ATLAS ZDC MC was never deployed in conjunction with the rest of the ATLAS MC chain, but instead in standalone mode with neutrons shot to the detector using a particle gun configured to distribute the particles taking into account the transverse Fermi momentum \cite{Moniz:1971mt}. In summary, currently there is no infrastructure available to: 1. realistically generate events describing both hard-scattering and nucleus fragmentation in \pxPb\ collisions and 2. run these events through a comprehensive simulation that includes both the ATLAS detector and the ZDC response. 
These aspects will be relevant for the discussion of the unfolding strategy, presented in Section~\ref{ssec:unf_strat}.


\subsection{Disabled HEC}
\label{sec:hec_issue}

2016 \pxPb\ data were affected by an issue in the hadronic calorimeter endcap, colloquially referred to as the ``disabled HEC'' or ``HEC issue''. The issue affected a set of channels in the 4$^{th}$ quadrant of the A-side (positive \ETA) of the HEC, that were disabled during the data taking operations because of a faulty LV power supply \cite{2016:HECJIRA}. The region of phase space affected was roughly [\ETA,$\phi$] $\in$ [1.5,3.2]$\times$[-$\pi$,-$\pi$/2], see Figure~\ref{fig:HECissue}, central panel. Jets reconstructed in this part of the detector are likely to have a mischaracterized response, since only the electromagnetic portion of the calorimeter was fully operational during the Run in this region. 
Therefore, this fraction of the phase space was removed from the analysis of \pxPb\, data (both beam orientations)
by rejecting events where at least one of the two leading jets was detected in the HEC problematic region, opportunely padded with a contour of thickness equal to R=0.4, \textit{i.e.} [\ETA,$\phi$] $\in$ [1.1,3.6]$\times$([$-\pi$,-$\pi$/2+0.4] $\cup$ [$\pi$-0.4,$\pi$]).

\begin{figure}[ht]
    \centering
    \includegraphics[width=0.32\linewidth,trim={0cm 0cm 0cm 0mm},clip]{sections/Datasets/fig/HEC/etaPhi_hec_nops_pp_data.pdf}
    \includegraphics[width=0.32\linewidth,trim={0cm 0cm 0cm 0mm},clip]{sections/Datasets/fig/HEC/etaPhi_hec_ps_Pbpdata_marked.pdf}
    \includegraphics[width=0.32\linewidth,trim={0cm 0cm 0cm 0mm},clip]{sections/Datasets/fig/HEC/etaPhi_hec_Pb_MC_marked.pdf}    
    \caption{
\textcolor{black}{
    $\phi$ vs \ETA\,distribution for 2017 \pxp\,data, 2016 \Pbxp\,data and 2016 \Pbxp\,\texttt{Pythia8} MC overlay. The \Pbxp\,orientation was chosen because of the boost ($\Delta y$ = +0.465) towards positive \ETA\,that helps in highlighting the dip due to the HEC issue. No (\ETA,$\phi$) cut is applied in any of the three distributions, in order to illustrate how the issue is not present in \pxp\,data (left), present in \Pbxp\,data (center) and reproduced in \Pbxp\,reconstructed MC (right). In the latter two distributions, the red (black) lines show the region corresponding to the (cut to reject events affected by the) HEC issue. All the distributions are obtained using jets with R~=~0.4 and \PT{} $>$ 20 GeV.
    }
    }
    \label{fig:HECissue}
\end{figure}
These cuts, illustrated by the black lines in Figure~\ref{fig:HECissue} central and right panels, carve out a rectangular shaped region in the [\ETA,$\phi$] of the jets.
This is taken into account at the level of the unfolding.



\subsection{Missing ZDC Module on Side C}
\label{sec:missing_zdc_mod}
During the 2016 \pxPb\ run, the ZDC side C module 0 (EM) was removed and replaced with the Arm2 calorimeter of the LHCf experiment \cite{Adriani:2135330}. Figure~\ref{fig:lhcf_v_zdc_acceptance} displays the transverse acceptance of both the ATLAS ZDC and LHCf.  The ATLAS ZDC EM Module is the first of 4 sequential calorimeter modules (counting from the perspective of looking at the ZDCs from the ATLAS interaction point); therefore, the LHCf calorimeter would see forward neutral products before they reach the side C ZDC.  The LHCf detector was put in place to measure the forward neutral products of the proton going side in the \pxPb\ orientation; however, in the \Pbxp\ orientation the LHCf detector was retracted into the "garage" position to avoid radiation damage. \cite{Adriani:2135330}  The "garage" position indicates an upwards vertical shift of 112 mm from the nominal position shown in Figure~\ref{fig:lhcf_v_zdc_acceptance}.  This shift places the full LHCf Arm2 calorimeter outside of the full acceptance of the ATLAS ZDC.  Still, the side C ZDC suffers from only having roughly 75\% of the material budget of side A. This effectively reduces the interaction length of the side C ZDC from 4.6 to 3.4, which in turn reduces the neutrons expected to shower in the ZDC from 99\% to 96.6\%. For this reason, in this analysis only Period B (\pxPb), is used to obtain the final results. 
For completeness, a comparison between the results from Period B (\pxPb) and C (\Pbxp) is reported in the Supporting Material, see Sections~\ref{auxsec:both_periods} and \ref{auxsec:both_periods_RCP}.

\begin{figure}[h]
    \centering
\begin{subfigure}[c]{0.45\linewidth}
    \includegraphics[width=1\linewidth]{sections//Datasets//fig/Screenshot 2024-06-27 at 11.44.20 AM.png}
\end{subfigure}
\begin{subfigure}[c]{0.45\linewidth}
        \includegraphics[width=1\linewidth]{sections//Datasets//fig/Screenshot 2024-06-27 at 11.22.36 AM.png}
\end{subfigure}
    \caption{Detector Geometry for LHCf Arm2 \cite{Cekmecelioglu:2867369} (left) and the ATLAS ZDC \cite{Jenni:1009649} (right). View is from the perspective of the ATLAS interaction point. In both figures the distances are given in mm.  The blue oval in the left image corresponds to the boundary of the yellow oval in the right image.  The active area of LHCf is represented by the filled blue squares.  The active area of the ZDC is given by the area bounded by the purple dashed line.}
    \label{fig:lhcf_v_zdc_acceptance}
\end{figure}

In the \pxPb\ orientation, the LHCf was in running position in runs 313063 and 313435, while for all other runs it was in "garage" position. Due to the asymmetric design of the LHCf module in the transverse plane, the effect on the acceptance of the ZDC on the proton going direction during these runs is non trivial.  Any measurement presented in this Note using the proton going ZDC will have these runs removed.




\clearpage{}
\clearpage{}\section{Centrality Determination}
\label{sec:centrality}
In \PbxPb\ collisions, the volume of the QGP droplet created is related to the impact parameter, or \textit{centrality}, of the collision. For this reason, many observables sensitive to QGP effects are usually studied in different centrality intervals to infer the property of the medium and understand its evolution. Centrality still plays an important role even in small systems such as \pxA\ collisions. Here it can be seen as an experimental classification of the collision geometry generally based on a measurement of the overall event activity in a rapidity region separated from the hard process of interest \cite{Perepelitsa:2014yta}. The centrality determination for the 2016 \pxPb\ data-set was carried out by K.Hill and D.Perepelitsa in 2018. All the details of these studies, used to define centrality in the analysis presented in this note, are available in \cite{Hill:2301540}. Main aspects relevant for the 8.16 \TeV\ dijet analysis are summarized in this section. 


\subsection{Centrality Selections}
In \pxPb\ collisions, the momentum asymmetry of the projectiles, resulting in a boost of the center-of-mass frame relative to the detector, lead to a corresponding shift in the rapidity distribution of soft particles. This shift corresponds to $\Delta y = \pm$0.465 for \Pbxp\ and \pxPb\ data taking, respectively. 



\begin{figure}[ht]
    \centering
    \includegraphics[width=0.45\linewidth,trim={0cm 0cm 0cm 0mm},clip]{sections/Centrality/fig/fCalECorr_pPb_.pdf}
    \includegraphics[width=0.45\linewidth,trim={0cm 0cm 0cm 0mm},clip]{sections/Centrality/fig/fCalECorr_Pbp_.pdf}
    \caption{Correlation between the transverse energy accumulated in the two FCal arms: $p$-going (side C for \pxPb\ and A for \Pbxp) and Pb-going (side A for \pxPb\ and C for \Pbxp). 
    Distributions obtained using the \texttt{HLT\_mb\_sptrk} trigger only and event selection analogous to the centrality analysis. 
    }
    \label{fig:FCal_correl}
\end{figure}

At low total energy in both FCals the correlation is roughly one-to-one, but as the energy in the Pb-going direction increases, the energy in the $p$-going direction seems to saturate. This implies a lack of sensitivity of the $p$-going direction to the overall soft particle production and, therefore, to the event centrality. For this reason, centrality was determined using solely the transverse energy registered in the Pb-going arm of the FCal (side A for \pxPb\ data, side C for \Pbxp\ data, respectively). The latter was correlated with the collision geometry thanks to Monte Carlo (MC) Glauber models. In particular, traditional MC Glauber and Glauber-Gribov models were considered in the analysis. All the details can be found in \cite{Hill:2301540}. Please note that, by using the Pb-going arm of the FCal, centrality is determined in a kinematic region totally separated from the region considered for the analysis presented in this Note since, as discussed in Section~\ref{sec:trig_evsel}, the jet triggers in the backward region (e.g. the Pb going side of the experiment) were not active for most of the data taking. 

Centrality selections taken from \cite{Hill:2301540} and used for this analysis are reported in Table~\ref{tab:cent_bin}.
\begin{table}[!ht]
\centering
\begin{tabular}{ | c | c | }
           \hline
           \rule{0pt}{2.5ex}  
           Centrality   &  \hspace{4mm} \ETPB\ [TeV] \\[0.5ex]
           \hline
           \hline
            0--1 \%      & 0.10898 - 6.00   \\
            \hline 
            1--5 \%      & 0.7953 - 0.10898 \\
            \hline
            5--10 \%     & 0.06518 - 0.7953 \\
            \hline
            10--20 \%    & 0.04916 - 0.06518 \\
            \hline
            20--30 \%    & 0.03866 - 0.04916 \\
            \hline
            30--40 \%    & 0.03046 - 0.03866 \\
            \hline
            40--50 \%    & 0.02362 - 03046 \\
            \hline
            50--60 \%    & 0.01769 - 0.02362 \\
            \hline
            60--70 \%    & 0.1244 - 0.01769 \\
            \hline
            70--80 \%    & 0.00733- 0.1244 \\
            \hline
            80--90 \%    & 0.00360 - 0.00733 \\
            \hline             
        \end{tabular}
        \vspace{2mm}
        \caption{Definition of centrality binning based on  the distribution of transverse energy in the Pb-going FCal (\ETPB) for a nominal efficiency value of 98\%. Taken from \cite{Hill:2301540}.}
        \label{tab:cent_bin}
\end{table}

Refer to \cite{Hill:2301540} for all details related to the determination of the centrality in the 2016 $p$+Pb data-set, and the associated uncertainties.

Centrality, as defined in Table~\ref{tab:cent_bin}, has been proven to be a non-robust estimator for the event geometry in \pxPb, as discussed in \cite{ATLAS:2023zfx}.
Still, this definition is partially used in this note to construct a central to peripheral ratio of average neutron multiplicities in the same event activity selection of the analysis presented in \cite{ATLAS:2023zfx}, to potentially add further information useful for the modeling of the \pxPb\ collisions. For this purpose, the same three centrality intervals utilized in the paper are used:
\begin{itemize}
    \item 0--10\%, Central events,  0.06518 < $\Sigma$\ETPB\,/ [TeV] < 6.00;
    \item 10--60\%, Mid-central events,  0.01769 < $\Sigma$\ETPB\,/ [TeV] < 0.06518;
    \item 60--90\%, Peripheral events, 0.00360 < $\Sigma$\ETPB\,/ [TeV] < 0.01769;
\end{itemize}

An example of FCal transverse energy distribution obtained from MB events and divided in these three different centrality intervals is shown in Figure~\ref{fig:MB_FCAL}. A comparison between FCal distribution in MB events and in dijet events selected for the analysis is reported in Figure~\ref{fig:MB_vs_dijet}.  The greater number of dijet events in most central collisions is due to higher \NCOLL\ in central events, enhancing the probability of a dijet event occurrence. 

\begin{figure}[ht]
    \centering
    \includegraphics[width=0.6\linewidth,trim={0cm 0cm 0cm 0mm},clip]{sections/Centrality/fig/FCalCent_hist.pdf}
    \caption{ 
    Distribution of $\Sigma$\ETPB\ in the FCal for \pxPb\ data taking, extracted from MB events where the \texttt{HLT\_mb\_sptrk} trigger was fired. The two different centrality bins used in the analysis are denoted by alternating shaded and un-shaded areas, namely 0--10\% (central - blue), 10--60\% (mid-central, white) and 60--90\% (peripheral, blue). 
    }
    \label{fig:MB_FCAL}
\end{figure}

\begin{figure}[ht]
    \centering
    \includegraphics[width=0.5\linewidth,trim={0cm 0cm 0cm 0mm},clip]{sections/Centrality/fig/MB_vs_dijet.png}
    \caption{ 
    Distributions of $\Sigma$\ETPB\ in the FCal for \pxPb\ data taking, extracted from MB (blue) and dijet (red) events. The \texttt{HLT\_mb\_sptrk} trigger was used to obtain the MB distribution. The dijet distribution was obtained using the same event selections applied in the analysis. 
    }
    \label{fig:MB_vs_dijet}
\end{figure}

\subsection{Number of Participating Nucleons}
\label{ssec:TABsec}
Within the centrality analysis reported in \cite{Hill:2301540}, the number of participating nucleons, \NPART\ was also estimated in the same centrality intervals. It is reported in Table~\ref{tab:npart} for two different MC Glauber models. 
The standard Glauber model uses a fixed nucleon-nucleon cross-section, while the Glauber-Gribov model allows the nucleon-nucleon cross-section to fluctuate on an event by event basis. For more details, refer to Chapter 5 of \cite{Hill:2301540}.

Like the centrality intervals, \NPART\ was determined in a finer binning compared to the one used for this analysis. The approach chosen in this case was to build a weighted average \cite{UnbWeight}, where the weights are set to the width of the centrality interval associated to \NPART. It is worth noting that values of \NPART\ evaluated in finely binned centrality intervals are characterized by asymmetric uncertainties, while their weighted average\footnote{In determining the uncertainty on the weighted average, a conservative approach is adopted by taking the larger of the uncertainty value of each of the narrow centrality class.} returns a symmetrized uncertainty.


\begin{table}
    \centering
    \begin{tabular}{|l|c|c|c|c|}
    \hline
    \multirow{2}{*}{\textbf{Centrality}} & \multicolumn{4}{c|}{\textbf{\NPART}} \\
& \multicolumn{2}{c|}{Standard Glauber} & \multicolumn{2}{c|}{Glauber-Gribov}  \\
\hline
    \rule{0pt}{2.5ex} 
    0--1\% & 19.10$^{+0.61}_{-1.28}$ & \multirow{3}{*}{16.37$\pm$1.00} & 25.60$^{+0.55}_{-1.61}$ & \multirow{3}{*}{19.49$\pm$2.17}  \\ [0.5ex]
    1--5\% & 16.94$^{+0.52}_{-0.90}$ & & 20.60$^{+0.47}_{-0.98}$ & \\[0.5ex]
    5--10\% & 15.37$^{+0.45}_{0.69}$ & & 17.37$^{+0.42}_{-0.62}$ &  \\[0.5ex]
    \hline
    \rule{0pt}{2.5ex} 
    10--20\% & 13.76$^{+0.37}_{0.51}$ & \multirow{5}{*}{10.41$\pm$1.18} & 14.52$^{+0.38}_{0.37}$ &  \multirow{5}{*}{10.07$\pm$1.46} \\[0.5ex]
20--30\% & 12.03$^{+0.28}_{-0.36}$ &  & 11.87$^{+0.34}_{-0.21}$ &   \\[0.5ex] 
    30--40\% & 10.41$^{+0.21}_{-0.24}$ & & 9.77$^{+0.32}_{-0.15}$ & \\[0.5ex]
    40--50\% & 8.76$^{+0.20}_{0.18}$ & & 7.92$^{+0.31}_{-0.11}$ &  \\[0.5ex]
    50--60\% & 7.08$^{+0.36}_{0.14}$ & & 6.25$^{+0.32}_{0.10}$ &  \\[0.5ex]
    \hline
    \rule{0pt}{2.5ex} 
    60--70\% & 5.46$^{+0.55}_{-0.13}$ & \multirow{3}{*}{4.21$\pm$0.69} & 4.81$^{+0.31}_{-0.09}$ & \multirow{3}{*}{3.75$\pm$0.58}  \\[0.5ex]
    70--80\% & 4.08$^{+0.64}_{-0.10}$ & & 3.64$^{+0.30}_{-0.07}$ & \\[0.5ex]
    80--90\% & 3.09$^{+0.54}_{0.07}$ & & 2.81$^{+0.24}_{-0.06}$ &  \\[0.5ex]
    \hline
    \rule{0pt}{2.5ex}
    0--90\%& \multicolumn{2}{c|}{8.80$^{+0.33}_{-0.18}$} & \multicolumn{2}{c|}{8.86$^{+0.22}_{-0.18}$}  \\[0.5ex]
    \hline
    \end{tabular}
    \vspace{0.5cm}
    \caption{ \textcolor{black}{ 
    Values of \NPART\,calculated using standard Glauber and Glauber-Gribov models for different centrality classes, taken from \cite{Hill:2301540}. The centrality-weighted mean for values corresponding to the centrality classes used in \cite{ATLAS:2023zfx} and in some part of this analysis, computed as described in the text, is also reported. These values currently represent a conservative estimate of \NPART. } 
    }
    \label{tab:npart}
\end{table}





%
\clearpage{}
\clearpage{}\section{Trigger Studies}
\label{sec:trig_evsel}


\subsection{ATLAS Jet Trigger in Run 2}
The ATLAS trigger system consists of a hardware-based component, the Level-1 (L1) trigger, and a software-based trigger system, the higher-level trigger (HLT). The HLT runs a simplified version of the ATLAS reconstruction software. At the HLT level, calibrated jets defined by the \antikt\ algorithm with R~=~0.4 are used. 
Each trigger item is composed of an L1 trigger and an HLT trigger. The L1 system first evaluates each given trigger if it is passed. To keep the trigger rate at an acceptable level, and to not saturate the available bandwidth, triggers with a high rate are only read out for a fraction of all events by setting a prescale value, characteristic of each trigger item. The prescale is a number applied online by the ATLAS trigger system. If an event satisfies the requirements of a given trigger, the prescales for the L1 and HLT triggers are determined. If the L1 or the HLT trigger is prescaled, the HLT trigger decision is not evaluated to optimize the computing time. Prescales for the L1 trigger and the HLT trigger are decided independently. 

Single-jet triggers are designed to fire with events characterized by large transverse energy depositions in the ATLAS calorimeter system. 
The HLT trigger name encodes the nominal threshold value (expressed in GeV) set for the HLT jet. Central and forward jets are constructed in two separate systems. The transition between central and forward jets occurs at |\ETA| = 3.1  and at |\ETA| = 3.2 for the L1-Calo system and HLT, respectively.
As the name suggests, central jet triggers cover the pseudorapidity region |\ETA| $<$ 3.2, while forward jet triggers cover the range
3.2 < |\ETA| < 4.9. The forward jet triggers are denoted with a suffix \texttt{320eta490}, highlighting their acceptance. In \pxPb\ data taking, where the collision system is asymmetric, the forward triggers are separated by being either positive or negative as a function of their \ETA\ coverage. The two orientations are distinguished by a letter \texttt{p} or \texttt{n} before \texttt{320eta490}, respectively. $\PTE^{99}$ represents the 99\% efficiency point for a given chain, calculated taking into account both L1 and HLT efficiencies, and is reported in GeV.
It is worth noting that a set of semi-forward triggers were enabled during the 2016 data-taking. These triggers cover 2.0 < |\ETA| < 3.2 and are denoted with a suffix \texttt{p200eta320} (\texttt{n200eta320}) if they cover the positive (negative) pseudorapidity region. 
This measurement will use only the central and forward jet triggers. 


\subsection{Jet triggers in the 2016 \pxPb~Run}
As discussed in Section~\ref{sec:data}, the 2016 \pxPb\, data set is characterized by two different beam orientations. Trigger recommendations for this data-taking are available on the Heavy Ion Trigger Menu Forum \cite{TriggerForum:pPb2016}. Two different trigger menus were deployed during the Run, with \texttt{ion} and \texttt{pp} triggers used as primary triggers, respectively. The switch between the two configurations did not happen in between period B (\pxPb) and C (\Pbxp) but after a few runs of period C. In addition, given the asymmetric collision system, trigger configurations are different for forward ($p$-going) and backward (Pb-going) arms of the detector\footnote{Forward triggers in the backward part of the detector were turned off for most of the data taking and therefore are not considered in this analysis.}, yielding three different trigger schemes. 
Please refer to Chapter~6 of \cite{Longo:2860813} for detailed information regarding the extensive trigger studies that were carried out for the dataset analyzed in this analysis.  Tables~\ref{tab:trigger_overview_pPb}--\ref{tab:trigger_overview_Pbp_pp} summarize the results of these trigger studies for the three trigger schemes present in this dataset.
Below we summarize how different fully efficient jet triggers are combined, using the same approach as \cite{ATLAS:2023zfx}.

\begin{table}[!ht]
\centering
\begin{tabular}{ | c | c | c | c | c |}
           \hline
           \textbf{Trigger Chain} &  \textbf{$\PTE^{99}$ [GeV]} & \textbf{$\eta$} & \textbf{$\mathscr{L}$ [ub$^{-1}$]} & \textbf{$p_{\TN{T}}^{99}$ Monitor}\\
            \hline
            \multicolumn{1}{|l|}{\texttt{HLT\_j30\_ion\_0eta490\_L1TE10}}   & 36  & [$-$4.9,4.9]   & 113.76 &  \multicolumn{1}{l|}{\scriptsize HLT\_mb\_sptrk} \\
            \hline
            \multicolumn{1}{|l|}{\texttt{HLT\_j40\_ion\_L1J5}}    & 44     & [$-$3.2,3.2]   & 468.1 &  \multicolumn{1}{l|}{\scriptsize HLT\_mb\_sptrk}\\
            \hline
            \multicolumn{1}{|l|}{\texttt{HLT\_j50\_ion\_L1J10}}    & 57    & [$-$3.2,3.2]   & 1120 &  \multicolumn{1}{l|}{\scriptsize HLT\_mb\_sptrk}\\
            \hline
            \multicolumn{1}{|l|}{\texttt{HLT\_j60\_ion\_L1J20}} & 70        & [$-$3.2,3.2]   & 1562 &  \multicolumn{1}{l|}{\scriptsize HLT\_mb\_sptrk}\\
            \hline
            \multicolumn{1}{|l|}{\texttt{HLT\_j75\_ion\_L1J20}}   & 80  & [$-$3.2,3.2]   & 3468 &  \multicolumn{1}{l|}{\scriptsize HLT\_mb\_sptrk}\\
            \hline
            \multicolumn{1}{|l|}{\texttt{HLT\_j100\_ion\_L1J20}}   & 111 & [$-$3.2,3.2]   & 56671 &  \multicolumn{1}{l|}{\scriptsize HLT\_mb\_sptrk}\\
            \hline
            \multicolumn{1}{|l|}{\texttt{HLT\_j15\_ion\_n320eta490\_L1MBTS\_1\_1}}   & 44  & [$-$3.2,$-$4.9]   & 77.5 &  \multicolumn{1}{l|}{\scriptsize HLT\_mb\_sptrk}\\
            \hline
            \multicolumn{1}{|l|}{\texttt{HLT\_j25\_ion\_n320eta490\_L1TE5}}   & 40  & [$-$3.2,$-$4.9]   & 247.7 &  \multicolumn{1}{l|}{\scriptsize HLT\_mb\_sptrk}\\
            \hline
            \multicolumn{1}{|l|}{\texttt{HLT\_j35\_ion\_n320eta490\_L1TE10}}   & 57  & [$-$3.2,$-$4.9]   & 179.9 &  \multicolumn{1}{l|}{\tiny HLT\_j25\_ion\_n320eta490\_L1TE5}\\
            \hline
            \multicolumn{1}{|l|}{\texttt{HLT\_j45\_ion\_n320eta490}}   & 65  & [$-$3.2,$-$4.9]   & 5288 &  \multicolumn{1}{l|}{\tiny HLT\_j25\_ion\_n320eta490\_L1TE5}\\
            \hline
            \multicolumn{1}{|l|}{\texttt{HLT\_j55\_ion\_n320eta490}}   & 80  & [$-$3.2,$-$4.9]   & 16635 &  \multicolumn{1}{l|}{\tiny HLT\_j25\_ion\_n320eta490\_L1TE5}\\
            \hline
            \multicolumn{1}{|l|}{\texttt{HLT\_j65\_ion\_n320eta490}}   & 85  & [$-$3.2,$-$4.9]   & 53571 &  \multicolumn{1}{l|}{\tiny HLT\_j25\_ion\_n320eta490\_L1TE5}\\
            \hline
        \end{tabular}
        \vspace{2mm}
        \caption{
        \textcolor{black}{
        Overview of \pxPb\,\texttt{ion} triggers used in the analysis. The sampled luminosity, $\mathscr{L}$, recorded by each trigger is reported. The L1 for triggers that do not specify this detail in the chain name are:  \texttt{L1\_J15.31ETA49} for \texttt{HLT\_j45\_ion\_n320eta490} and \texttt{HLT\_j55\_ion\_n320eta490}, and \texttt{L1\_J20.31ETA49} for \texttt{HLT\_j65\_ion\_n320eta490}. $\PTE^{99}$ represents the 99\% efficiency point for a given chain, calculated taking into account both L1 and HLT efficiencies, and is reported in GeV.
        }
        }
        \label{tab:trigger_overview_pPb}
\end{table}

\begin{table}[!ht]
\centering
\begin{tabular}{ | c | c | c | c | c |}
           \hline
           \textbf{Trigger Chain} &  \textbf{$\PTE^{99}$ [GeV]} & \textbf{$\eta$} & \textbf{$\mathscr{L}$ [ub$^{-1}$]} & \textbf{$p_{\TN{T}}^{99}$ Monitor}\\
            \hline
            \multicolumn{1}{|l|}{\texttt{HLT\_j30\_ion\_0eta490\_L1TE10}}   & 38  & [$-$4.9,4.9]   & 30.13 &  \multicolumn{1}{l|}{\scriptsize HLT\_mb\_sptrk} \\
            \hline
            \multicolumn{1}{|l|}{\texttt{HLT\_j40\_ion\_L1J5}}    & 50     & [$-$3.2,3.2]   & 120.9 &  \multicolumn{1}{l|}{\scriptsize HLT\_mb\_sptrk}\\
            \hline
            \multicolumn{1}{|l|}{\texttt{HLT\_j50\_ion\_L1J10}}    & 57    & [$-$3.2,3.2]   & 599.6 &  \multicolumn{1}{l|}{\scriptsize HLT\_mb\_sptrk}\\
            \hline
            \multicolumn{1}{|l|}{\texttt{HLT\_j60\_ion\_L1J20}} & 75        & [$-$3.2,3.2]   & 745.3 &  \multicolumn{1}{l|}{\scriptsize HLT\_mb\_sptrk}\\
            \hline
            \multicolumn{1}{|l|}{\texttt{HLT\_j75\_ion\_L1J20}}   & 85  & [$-$3.2,3.2]   & 1223 &  \multicolumn{1}{l|}{\scriptsize HLT\_mb\_sptrk}\\
            \hline
            \multicolumn{1}{|l|}{\texttt{HLT\_j100\_ion\_L1J20}}   & 111 & [$-$3.2,3.2]   & 16674 &  \multicolumn{1}{l|}{\scriptsize HLT\_mb\_sptrk}\\
            \hline
            \multicolumn{1}{|l|}{\texttt{HLT\_j15\_ion\_p320eta490\_L1MBTS\_1\_1}}   & 44  & [3.2,4.9]   & 27.6 &  \multicolumn{1}{l|}{\scriptsize HLT\_mb\_sptrk}\\
            \hline
            \multicolumn{1}{|l|}{\texttt{HLT\_j25\_ion\_p320eta490\_L1TE5}}   & 31  & [3.2,4.9]   & 115.4 &  \multicolumn{1}{l|}{\scriptsize HLT\_mb\_sptrk}\\
            \hline
            \multicolumn{1}{|l|}{\texttt{HLT\_j35\_ion\_p320eta490\_L1TE10}}   & 44  & [3.2,4.9]   & 277.2 &  \multicolumn{1}{l|}{\tiny HLT\_j25\_ion\_p320eta490\_L1TE5}\\
            \hline
            \multicolumn{1}{|l|}{\texttt{HLT\_j45\_ion\_p320eta490}}   & 57  & [3.2,4.9]   & 2628 &  \multicolumn{1}{l|}{\tiny HLT\_j25\_ion\_p320eta490\_L1TE5}\\
            \hline
            \multicolumn{1}{|l|}{\texttt{HLT\_j55\_ion\_p320eta490}}   & 61  & [3.2,4.9]   & 8612 &  \multicolumn{1}{l|}{\tiny HLT\_j25\_ion\_p320eta490\_L1TE5}\\
            \hline
            \multicolumn{1}{|l|}{\texttt{HLT\_j65\_ion\_p320eta490}}   & 75  & [3.2,4.9]   & 15559 &  \multicolumn{1}{l|}{\tiny HLT\_j25\_ion\_p320eta490\_L1TE5}\\
            \hline
        \end{tabular}
        \vspace{2mm}
        \caption{
        \textcolor{black}{
        Overview of \Pbxp\,\texttt{ion} triggers used in the analysis. The sampled luminosity, $\mathscr{L}$, recorded by each trigger is reported. The L1 for triggers that do not specify this detail in the chain name are:  \texttt{L1\_J15.31ETA49} for \texttt{HLT\_j45\_ion\_p320eta490} and \texttt{HLT\_j55\_ion\_p320eta490}, and \texttt{L1\_J20.31ETA49} for \texttt{HLT\_j65\_ion\_p320eta490}. $\PTE^{99}$ represents the 99\% efficiency point for a given chain, calculated taking into account both L1 and HLT efficiencies, and is reported in GeV. 
}
        }
        \label{tab:trigger_overview_Pbp_ion}
\end{table}

\begin{table}[!ht]
\centering
\begin{tabular}{ | c | c | c | c | c |}
           \hline
           \textbf{Trigger Chain} &  \textbf{$\PTE^{99}$ [GeV]} & \textbf{$\eta$} & \textbf{$\mathscr{L}$ [ub$^{-1}$]} &\textbf{$p_{\TN{T}}^{99}$ Monitor}\\
            \hline
            \multicolumn{1}{|l|}{\texttt{HLT\_j30\_0eta490\_L1TE10}}   & \color{black}{40}  & [$-$4.9,4.9]   & 108.48 &  \multicolumn{1}{l|}{\scriptsize HLT\_mb\_sptrk} \\
            \hline
            \multicolumn{1}{|l|}{\texttt{HLT\_j40\_L1J5}}    & 54     & [$-$3.2,3.2]   & 463.8 &  \multicolumn{1}{l|}{\scriptsize HLT\_mb\_sptrk}\\
            \hline
            \multicolumn{1}{|l|}{\texttt{HLT\_j50\_L1J10}}    & 61    & [$-$3.2,3.2]   & 2169 &  \multicolumn{1}{l|}{\scriptsize HLT\_mb\_sptrk}\\
            \hline
            \multicolumn{1}{|l|}{\texttt{HLT\_j60}} & 75        & [$-$3.2,3.2]   & 3036 &  \multicolumn{1}{l|}{\scriptsize HLT\_mb\_sptrk}\\
            \hline
            \multicolumn{1}{|l|}{\texttt{HLT\_j75\_L1J20}}   & 85  & [$-$3.2,3.2]   & 21717 &  \multicolumn{1}{l|}{\scriptsize HLT\_mb\_sptrk}\\
            \hline
            \multicolumn{1}{|l|}{\texttt{HLT\_j90\_L1J20}}   & 104  & [$-$3.2,3.2]   & 68464 &  \multicolumn{1}{l|}{\scriptsize HLT\_mb\_sptrk}\\
            \hline
            \multicolumn{1}{|l|}{\texttt{HLT\_j100\_L1J20}}   & 111 & [$-$3.2,3.2]   & 86453 &  \multicolumn{1}{l|}{\scriptsize HLT\_mb\_sptrk}\\
            \hline
            \multicolumn{1}{|l|}{\texttt{HLT\_j15\_p320eta490\_L1MBTS\_1\_1}}   & 24  & [3.2,4.9]   & 61.75 &  \multicolumn{1}{l|}{\scriptsize HLT\_mb\_sptrk}\\
            \hline
            \multicolumn{1}{|l|}{\texttt{HLT\_j25\_p320eta490\_L1TE5}}   & 34  & [3.2,4.9]   & 310.9 &  \multicolumn{1}{l|}{\scriptsize HLT\_mb\_sptrk}\\
            \hline
            \multicolumn{1}{|l|}{\texttt{HLT\_j35\_p320eta490\_L1TE10}}   & 47  & [3.2,4.9]   & 685.2 &  \multicolumn{1}{l|}{\scriptsize HLT\_j25\_p320eta490\_L1TE5}\\
            \hline
            \multicolumn{1}{|l|}{\texttt{HLT\_j45\_p320eta490}}   & 65  & [3.2,4.9]   & 8274 &  \multicolumn{1}{l|}{\scriptsize HLT\_j25\_p320eta490\_L1TE5}\\
            \hline
            \multicolumn{1}{|l|}{\texttt{HLT\_j55\_p320eta490}}   & 80  & [3.2,4.9]   & 20076 &  \multicolumn{1}{l|}{\scriptsize HLT\_j25\_p320eta490\_L1TE5}\\
            \hline
            \multicolumn{1}{|l|}{\texttt{HLT\_j65\_p320eta490}}   & 91  & [3.2,4.9]   & 56342 &  \multicolumn{1}{l|}{\scriptsize HLT\_j25\_p320eta490\_L1TE5}\\
            \hline
        \end{tabular}
        \vspace{2mm}
        \caption{
        \textcolor{black}{
        Overview of \Pbxp\,\texttt{pp} triggers used in the analysis. The sampled luminosity, $\mathscr{L}$, recorded by each trigger is reported. The L1 for triggers that do not specify this detail in the chain name are: \texttt{L1\_J15.31ETA49} for \texttt{HLT\_j45\_p320eta490} and \texttt{HLT\_j55\_p320eta490}, and \texttt{L1\_J20.31ETA49} for \texttt{HLT\_j65\_p320eta490}. $\PTE^{99}$ represents the 99\% efficiency point for a given chain, calculated taking into account both L1 and HLT efficiencies, and is reported in GeV. 
        }
        }
        \label{tab:trigger_overview_Pbp_pp}
\end{table}

\subsection{Trigger Strategy}
\label{sec:FTS}
Using the information outlined above, a well-defined (\ETA,\PT{}) map of triggers to be used in this analysis was constructed for each of the three trigger schemes used in the \pxPb\, data taking. Events are considered only if the leading jet in the trigger has fired the corresponding trigger chosen for the (\ETA,\PT{}) region where the jet is reconstructed. 
The trigger selection for |\ETA|$<$3.0 and for |\ETA|$>$3.5 
\textcolor{black}{
above a \PT{ }of 40 GeV
}
is straightforward. In these regions, the fully efficient efficient trigger associated with the lowest prescale value is chosen. Efficiency values used to determine what trigger to choose in each (\ETA,\PT{}) region are reported in Tables~\ref{tab:trigger_overview_pPb}, \ref{tab:trigger_overview_Pbp_ion}, and \ref{tab:trigger_overview_Pbp_pp}. For triggers covering the same \ETA\,region, prescale values decrease as the \PT{ }threshold increases. 

A different treatment is requested in the \ETA\,region corresponding to the transition between central and forward triggers, i.e. around \ETA = $\pm$3.2\footnote{Sign here is determined by the beam orientation, since this analysis covers only central and forward regions, discarding the backward where the energy accumulated in the FCal is used for centrality determination.}. To handle the trigger analysis in this region, an inclusive approach (see \cite{Lendermann:2009ah} and references therein), hereafter referred to as the Fully Efficient Trigger Combination (FETC) method, was chosen. 
In the FETC approach, a combined weight based on all the considered trigger items for the given kinematic region is determined for the entire event sample. For each event, at least one actual trigger item bit from the set of considered items is required to be set. The weight ($w_{j,k}$) calculation for an event $j$ in a given run $k$ is based on the probability to accept the event after the application of the prescale to all the $i$-th trigger items considered. 
If $r_{i,j}$ is the bit for the trigger item $i$ in the event $j$, and the same trigger item is prescaled by a factor $p_{i,k}$ during the run $k$, the probability that the event $j$ in the run $k$ is accepted by the trigger item $i$ after the prescale can be expressed as
\begin{equation}
    q_{i,j,k} = \frac{r_{i,j}}{p_{i,k}}
\end{equation}
Assuming the prescale decisions for each trigger item are independent from each other, the probability that at least one of the $N_{\TN{items}}$ trigger items accepts the event is given by 
\begin{equation}
    q_{j,k} = 1 - \prod^{N_{\TN{items}}}_{i=1} \biggl ( 1-\frac{r_{i,j}}{p_{i,k}}\biggr )
\end{equation}
Starting from this, one can define the run dependent weight for event $j$ as 
\begin{equation}
    w_{j,k} = \frac{1}{q_{j,k}}
\end{equation}

When combining all the considerations outlined above, one obtains a (\ETA,\PT{}) map that provides, per each jet with \PT{}$>$ 40 GeV
and -4.5 $<$ \ETA $<$ 3.1 (-3.1 $<$ \ETA $<$ 4.5), the jet trigger(s) considered for the analysis of the \pxPb\,(\Pbxp) data. The three maps corresponding to the trigger configurations used in the \pxPb\,run are shown in Figure~\ref{fig:trig_map_pPb}. The analysis reported in this note will use dijet events with a leading jet 
\PT{}$>$ 40 GeV and a sub-leading jet \PT{}$>$ 30 GeV. Only the dataset using \pxPb-ion trigger scheme is considered for the analysis presented in the CONF note, while the other two periods are used just for studies included in the Appendix of this Note and reported here for completeness. 



\begin{figure}[ht]
    \centering
    \includegraphics[width=0.48\linewidth,trim={0cm 0cm 0cm 0mm},clip]{sections/TriggerAndES/fig/TriggerThresholdMap_pPb.pdf}
    \includegraphics[width=0.48\linewidth,trim={0cm 0cm 0cm 0mm},clip]{sections/TriggerAndES/fig/TriggerThresholdMap_Pb_ion.pdf}  
    \includegraphics[width=0.48\linewidth,trim={0cm 0cm 0cm 0mm},clip]{sections/TriggerAndES/fig/TriggerThresholdMap_Pbp_pp.pdf} 
    \caption{(\ETA,\PT) maps identifying the trigger selected for analysis of an event where the leading jet is characterized by a given \ETA and \PT. Different colors corresponds to different triggers. Overlapping colors represent region where FETC method is applied. The three panels correspond to the three trigger configuration of the 2016 \pxPb\,data: \pxPb\, (top left), \Pbxp\,\texttt{ion} triggers (top right), \Pbxp\, \texttt{pp} triggers (bottom). } 
    \label{fig:trig_map_pPb}
\end{figure}\clearpage{}
\clearpage{}\section{Jet Selection and Reconstruction Performance} 
\label{sec:JetReco}

When studying HI collisions with the ATLAS detector, jets are reconstructed using the \antikt\ jet clustering algorithm \cite{Cacciari_2008}. This algorithm uses collections of four-vectors corresponding to a set of tracks of charged particles in the inner tracking system (aka ``track jets''), massless calorimeter towers with size $\Delta \eta \times \Delta \phi = 0.1 \times 0.1$ (aka ``HI jets''), or sets of simulated particles at the generator level (aka ``truth jets''). The measurement presented in this note makes use of HI jets reconstructed using the \antikt\ algorithm with R~=~0.4. Truth jets are used, together with the reconstructed ones, to carry out several performance studies reported in this note. 
A detailed review of the performance of jet reconstruction in the HI case can be found in \cite{ATLAS-CONF-2015-016}. This note will only discuss the main features of the HI jet reconstruction and their implications on the dijet cross-section measurement.

\subsection{Underlying Event Subtraction}
To reconstruct jets in HI collisions, a large background arising from the underlying event (UE) must be subtracted from each tower to avoid large rates of fake jets \cite{PhysRevC.86.024908}. The energy registered in each cell can be decomposed into contributions from a jet and from the UE as
\begin{equation}
    \frac{d^2E_{\TN{T}}^{\TN{total}}}{d\eta d\phi}=\frac{d^2E_{\TN{T}}^{\TN{jet}}}{d\eta d\phi}+\frac{d^2E_{\TN{T}}^{\TN{UE}}}{d\eta d\phi}
\end{equation}
Only the first contribution should be used to reconstruct jets. Therefore, $\frac{d^2E_{\TN{T}}^{\TN{UE}}}{d\eta d\phi}$ must be estimated and subtracted. 
The UE subtraction procedure makes use of the UE average transverse energy density, $\rho (\eta,\phi)$. 
Additional information on the subtraction of elliptic flow during the HI jet reconstruction used for ATLAS Run 2 data can be found in \cite{Rinn:2711618, Rybar:2688696, Longo:2860813} and in the references therein. 

The average transverse energy density is estimated using an iterative procedure that excludes regions populated by jets. A detailed description of this procedure can be found in \cite{Rinn:2711618}. 


\subsection{Jet Reconstruction and Calibration}
Individual calorimeter cells are themselves calibrated, but there is a variety of factors (\textit{e.g.} passive material in the detector, gaps between cells, signal losses, etc. \cite{PERF-2011-03}) that allow only for a partial measurement of the energy deposited by a single jet in the detector. To compensate for these experimental effects, the energy scale of reconstructed jets needs to be calibrated to the truth energy scale. 
The calibration sequence for HI jets is presented in \cite{Krivos:2706256}. 
Here we briefly review the main steps of this procedure. 
After the UE subtraction, two different types of calibrations are applied, namely: 
\begin{enumerate}
    \item EtaJES MC-based calibration, to recover the truth energy scale in MC. 
    \item HI/EMTopo cross-calibration and in-situ / $\eta$-intercalibration, to account for differences between MC and data. 
\end{enumerate}

The EtaJES calibration is typically the larger between the two and is derived from \texttt{Pythia8} simulated dijet events. The scope of this calibration is to correct for the energy scale component that is common to both data and MC. The per-jet energy correction factor is computed using a numerical inversion procedure \cite{Kosek:1742072, Ouellette:2655891}, see \cite{Cukierman:2016dkb} for a review of the procedure's math. A typical jet energy scale correction, $E_{\TN{T}}^{\TN{true}}/E_{\TN{T}}^{\TN{reco}}$, ranges from $\sim$ 1.2 to 2.0. Within this step of the calibration procedure a small pseudorapidity correction is also applied to the jet reconstructed \ETA, to counteract the implicit assumption that the $z$-vertex\footnote{Position of the vertex associated to the jet along the beam axis.} = 0 made during the reconstruction of the data. For this reason, this procedure is called ``origin correction''.

The second calibration component is only applied to data to account for discrepancies between MC and data. The in-situ and \ETA-intercalibrations implement residual corrections to the jet energy scale based on the \PT\ balance in vector boson ($\gamma$, Z$^0$) - jet events (the \textit{in-situ} correction) and dijet events where mid-rapity jets are used to calibrate the forward rapidity jets (the \ETA-intercalibration).
Both factors have been investigated in detail for EMTopo jets in \pxp\ within the JetEtMiss group, but they cannot be directly applied to HI jets because of the different type of constituents used in the two cases.  The cross-calibration, derived using the direct balance method in events containing a Z boson or photon opposite to a jet, is used to bridge between HI and fully calibrated EMTopo jets by accounting for additional differences in response between them \cite{Ouellette:2655891}. Thanks to the cross-calibration, it is possible to characterize the performance of the HI jets within the EMTopo jet framework. This procedure allows using of the baseline component as recommended by JetEtMiss group. 



\subsection{Jet Reconstruction Performance}
Criteria used to select jets in this analysis are discussed in Section~\ref{sec:data}. All the jet performance studies were carried out for both the beam orientations of \pxPb\ data taking.
All the jet performance studies presented in this note were carried out using the \texttt{HIJetValTools} package \cite{2022:HIJetValTools} modified to accommodate peculiarities of the \pxPb\ data taking \cite{2022:HIJetValTools-pPb} (beam dependent configuration, centrality determination using only one side of the FCal, triggers etc.). 

The analysis of jet reconstructed performance is carried out segmenting the detector into 6 different \ETA\ regions, characterized by the following limits: -4.5, -3.2, -1.5, 0, 1.5, 3.2, 4.5. In the case of \pxPb\ collisions, the Pb going \ETA\ bin is excluded from the analysis (\textit{i.e.} [-4.5, -3.2] for \Pbxp\ data, and [3.2, 4.5] for \pxPb\ collisions. In addition to this choice, determined by the lack of triggers covering this region and by the intention of avoiding contamination from jets in the portion of the FCal used for centrality determination, an additional reduction of \ETA\ = 0.1 is applied to the outermost surviving backward bin, to avoid the trigger \textit{transition} region discussed in Section~\ref{sec:trig_evsel}, resulting in the following \ETA\ binning: 
\begin{itemize}
    \item \pxPb : -4.5, -3.2, -1.5, 0, 1.5, 3.1;
    \item \Pbxp : -3.1, -1.5, 0, 1.5, 3.2, 4.5;
\end{itemize}
Please note that this binning was adopted only for the sake of the jet performance studies reported in this section. For \pxA\ collisions, the centrality dependence of the jet reconstruction performance was also investigated. Three centrality bins, namely \textcolor{black}{
0--10\%, 10--60\% 
}
and 60--90\%, were defined for the sake of these studies, analogously to what is done for the analysis.

The jet reconstruction efficiency was evaluated using the \texttt{Pythia8} MC samples listed in Section~\ref{sec:data} by determining how often a reconstructed jet is associated with a truth one. 
The results, binned in terms of event centrality, were obtained for both \pxPb\ and \Pbxp\ data taking periods and are reported in Figures \ref{fig:pPb_jetReco} and \ref{fig:Pbp_jetReco}, respectively. In both cases, a mild dependence on \ETA\ and centrality can be noted, with the 99\% efficiency threshold varying by a few GeV between the forward-backward regions and central and peripheral collisions. 


\begin{figure}[ht]
    \centering
    \includegraphics[width=0.325\linewidth,trim={0cm 0cm 0cm 0mm},clip]{sections/JetPerformance/fig/pPb_jetRecoEfficiency/jetRecoEff_Cent0_QM2023_20230604.pdf}
     \includegraphics[width=0.325\linewidth,trim={0cm 0cm 0cm 0mm},clip]{sections/JetPerformance/fig/pPb_jetRecoEfficiency/jetRecoEff_Cent1_QM2023_20230604.pdf}
    \includegraphics[width=0.325\linewidth,trim={0cm 0cm 0cm 0mm},clip]{sections/JetPerformance/fig/pPb_jetRecoEfficiency/jetRecoEff_Cent2_QM2023_20230604.pdf}
    \caption{
    Jet reconstruction efficiency in \pxPb\ data-taking configuration, evaluated in different \ETA\ regions of the ATLAS detector for central (0--10\%, left), mid-central (10--60\%, center) and peripheral (60--90\%, right) collisions. Corresponding 99\% efficiency \PTT\ thresholds are also reported for each region and centrality interval.
    }
    \label{fig:pPb_jetReco}
\end{figure}
\begin{figure}[ht]
    \centering
    \includegraphics[width=0.325\linewidth,trim={0cm 0cm 0cm 0mm},clip]{sections/JetPerformance/fig/Pbp_jetRecoEfficiency/jetRecoEff_Cent0_QM2023_20230604.pdf}
     \includegraphics[width=0.325\linewidth,trim={0cm 0cm 0cm 0mm},clip]{sections/JetPerformance/fig/Pbp_jetRecoEfficiency/jetRecoEff_Cent1_QM2023_20230604.pdf}
    \includegraphics[width=0.325\linewidth,trim={0cm 0cm 0cm 0mm},clip]{sections/JetPerformance/fig/Pbp_jetRecoEfficiency/jetRecoEff_Cent2_QM2023_20230604.pdf}
    \caption{
    Jet reconstruction efficiency in \Pbxp\ data-taking configuration, evaluated in different \ETA\ regions of the ATLAS detector for central (0--10\%, left), mid-central (10--60\%, center) and peripheral (60--90\%, right) collisions. Corresponding 99\% efficiency \PTT\ thresholds are also reported for each region and centrality interval.
    }
    \label{fig:Pbp_jetReco}
\end{figure}


The post-calibration jet energy scale (JES), calculated as \JES\ was studied as a function of truth jet \PT, centrality, pseudorapidity, and collision type. Closure is achieved when the ratio is equal to 1, meaning that, on average, \PTR\ matches \PTT. To quantitatively evaluate the closure, the distribution of \JES\ is constructed in each \PTT\ bin and then fit using a Gaussian function. The fit is carried out in two iterations, with the second one restricted to $\mu\pm1.5\sigma$, with $\mu$ and $\sigma$ being in this case the mean and the standard deviation obtained from the first iteration of the fit, respectively. The jet energy resolution (JER) is evaluated as the ratio $\sigma/\mu$ constructed using the parameters obtained from the second fit iteration. 

The JES and JER for jets detected in \pxPb\ and \Pbxp\ collisions are reported in Figures~\ref{fig:pPb_JERJESfit} and \ref{fig:Pbp_JERJESfit}, respectively. The top panel shows the JER, while the bottom one displays the JES closure. It should be noted that the fit procedure suffers from instability due to some asymmetry in the \PTR/\PTT\ distribution at low \PTT. In regions where good convergence of the fit can be achieved (\textit{e.g.} above \PTT$\sim$28-30 GeV), the JES closure is found to be within 3\%, with the most forward bin showing the largest non-closure values. A mild JER dependence on centrality can be noticed for jets with \PTT\ below 80 GeV. 

The two beam orientations show compatible results when compared within same the \ETA\ region of the detector. For the \Pbxp\ orientation, the statistics entering the studies in the semi-forward (1.5~<~\ETA~<~3.2) and forward (3.2~<~\ETA~<~4.5) regions are lower compared to \pxPb\ because of the rejection of the (\ETA,$\phi$) region corresponding to the HEC issue in 2016, discussed in Section~\ref{sec:hec_issue}.

A worsening of the JER can be observed in both forward directions for \PTT\ $>$ 40 GeV. This behavior was observed also in other analyses including jets in the forward region, see for instance the dijet analysis in \pxp\ at 7 TeV \cite{ATLAS:2014hvo,Baker:1328678}. An illustration of the JES and JER from that analysis in reported is Figure~\ref{fig:ppOldJES} for completeness. The observed worsening of the detector response in this region is due to a transition between two calorimeter systems with different energy responses determined by different calorimeter geometry and/or technology. This artificially increases the energy of one side of the jet with respect to the other, altering the reconstructed four-momentum \cite{ATLAS:2017bje}. 

\begin{figure}[ht]
    \centering
    \includegraphics[width=0.325\linewidth,trim={0cm 0cm 0cm 0mm},clip]{sections/JetPerformance/fig/pPb_JESJER/JESJER_010.pdf}
     \includegraphics[width=0.325\linewidth,trim={0cm 0cm 0cm 0mm},clip]{sections/JetPerformance/fig/pPb_JESJER/JESJER_1060.pdf}
    \includegraphics[width=0.325\linewidth,trim={0cm 0cm 0cm 0mm},clip]{sections/JetPerformance/fig/pPb_JESJER/JESJER_6090.pdf}
    \caption{
    Post-calibration jet energy resolution (top) and jet energy scale (bottom) evaluated in \pxPb\ collisions ($p$ traveling towards negative \ETA\,) as a function of different \ETA\,of the leading jet and of different centrality values characterizing the \pxPb\ reaction. The JES and JER values were obtained using the mean ($\mu$) and standard deviation ($\sigma$) of the $\langle \PTRE / \PTTE \rangle$ distribution evaluated by applying a Gaussian fit to it. The error bars displayed corresponds to the uncertainty resulting from the fit on each parameter. Some issue with the convergence of the fit can be noticed at high \PT{} in the forward region (\ETA > 3.2), near to the edges of the detector acceptance, where statistics start to be low.
    }
    \label{fig:pPb_JERJESfit}
\end{figure}


\begin{figure}[ht]
    \centering
    \includegraphics[width=0.325\linewidth,trim={0cm 0cm 0cm 0mm},clip]{sections/JetPerformance/fig/Pbp_JESJER/JESJER_010.pdf}
     \includegraphics[width=0.325\linewidth,trim={0cm 0cm 0cm 0mm},clip]{sections/JetPerformance/fig/Pbp_JESJER/JESJER_1060.pdf}
    \includegraphics[width=0.325\linewidth,trim={0cm 0cm 0cm 0mm},clip]{sections/JetPerformance/fig/Pbp_JESJER/JESJER_6090.pdf}
    \caption{
    Post-calibration jet energy resolution (top) and jet energy scale (bottom) evaluated in \Pbxp\ collisions ($p$ traveling towards positive \ETA) as a function of different \ETA\ of the leading jet and of different centrality values characterizing the \Pbxp\ reaction. The JES and JER values were obtained using the mean ($\mu$) and standard deviation ($\sigma$) of the $\langle \PTRE / \PTTE \rangle$ distribution evaluated by applying a Gaussian fit to it. The error bars displayed corresponds to the uncertainty resulting from the fit on each parameter. Some issue with the convergence of the fit can be noticed at high \PT{} in the forward region (\ETA > 3.2), near to the edges of the detector acceptance, where statistics start to be low.
    }
    \label{fig:Pbp_JERJESfit}
\end{figure}

\begin{figure}[ht]
    \centering
    \includegraphics[width=0.95\linewidth,trim={0cm 0cm 0cm 0mm},clip]{sections/JetPerformance/fig/pp7TeV_JES_JER.pdf}
     \caption{
     JES (left) and JER (right) in function of \PTT\ for 7 TeV \pxp\ data. The plots show results in seven different rapidity ranges. Worse performance is observed for $3.2 < |\eta| < 3.6$. Plot taken from \cite{Baker:1328678}.
     }
    \label{fig:ppOldJES}
\end{figure}


\clearpage

\clearpage{}
\clearpage{}\section{Data Analysis} 
\label{sec:analysis}

This section describes the analysis of \EZDC, \EZDCp, \ETPB , and \ETP\, as a function of the event kinematics. 


\subsection{Event Selection}
\label{sec:event_selec}
To be included in the analysis, events must be pass the relevant good run list (GRL) selection.  A new GRL for the 2016 8.16 TeV \pxPb\,data was created for this analysis to remove periods where the ZDC had an intolerable defect.  This GRL can be found at:
\begin{quote}
\tiny
    \texttt{GoodRunsLists/data16\_hip/20240611/data16\_hip8TeV.periodAllYear\_HEAD\_Unknown\_PHYS\_HeavyIonP\_All\_Good\_WITH\_IGNORES.xml}. 
\end{quote}

The construction of jet spectra for the three different trigger schemes used during the \pxPb\ data taking used in this analysis is the same as detailed in \cite{Longo:2860813}. Jets are reconstructed using the \antikt\ jet algorithm with a radius parameter of R~=~0.4. 
This analysis considers the measured leading dijet pair constructed from the two highest \PT{} jets in the event with reconstructed $p_{\TN{T,1}}>40$ GeV, $p_{\TN{T,2}}>30$ GeV.
Both jets are required to pass the LooseBad jet cleaning cut and the event is required to have a primary vertex. A dijet event is considered if the leading jet fired the lowest prescale, fully-efficient (e.g. with efficiency > 99\%) trigger corresponding to the jet \ETA\,and \PT. 
The measured leading dijet pair is considered if both jets are within -4.5~$<$~\ETA~$<$~2.8 (-2.8~$<$~\ETA~$<$~4.5) for the \pxPb\ (\Pbxp) beam orientation. Pseudorapidity cuts are included for both beam orientations, but this note will concentrate on the \pxPb\ side due to the missing module in the ZDC as mentioned in Sec.~\ref{sec:missing_zdc_mod}. The p-going \ETA\ cut is set to avoid considering jets which are not fully contained in the FCal acceptance. The Pb-going \ETA\ cut is imposed by the absence of jet triggers in the backward region during the measurement and is selected to avoid jets biasing the centrality determination in the Pb-going arm of the FCal. 

Events with either the leading or the sub-leading jet reconstructed in the HEC issue region (see Section~\ref{sec:hec_issue}) are not considered in the analysis. Events entering the analysis must have at least two such jets which meet the above criteria, allowing for construction of the initial state kinematic variables defined by Equation~\ref{eq:xp} and Equation~\ref{eq:xPb}.

To have a meaningful determination of per-event centrality, it is essential to analyze events where only one hadronic \pxPb\ collision occurred. To meet this requirement, an effective pile-up rejection strategy is needed, in particular for data with high $\mu$ (by Heavy Ion standards) like those analyzed in this note ($\mu$ ranging from 0.15 to 0.3 in most of the runs). The same pile-up rejection strategy described in \cite{Hill:2301540}, specifically established for this data-set, is chosen for this analysis. Events characterized by more than one vertex with 7 or more associated tracks are rejected at the level of event selection. By applying this pile-up rejection, residual pile-up contributions can be still noticed in the ZDC, see Figure~\ref{fig:result_corr_mb_ppb} where the correlation between \EZDC\ and \ETPB\ is demonstrated. 

\begin{figure}[ht]\centering
\begin{subfigure}{.49\linewidth}
\includegraphics[width=\linewidth]{sections/Results/fig/minbias/ZdcFCal_bin_log_HLT_mb_sptrk_L1MBTS_1.pdf}
\end{subfigure}
\begin{subfigure}{.49\linewidth}
\includegraphics[width=\linewidth]{sections/Results/fig/minbias/ZdcFCal_bin_lin_HLT_mb_sptrk_L1MBTS_1.pdf}
\end{subfigure}
\caption{Correlation between the Pb going neutral forward energy (\EZDC) and the transverse energy accumulated in the Pb-going FCal arm (\ETPB) for the \pxPb\ orientation.  \ETPB\ is binned a in logarithmic (left) and linear (right) fashion. \EZDC\ is normalized to the single neutron peak in order to display energy in terms of number of neutrons.}  
\label{fig:result_corr_mb_ppb}
\end{figure}

This residual pile-up is not seen in the \ETPB\ spectrum, indicating that these backgrounds come from interactions that either only produce particles at very far forward rapidities or happen away from the IP.
This residual was found to be present even when removing all the secondary vertices, and is suspected to be related to UPC \pxPb\ events. 
A procedure to characterize and remove this additional pile-up from the ZDC distributions is discussed in Section~\ref{ssec:pup_template}. 

ATLAS Liquid Argon calorimeter~\cite{ATLAS-TDR-02}  cells receive ionization pulses resultant from the passage of charged particle tracks. These analog pulses are delivered to a Front-End Board (FEB) where they are amplified and shaped.  The characteristic shape of these FEB pulses is a large positive amplitude pulse lasting around 100 ns proceeded immediately by a smaller negative amplitude pulse lasting around 400 ns.  In the 2016 8.16 \TeV\ \pxPb\ dataset, bunches of colliding particles circulate with 100 ns spacing between them. The measurement of \ETPB\ is then biased negatively in events with preceding bunches having large event activity. Due to the physical constraints of transmission lines and the ZDCs location far from the interaction point, pulses seen in the ZDC have long characteristic decay time.  Therefore, by sampling the ZDCs ADC value just before a pulse arrives and comparing that to a common baseline it is possible to identify these events with preceding large event activity.  More details on this selection are given in Section~\ref{sec:oot_zdc}.


Due to the difference in Z between the Pb ion (82) and the proton (1), most of the events from UPC \pxPb\ collisions are expected to involve the emission of a virtual photon from the Pb ion to the $p$, with subsequent low activity on the Pb-going side of the FCal. By applying a 0--90\% selection for the per-event dijet yield analysis, one can expect a substantial fraction of these events to be effectively rejected.  Diffractive events are associated with significant rapidity gaps in the particle production. To further enforce the UPC events rejection, a cut based on rapidity gaps, equal to one utilized in the centrality analysis \cite{Hill:2301540}, where events with $\Delta \eta_\mathrm{gap}^\mathrm{Pb} > 1.8$ were rejected.  In this analysis, such a cut rejects only a handful of events in the 0-90\% centrality class, as demonstrated by Figure~\ref{fig:gap_ana}.  



Unfolding of experimental effects in both the central ATLAS detector and the ZDC is non-trivial due to absence of a comprehensive MC simulation of both hard-scattering and nuclear breakup. In order to unfold for jet-related effects without losing the correlation with the ZDC energy, a 2D unfolding approach in \xp\ and the calorimetric energy of interest (\EZDC,\EZDCp,\ETPB\ and \ETP) was chosen, with no experimental effects unfolded for the second, given the lack of adequate MC tools for both the calorimeters. More discussion on the unfolding can be found in Section~\ref{ssec:unf_strat}.




\begin{figure}
\centering
\begin{minipage}{.42\textwidth}
  \centering
  \includegraphics[width=.8\linewidth]{sections/Analysis/fig/gapcut/Gap1D.png}
\end{minipage}\begin{minipage}{.58\textwidth}
  \centering
  \includegraphics[width=.9\linewidth]{sections/Analysis/fig/gapcut/Gap2D.png}
\end{minipage}
\begin{minipage}{.5\textwidth}
  \centering
  \includegraphics[width=.9\linewidth]{sections/Analysis/fig/gapcut/gap_zdc_mb_HLT_mb_sptrk_L1MBTS_1.pdf}
\end{minipage}
\caption{ 
(left top) 1-dimensional nucleus-going rapidity gap distribution. The red dotted line shows the cut applied in the analysis. (right top) Correlation between \ETPB\ and $\Delta \eta_\mathrm{gaps}^\mathrm{Pb}$. (bottom) Correlation between \EZDC\ and $\Delta \eta_\mathrm{gaps}^\mathrm{Pb}$. The vertical (horizontal) red dotted line represents the cut applied in the analysis for the rapidity gap (centrality) selection.  
}
\label{fig:gap_ana}
\end{figure}
\clearpage
\subsection{ZDC Analysis}
\label{sec:zdc_calib}
The ZDC samples each neutron with a certain resolution; therefore, the calorimeter response to events with multiple neutrons is approximately the sum of normally distributed random single neutrons.  The sum of $X$ normally distributed random variables with mean $\mu$ and variance $\sigma^2$ is distributed as $\mathcal{N}\left(X\mu,X\sigma^2\right)$.  The single neutron peak mean and variance are extracted from calibrated data on a run by run basis. 
Two types of fits are done for each run, they are as follows.
\begin{enumerate}
    \item Simultaneous fit to multiple neutron peaks (Equation \ref{eqn:1nGaussFits})
    \item Simultaneous fit to multiple neutron peaks + background parameterization (Equations \ref{eqn:1nGaussFits} + \ref{eqn:1nBKGDFits})
\end{enumerate}

\begin{equation}
    \label{eqn:1nGaussFits}
    f_n(E) = \sum_{n=1}^{6} \frac{A_n}{\sqrt{2\pi n\sigma}} \exp{\frac{\left(E-n\mu\right)^2}{2n\sigma^2}}
\end{equation}
\begin{equation}
    \label{eqn:1nBKGDFits}
    f_{\mathrm{bkgd}}(E) = \frac{\left(B_1-B_3E\right)E}{E+B_2}
\end{equation}

The background term in the ZDC fits is necessary for the fit to converge well around points in the immediate vicinity of the one neutron peak. The functional form of Equation~\ref{eqn:1nBKGDFits} was chosen so that it goes to zero with energy and decreases towards the high edge of this plot. Note that the relative amplitudes of successive neutron peaks is certainly not physical when using the background fit; however, the extracted $\mu_{1n}$ and $\sigma_{1n}$, which enter the analysis, are clearly more accurate. Figure~\ref{fig:result_fits_313100} shows the result of both of these fitting methods. For a tabulation of the fit results for all runs in the 2016 \pxPb\ dataset, see the Appendix of \cite{Hoppesch:2900504}. 
\begin{figure}[ht]\centering
\begin{subfigure}{.49\linewidth}
\includegraphics[width=\linewidth]{sections/Analysis/fig/zdccalib/Run313100v16BKGD.pdf}
\end{subfigure}
\begin{subfigure}{.49\linewidth}
\includegraphics[width=\linewidth]{sections/Analysis/fig/zdccalib/Run313100v16.pdf}
\end{subfigure}
\caption{Run 313100 neutron energy distribution with (left) and without (right) a background term}
\label{fig:result_fits_313100}
\end{figure}

\subsubsection{Calibration of proton side ZDC}
Due to the difference in energy scales between \EZDC and \EZDCp, the modules facing the Pb fragments were operated at a lower High Voltage (HV) than the modules facing the proton fragments.  Since, the energy scale in the ZDC is set by the single neutron peak, there is no way to set an energy scale on the proton going side from only proton data.  In order to calibrate the proton going ZDC in the \pxPb\ orientation, the first run of the \Pbxp\ orientation (313572) had the Pb going ZDC configured with the proton HV from the previous \pxPb\ orientation.  This strategy was employed specifically to calibrate the proton going side, as the Pb going ZDC had a very low saturation point during that run. This allows a setting of an absolute scale for the proton going ZDC for the \pxPb\ run. Unfortunately, this special run was only done for one beam orientation, which is why there can only be calibrated proton going ZDC energies for the \pxPb orientation.  


\subsection{ZDC Variance scaling with Energy}
\label{sec:zdc_mc_reso_studies}
The assumption that \EZDC\ of an event with $X$ spectator neutrons can be modeled as a sum of normally distributed beam energy neutrons is evaluated by using a toy Monte Carlo simulation.  This involved the production of three samples
\begin{itemize}
    \item 500k single neutron events with beam energy of 2.56 \TeV\
    \item 100k five neutron events with a beam energy of 2.56 \TeV\
    \item 50k ten neutron events with a beam energy of 2.56 \TeV\
\end{itemize}
As explained in Section~\ref{sec:zdc_mc_event_gen_issues} the neutrons are given a nuclear Fermi momentum \cite{Moniz:1971mt} randomly drawn from the Fermi sphere.  The procedure for analyzing the neutron variance is as follows,
\begin{enumerate}
    \item Extract a single neutron standard deviation, $\sigma_{1n}^{R}$, from the single neutron sample by fitting the reconstructed energy to a normal distribution.
    \item Fabricate a convolved five (ten) neutron spectrum from the summed energy of five (ten) \textbf{i.i.d.} single neutron events.  
\item Extract a five (ten) neutron standard deviation, $\sigma_{5n}^{R}$  $\left(\sigma_{10n}^{R}\right)$, from the real five (ten) neutron sample.  Verify that $\sigma_{5n}^{R} = \sigma_{\sum_i^5;1n}^{R}$ $\left(\sigma_{10n}^{R} = \sigma_{\sum_i^{10};1n}^{R}\right)$ 
\end{enumerate}

\begin{table}[h]
\centering
\begin{tabular}{|c|c|c|c|c|}
\hline
$\sigma_{1n}^{R} $ & $\sigma_{\sum_i^5;1n}^{R}$  & $\sigma_{5n}^{R}$ & $\sigma_{\sum_i^{10};1n}^{R}$  & $\sigma_{10n}^{R}$ \\
\hline
\hline
0.546 $\pm$~0.0021 & 1.289 $\pm$~0.0074 & 1.280 $\pm$~0.0069 & 1.906 $\pm$~0.011 & 1.861 $\pm$~0.014 \\
\hline
\end{tabular}
\caption{Standard deviation for each ZDC Toy MC Sample with fit parameter uncertainties.}
\label{tab:std_dev_toy_mc}
\end{table}

The results of this study can be found in Figure~\ref{fig:analysis_zdc_MC_study} and the variances are tabulated in Table~\ref{tab:std_dev_toy_mc}. This study indicates that the variance of the convolved sample is slightly larger than the true sample varience, and therefore that estimating the variance on N neutrons as a convolution of single neutrons gives a conservative estimate for the ZDC resolution above the single neutron peak. This study is a work in progress as the ZDC Monte Carlo simulation is under active development and validation.

\begin{figure}[ht]\centering
\begin{subfigure}[c]{.49\linewidth}
\includegraphics[width=\linewidth]{sections/Analysis/fig/zdcMC/ToyMC_Reco_SideC_5neutrons.pdf}
\end{subfigure}
\begin{subfigure}[c]{.49\linewidth}
\includegraphics[width=\linewidth]{sections/Analysis/fig/zdcMC/ToyMC_Truth_SideC_5neutrons.pdf}
\end{subfigure}
\begin{subfigure}[c]{.49\linewidth}
\includegraphics[width=\linewidth]{sections/Analysis/fig/zdcMC/ToyMC_Reco_SideC_10neutrons.pdf}
\end{subfigure}
\begin{subfigure}[c]{.49\linewidth}
\includegraphics[width=\linewidth]{sections/Analysis/fig/zdcMC/ToyMC_Truth_SideC_10neutrons.pdf}
\end{subfigure}
\caption{Reconstructed \EZDC\ spectrum for five neutron Monte Carlo sample (top left) and ten neutron Monte Carlo (bottom left). The black spectrum is the single neutron reconstructed energy.  The blue spectrum is the convolved sample, fabricated from single neutron events. The red spectrum is the sample. The truth neutron energy spectrum between the convolved sample and the real sample are consistent for both the five (top right) and ten (bottom right) neutron productions.}
\label{fig:analysis_zdc_MC_study}
\end{figure}
\clearpage
\subsection{Pileup Template Fit}
\label{ssec:pup_template}
A template fitting method is used to subtract a residual pileup contribution (see Figure~\ref{fig:zdc_spectrum_nominal_pileup_cut}) to the \EZDC spectrum.  

\begin{figure}
    \centering
    \includegraphics[width=0.4\linewidth]{sections//Analysis//fig//zdcpileup/ZdcSpectrum_DJSample.pdf}
    \caption{Distribution of the \EZDC\ with a dijet Selection.  The tail going towards large \EZDC\ is indicative of residual pileup (i.e. pileup not removed by the nominal pileup selection, see section \ref{sec:event_selec}). The residual pileup contribution remains even with a stricter vertex based pileup cut (e.g. requiring only one vertex in the event). The inset shows a zoomed view of the single neutron peak.}
    \label{fig:zdc_spectrum_nominal_pileup_cut}
\end{figure}

The template fitting procedure draws inspiration from a prior note, \cite{ATL-COM-PHYS-2022-1138}.  Events in the \EZDC\ spectrum are characterized by having either one interaction (`Good') or multiple interactions (`Pileup') in a bunch crossing.  Thus the spectrum can be decomposed as,
\begin{equation}
    \label{eqn:pileup_spectrum_sum}
    \left(\frac{\mathrm{dN}}{\mathrm{d}\EEZDC}\right)_{\mathrm{Total}} = \left(\frac{\mathrm{dN}}{\mathrm{d}\EEZDC}\right)_{\mathrm{Good}} + \left(\frac{\mathrm{dN}}{\mathrm{d}\EEZDC}\right)_{\mathrm{Pileup}} 
\end{equation}

The pileup contribution in Equation~\ref{eqn:pileup_spectrum_sum} is estimated by considering the convolution of events in the `Good' distribution.  For every X pileup interaction added to an event, the event is weighed by the Poisson probability of seeing X extra events in the same bunch crossing.  The contribution from pileup events containing one, two, and three additional interactions is estimated as follows.
\begin{equation}
    \label{eqn:pileup_convolution_weights}
    \begin{split}  
    \left(\frac{\mathrm{dN}}{\mathrm{d}\EEZDC}\right)_{\mathrm{Pileup}} & = \left[\left(\frac{\mathrm{dN}}{\mathrm{d}\EEZDC}\right)_{\mathrm{Good}} \circledast \left(\frac{\mathrm{dN}}{\mathrm{d}\EEZDC}\right)_{\mathrm{Good}}\right] \times \langle \mu  e^{- \mu } \rangle \\
    \left(\frac{\mathrm{dN}}{\mathrm{d}\EEZDC}\right)_{\mathrm{Pileup}} & = \left[\left(\frac{\mathrm{dN}}{\mathrm{d}\EEZDC}\right)_{\mathrm{Good}} \circledast \left(\frac{\mathrm{dN}}{\mathrm{d}\EEZDC}\right)_{\mathrm{Good}}\circledast \left(\frac{\mathrm{dN}}{\mathrm{d}\EEZDC}\right)_{\mathrm{Good}}\right] \times \langle \frac{\mu^2 e^{ -\mu }}{2}\rangle \\
    \left(\frac{\mathrm{dN}}{\mathrm{d}\EEZDC}\right)_{\mathrm{Pileup}} & = \left[\left(\frac{\mathrm{dN}}{\mathrm{d}\EEZDC}\right)_{\mathrm{Good}} \circledast \left(\frac{\mathrm{dN}}{\mathrm{d}\EEZDC}\right)_{\mathrm{Good}} \circledast \left(\frac{\mathrm{dN}}{\mathrm{d}\EEZDC}\right)_{\mathrm{Good}} \circledast \left(\frac{\mathrm{dN}}{\mathrm{d}\EEZDC}\right)_{\mathrm{Good}} \right] \times \langle \frac{\mu^3 e^{ -\mu }}{6} \rangle 
    \end{split}
\end{equation}

To get an unbiased convolved spectrum, it is important to consider that the dijet selection used in this analysis has a characteristic \EZDC spectrum shape (See Figure~\ref{fig:result_zdcall_ppb} compared to Figure~\ref{fig:result_zdcSpectra_pPb}).  The first convolved term in Equation~\ref{eqn:pileup_convolution_weights} will come from the dijet \EZDC\ spectrum, while the following convolved terms will be drawn from a MinBias \EZDC\ spectrum.

To define the distribution of 'Good' events in data (i.e. the events that enter the convolution calculation), a subset of events that both pass the nominal pileup cut and have a relatively low $\mu$ are selected.  Figure~\ref{fig:pileup_zdc_mu_scan} shows how events with a relatively low $\mu$ exhibit a greatly reduced high \EZDC tail, indicating that the rate of residual pileup is low. Figure~\ref{fig:pileup_convolved_zdc} shows how the total convolved distribution is built from the composite parts in Equation~\ref{eqn:pileup_convolution_weights}.  That figure also shows how the convolved sum is compatible in shape with the spectrum of events rejected by the nominal pileup cut.

\begin{figure}[ht]\centering
\begin{subfigure}{.49\linewidth}
\includegraphics[width=\linewidth]{sections/Analysis/fig/zdcpileup/zdc_mu.pdf}
\end{subfigure}
\begin{subfigure}{.49\linewidth}
\includegraphics[width=\linewidth]{sections/Analysis/fig/zdcpileup/zdc_mu_nominal.pdf}
\end{subfigure}
\caption{Distribution of \EZDC\ with a dijet selection with bins of $\mu$. All events (left) and only events that pass the nominal pileup cut (right).}
\label{fig:pileup_zdc_mu_scan}
\end{figure}

\begin{figure}[ht]\centering
\begin{subfigure}{.49\linewidth}
\includegraphics[width=\linewidth]{sections/Analysis/fig/zdcpileup/zdc_convolved_dist.pdf}
\end{subfigure}
\begin{subfigure}{.49\linewidth}
\includegraphics[width=\linewidth]{sections/Analysis/fig/zdcpileup/zdc_convolved_comp_nominal_dist.pdf}
\end{subfigure}
\caption{Distribution of convolved pileup distribution (left) for two, three, and four interaction pileup contributions separately and summed. Convolved pileup distribution plotted along with events that failed the nominal pileup cut. (right) Note that the compatible shape between the two distributions.}
\label{fig:pileup_convolved_zdc}
\end{figure}

With the convolved distribution from Figure~\ref{fig:pileup_convolved_zdc} it is possible to template fit the residual contribution in data.  The template fitting procedure involves fitting the \EZDC\ spectrum to a weighted sum of a self-normalized non pileup (good) spectrum, and a self-normalized pileup spectrum (convolved).  The functional form of this template is given in Equation~\ref{eqn:template_fit}.  The template has two free parameters
\begin{enumerate}
    \item \textbf{A}: sets the overall scale of the fit to equal the data
    \item \textbf{p}: controls the relative balance of the good and pileup distribution.  Defined such that $0\leq \mathrm{p} \leq 1$. 
\end{enumerate}
\begin{equation}
    \label{eqn:template_fit}
    \left(\frac{\mathrm{dN}}{\mathrm{d}\EEZDC}\right)_{\mathrm{Total}} = A\times\left(p\times\frac{1}{\mathrm{N_{Good}}}\times\left(\frac{\mathrm{dN}}{\mathrm{d}\EEZDC}\right)_{\mathrm{Good}} + (1-p)\times\frac{1}{\mathrm{N_{Pileup}}}\times\left(\frac{\mathrm{dN}}{\mathrm{d}\EEZDC}\right)_{\mathrm{Pileup}}\right)
\end{equation}

Once the distribution is fit with the template fit, a pileup-free distribution is obtained by subtracting the pileup template distribution scaled by an amplitude of $\mathrm{(1-p)\times A}$.  Bins that are over subtracted (i.e. give negative values in the \EZDC\ spectrum) by the template fit are set to zero. Figure~\ref{fig:pileup_template_fit_example} shows an example of an \EZDC\ spectrum before and after a template fit. The residual pileup in this sample makes up roughly 0.01\% of the sample. Finally, a statistical cut is implemented to remove bins that fluctuate above the template fit.

\begin{figure}
    \centering
    \includegraphics[width=0.7\linewidth]{sections//Analysis//fig//zdcpileup/zdc_template_fit_example.pdf}
    \caption{Distribution of \EZDC before (blue) and after (green) the template fitting procedure.  The black colored line is the template fit, given in eqn. \ref{eqn:template_fit}.  The fit parameters converge to p = 0.99073 \& A = 1.6e8  }
    \label{fig:pileup_template_fit_example}
\end{figure}

The impact of the residual pile-up removal from the ZDC on $\EAZDC$ was also scrutinized. These studies are reported in the Auxiliary Material, Section~\ref{aux:zdc_ref}.

\subsection{Study and rejection of Out of Time Pileup using the ZDC}
\label{sec:oot_zdc}
As has been seen in previous analyses, the FCal transverse energy distributions show a significant negative energy tail, see Figure~\ref{fig:fcal_pulse}, right panel. This tail can be explained in terms of out-of-time (OOT) pileup, in which the pulses from previous beam crossings overlap in the electronics. The calorimeter pulses (see Figure~\ref{fig:fcal_pulse}, left panel) are shaped in a way that cancels this effect on average, but can give rise to significant negative energy readout for events closely following others in time. 

\begin{figure}
    \centering
    \includegraphics[width=0.4\linewidth]{sections/Analysis/fig/oot_cut/FCal_pulse.png}
    \includegraphics[width=0.53\linewidth]{sections/Analysis/fig/oot_cut/fcal_raw.pdf}
    \caption{Left: The amplitude vs. time for the triangular pulse shape from the LAr calorimeter, overlaid with the
bipolar-shaped and sampled pulse shape. Taken from \cite{ATLAS-TDR-22}. Right: Distribution of FCal \ETPB\ in MB events in Run 313259. A clear negative tail can be observed.}
    \label{fig:fcal_pulse}
\end{figure}

The fraction of such events is limited and further rejected by selecting events within the 0–90\% centrality range. However, OOT contributions may still distort positive energy measurements in the FCal, biasing them toward artificially lower values. Therefore, rejecting events affected by OOT pile-up is essential for accurately characterizing the FCal energy distribution. This can be achieved by leveraging OOT pile-up information from the ZDC to exclude events with significant activity recorded in the preceding interaction.

To identify such events, the first ZDC sample within the readout window, known as the \textit{pre-sample}, was analyzed to flag events with clear predecessor activity. The ZDC employs a duplicated readout system to measure up to 100 spectator neutrons in pile-up events while preserving granularity at low neutron counts for data-driven energy scale calibration. The detector signal is split and processed through separate electronic channels with high-gain (HG, ×10) and low-gain (LG, ×1) settings. To ensure unambiguous detection of significant OOT pile-up without complications from gain differences, only LG data were analyzed.

Figure~\ref{fig:zdc_presample} presents the pre-sample values for each ZDC module on the Pb-going side, while Figure~\ref{fig:zdc_presample_sum} shows the sum of pre-samples across all four ZDC modules.

\begin{figure}
    \centering
    \includegraphics[width=0.75\linewidth]{sections/Analysis/fig/oot_cut/pre-sample-LG-ZDC.pdf}
    \caption{The distribution of pre-sample values for the four ZDC modules on the Pb-going side in the 2016 data.}
    \label{fig:zdc_presample}
\end{figure}
\begin{figure}
    \centering
    \includegraphics[width=0.75\linewidth]{sections/Analysis/fig/oot_cut/sum_presample.pdf}
    \caption{The sum of pre-sample values for the four ZDC modules on the Pb-going side in the 2016 data. The blue curve represents results obtained choosing only events from first bunch in trains, that are not affected by OOT pile-up. }
    \label{fig:zdc_presample_sum}
\end{figure}

The sum of the pre-samples in the Pb-going ZDC module is shown for two different selections: events occurring in first bunch in train, and events occurring in other bunches. Since the first bunches have unfilled bunches ahead of them, they are not expected to be affected by OOT contribution. Indeed, the distribution of the pre-samples for them is peaked around 400 ADCs, which roughly corresponds to the sum of the readout baselines in absence of signal. Therefore, the first bunch in train distribution can be used to derive an OOT pile-up cut based on the ZDC pre-samples. 
Given the steep edges of the first bunch in train distribution, the fine-tuning of the cut may significantly impact its performance. Therefore, different cuts were tried, examining their effectiveness in removing the negative tails of the FCal distribution. The results of this exercise are reported in Figure~\ref{fig:zdc_presample_sum_cut_impact}. The corresponding impact on the statistics selected in MB events and in the dijet sample studied in the analysis is reported in Table~\ref{tab:oot_num_impact}. The 425 cut was chosen for the analysis. {\color{red} Matt here we should write how and quote the fraction of dijet events that it removes (on top of the CONF selection}. 

\begin{figure}
    \centering
    \includegraphics[width=0.75\linewidth]{sections/Analysis/fig/oot_cut/OOT_Cut_Impact.pdf}
    \caption{MB FCal distribution obtained without (black) and with different OOT pile-up cuts applied on the ZDC pre-sample sum.}
    \label{fig:zdc_presample_sum_cut_impact}
\end{figure}

\begin{table}
\begin{center}
\begin{tabular}{|l|c|c|} 
 \hline
 \textbf{Cut on ZDC pre-sample sum} & \textbf{\% of events removed in MB} & \textbf{\% of dijet events removed }\\ 
 \hline\hline
 400 & 24.1\% & {\color{red} Matt} \\ 
 \hline
 425 & 12.8\% & {\color{red} Matt} \\
 \hline
 450 & 9.2\% & {\color{red} Matt}\\
 \hline
\end{tabular}
\caption{Fraction of events removed by the OOT cut on the ZDC pre-sample sum for different values of the cut. }
\label{tab:oot_num_impact}
\end{center}
\end{table}

The impact of the OOT pile-up cut on the average FCal energy extracted in the analysis was also assessed. The results are reported in Figure~\ref{fig:w_wo_oot}. The impact of the cut on the average energy in the calorimeter is marginal ($\sim$1\%), and leads to a systematic shift towards higher values. This behavior is expected from the removal of events affected by negative tails affecting the FCal energy. 

\begin{figure}
    \centering
    \includegraphics[width=0.75\linewidth]{sections/Analysis/fig/oot_cut/OOT_impact_FCal.pdf}
    \caption{Distribution of \EATPB\ as presented in the preliminary CONF results \cite{ATLAS-CONF-2024-013}, without OOT pile-up cut (full markers), and with the newly developed ZDC-based OOT pile-up cut applied (open markers). }
    \label{fig:w_wo_oot}
\end{figure}











\clearpage

\subsection{Unfolding Strategy}
\label{ssec:unf_strat}
A 2D unfolding strategy has been established for this analysis. The two variables participating in the unfolding procedure are \xp\ and, in turn, the calorimeter energy of interest (\EZDC, \EZDCp, \ETPB\ and \ETP). Because of the absence of FCal truth information in the MC (most of the energy deposited in the calorimeter comes from the overlay MB event) and the lack of a neutron emission simulation in ZDC events, as well as a fully commissioned Run 2 ZDC Monte Carlo, the second variable is propagated through the unfolding procedure with fully diagonal migration matrices, effectively leaving the calorimeter distribution unaltered if not for effects due to \xp\ migration. This approach relies on the assumption that the ATLAS detector response to a jet is not correlated with the response of the ZDC to forward neutral particles, and was chosen as a compromise to avoid a fully not unfolded result, that would have left also the jet related part of the measurement without correction for experimental effects despite refined techniques available to remove those from the data. On the other hand, it's unclear what the feasibility of unfolding calorimeteric information is at the moment, and if any effort in that direction would lead to a significant improvement of the results not immediately compensated by increasing uncertainties due to all the factors listed above. 

When the calorimetric target of the unfolding is the FCal energy (\ETPB\ or \ETP), a MC input to populate the matrices is available from the overlay MC. In order to equalize the FCal distribution in the MC and in the data, affected by discrepancies described in Section~\ref{sec:centrality} and shown in Figure~\ref{fig:MB_vs_dijet}, a re-weighting of the FCal distribution is applied prior to unfolding. This step effectively reproduce the  FCal data, that are then used as both truth and reconstructed in populating the migration matrices. 

On the other hand, the Pythia8 dijet overlay MC does not include any information on nuclear breakup or the ZDC. In order to propagate the data through the unfolding distribution: 
\begin{itemize}
    \item The data passing the event selection described in Section~\ref{sec:event_selec} are used to construct \EZDC\ distributions in different \xp\ bins defined for analysis. 
    \item In case of \EZDC\, the procedure to remove residual pile-up in the ZDC, outlined in Section~\ref{ssec:pup_template}, is applied to each distribution, to obtain a pile-up free ZDC distribution in each bin. 
    \item For each dijet MC event, depending on the reconstructed value of \xp\, the corresponding pile-up free \EZDC\ distribution is used to generate the ZDC energy associated to the event. 
    \item This energy is used as both truth and reconstructed value once filling the response matrices in the 2D unfolding. 
\end{itemize}
Two relevant aspects have to be underlined about this approach, namely 1. the randomization process clearly depends on the seed used to sample each data distribution and 2. the neutrons are attached to dijet MC events generated in different JZ slices, that introduces fluctuations between the MC and the data distribution. These factors produce fluctuations in different \xp\ bins depending on the randomization. In order to determine the final \EZDC\ distribution after unfolding, the unfolding procedure is repeated for 50 replicas each using a different randomization seed. The average of the unfolded results in each kinematic bin is taken as the unfolded point, while the standard deviation is taken as the estimator of the statistical uncertainty on the data. 

Response matrices are filled using reco-truth matched dijets, allowing for calculation of $x_{p}^{\text{reco}}$ and $x_{p}^{\text{truth}}$ via the left portion of Equation~\ref{eq:xp}. 
For each dijet event, an $\text{E}_\text{ZDC}^\text{Pb-side}$ is sampled from the data distribution in a given bin of \xp, using the reco-level \xp. 
In this way, the response matrix for each centrality interval is populated with ($x_{p}^{\text{reco}}$,~$x_{p}^{\text{truth}}$,~$\text{E}_\text{ZDC,reco}^\text{Pb-side}$,~$\text{E}_\text{ZDC,truth}^\text{Pb-side}$) entries. A ``projection'' of the 4D response matrix is shown in Figure~\ref{fig:pPb_resp0}, where the major axes correspond to $\text{E}_\text{ZDC}^\text{Pb-side}$ and the minor axes represent \xp. Refer to Section~\ref{sec:binning} for the binning used to construct each axis. Figure~\ref{fig:pPb_respInt} represents the object in Figure~\ref{fig:pPb_resp0} integrated over every bin in \EZDC, in order to better see the off diagonal elements.




\begin{figure}
    \centering
    \includegraphics[width=0.75\linewidth]{sections/Analysis/fig/responsematrix/avg_resp.png}
    \caption{2D Unfolding response matrix for $0-90\%$ \pxPb\ collisions.}
    \label{fig:pPb_resp0}
\end{figure}

In order to verify that the unfolding is able to correct for detector effects and extract the underlying truth information, an unfolding closure test was performed. The unfolding closure was examined by applying the full unfolding procedure to the reconstructed MC and comparing to the truth distributions. This check, referred to as the `technical-closure' test, is used to ensure that no significant flaws exist within the unfolding portion of the analysis procedure. The results of this study for $0-90\%$ collisions are shown in Figures~\ref{fig:nc_example0}. As expected, the closure is found to be exactly unity.

\begin{figure}
    \centering
    \includegraphics[width=0.75\linewidth]{sections/Analysis/fig/responsematrix/1d_resp.pdf}
    \caption{2D Unfolding response matrix integrated over each \EZDC\ bin for $0-90\%$ \pxPb\ collisions.}
    \label{fig:pPb_respInt}
\end{figure}

\begin{figure}
    \centering
    \includegraphics[width=0.75\linewidth]{sections/Analysis/fig/nonclosure/ncfit_pta_cent[0]yb[0]ys[0]_nRep0_nIter0_trial4.pdf}
    \caption{The technical-closure test in $0-90\%$ \pxPb\ collisions. A pol0 fit (red-dashed line) is shown indicating that closure is achieved.}
    \label{fig:nc_example0}
\end{figure}

\newpage
\subsection{Unfolding Iterations} 
\label{sec:convergence}
Bayesian unfolding is an iterative procedure, and therefore, the optimal number of iterations must be decided. The two contributing factors to this decision are the convergence of the unfolding ($\sigma_{\text{conv}}$), and the growth of the statistical uncertainty ($\sigma_{\text{stat}}$). 
The convergence is quantified as the sum over all bins of the difference between a bins value at iteration, $i$, and the value at the previous iteration, $i-1$:
\begin{equation}
    \sigma^{\text{conv}}_i = 
    \sqrt{
   \sum_{k}^{N_{\text{bins}}}  
    \big( 
    N_{i,k}  - N_{i-1,k}
    \big)^2
    }
\end{equation}
where $N_{i,k}$ is the value of a given bin $k$, at iteration $i$. Contrary to the statistical component, this contribution will decrease as a function of iterations, as successive unfoldings converge to the ``truth'' distribution.
The statistical uncertainty for the $i$-th iteration is summed over all bins, 
\begin{equation}
     \sigma^{\text{stat}}_i = 
    \sqrt{  \sum_{k}^{N_{\text{bins}}}  
    \big( 
    \sigma^{\text{stat}}_{i,k}
    \big)^2 . }
\end{equation}



The quadrature sum of the two components at each iteration ($\sigma^{\text{total}}_i$) is computed,
\begin{equation}
    \sigma^{\text{total}}_i = \sqrt{  (\sigma^{\text{conv}}_i)^2  +    (\sigma^{\text{stat}}_i)^2      } ,
\end{equation}
and the iteration that produces the minimum value is chosen. 
The results of this exercise are shown in Figure~\ref{fig:unfold_iter}, with two unfolding iterations being chosen for this analysis.  The convergence term cannot be determined for only one unfolding iteration because there is no previous iteration. If we compute the convergence term with the raw data being the previous iteration for the first unfolding, the first convergence term would be off the scale of this plot.  
\begin{figure}
    \centering
    \includegraphics[width=0.85\linewidth]{sections/Analysis/fig/nonclosure/zdc_xp_ratiotest.pdf}
    \caption{Unfolding iteration optimization for $0-90\%$ \pxPb\ collisions.  Please note that the convergence has been scaled by a factor of 10 to be visualized more easily.  The limited bin-to-bin migration is expected from Figure~\ref{fig:unfresult_zdc_xpbin}, where one can see a good agreement between unfolded and not-unfolded results.}
    \label{fig:unfold_iter}
\end{figure}



\subsection{Statistical Uncertainties} 
Statistical uncertainties for the unfolding come from both the reconstructed data and the response matrices. A bootstrapping method with 100 toys is used in order to ensure statistical significance of the results. See \cite{ATL-PHYS-PUB-2021-011} for a detailed description of both the method and the ROOT \texttt{TH1Bootstrap} class. The weights for each event in a given toy are sampled from a Poisson distribution with a mean of one, producing statistically correlated samples. 

The data contribution to the statistical uncertainty is evaluated as follows:
\begin{enumerate}
    \item Nominal data spectrum is smeared, using Poisson weights, producing 100 toy distributions. 
    \item RMS of the toy distributions is calculated and is taken as the error of the nominal data spectrum. 
    \item Nominal data spectrum is unfolded using the nominal response matrix.
    \item Error on the unfolded result is taken as the data contribution to the statistical uncertainty.
\end{enumerate}
The MC contribution to the statistical uncertainty is evaluated as follows: 
\begin{enumerate}
    \item 100 response matrices are produced with weights obtained by sampling from a Poisson distribution with a mean of one.
    \item Nominal data spectrum is unfolded 100 times, using the smeared response matrices. 
    \item RMS of the 100 unfolded data spectra is taken as the MC contribution to the statistical uncertainty.
\end{enumerate}
The final statistical uncertainty is calculated by taking the quadrature sum of the bootstrapped data and MC statistical uncertainties. 


\clearpage{}
\clearpage{}\section{Systematic Uncertainties} 
\label{sec:systematics}
The main sources of systematics that were scrutinized for the analysis are related to: 
\begin{itemize}
    \item Jet Energy Scale (baseline, cross-calib, flavor)
    \item Jet Energy Resolution
    \item Unfolding prior (for both jets and calorimetric information)
    \item HEC rejection
\end{itemize}

Systematics have been assessed on both the unfolded calorimetric energy distributions (ZDC and FCal) and their means. In the first case, the analysis procedure is repeated applying changes only to a given systematic source, while keeping the rest of the analysis framework untouched. The difference between a given  systematic test and the nominal analysis is taken as the systematic variation. Systematic uncertainties related to the JES and the unfolding prior are asymmetric, while those associated to the JER and the HEC exclusion are symmetric. 

An overview of all the systematic checks applied in this analysis to the unfolded data is offered below. These types of systematic sources ar similar to those investigated for \cite{Longo:2860813}. Please note that only the distribution of systematics on $\EAZDC$ and \EATPB\ are reported in the main body of this note, while distributions for all the components of the systematics on each of the \xp\ bins for \EZDC\ and \ETPB\ can be found in the Supporting Material, Sections~\ref{sec:zdc_panel_sys} and \ref{sec:fcal_panel_sys}, respectively. 


\subsection{Jet Energy Scale}
The first component of the JES uncertainty is taken from 13 \TeV\ \pxp\ JES uncertainties for EMTopo jets, as recommended by the JetEtMiss group. HI jets are first calibrated with respect to jets reconstructed using the nominal \pxp\ algorithm (EMTopo), referred to as cross-calibration. This calibrated JES is then adapted to the center-of-mass energy used in the HI analysis \cite{2016:JetCali}. These uncertainties are retrieved using the JetUncertainties package with the following configuration file:
\begin{quote}
\centering
\texttt{R4\_GlobalReduction\_SimpleJER.config} ,
\end{quote}
which corresponds to $\sim$20 JES and 8 JER nuisance parameters \cite{2018:JESJER_NP}. Variations which are not related to the analysis have been left out of their respective systemic uncertainty and are listed below:
\begin{quote}
\centering
   \texttt{JET\_EtaIntercalibration\_NonClosure\_2018data}: Only applicable to 2018 data 
\end{quote}


The second JES component is specific to heavy ion jets and collision energies (for flavor-related uncertainties) and is evaluated using the Heavy Ion Jet recommendation \cite{2018:JESuncert}. The jet flavor uncertainty has both an energy-independent (flavour response) and energy-dependent term (flavour fraction). The uncertainties are used to produce variations by shifting the jet \PT\, according to 
\begin{equation}
    p_{\text{T}}^{*,\text{reco}} = p_{\text{T}}^{\text{reco}}\big( 1 \pm U^{\text{JES}} (p_{\text{T}} , \eta) \big), 
\end{equation}
where $\eta$ accounts for the boost of the \pxPb\,system in the case of the flavor fraction. The flavor response represents a property of the detector and therefore does not require boosting. 
The relative JES systematic uncertainties broken down by individual components are displayed in Figures ~\ref{fig:zdc_jes_comps} and ~\ref{fig:fcal_jes_comps} for \AZDC\ and \EATP\, respectively. The overall systematic contribution due to the JES is within  2\% in most of the kinematic bins, except from the edges of the \xp\ phase space. 

\begin{figure}
    \centering
    \includegraphics[width=1\linewidth]{sections//SupportMaterial//fig//syst_comps/systematic[JET]_zdc_comps.pdf}
    \caption{The components of the JES systematic uncertainty on the \AZDC\ 
a function of \xp\ for \SNN~=~8.16~\TeV\ collisions. The black dashed line shows the total JES systematic uncertainty.}
    \label{fig:zdc_jes_comps}
\end{figure}
\begin{figure}
    \centering
    \includegraphics[width=1\linewidth]{sections//SupportMaterial//fig//syst_comps/systematic[JET]_fcal_comps.pdf}
        \caption{The components of the JES systematic uncertainty on the \EATPB\ 
    as a function of \xp\ for \SNN~=~8.16~\TeV\ collisions. The black dashed line shows the total JES systematic uncertainty.}
    \label{fig:fcal_jes_comps}
\end{figure}

\subsection{Jet Energy Resolution}

The JER systematic accounts for differences between the jet energy resolution in data and MC, where the fractional resolution, $\sigma_{\text{JER}}^{\text{syst}}$, is retrieved as a function of reco jet \PT\ and $\eta$ using the JetUncertainties tool~\cite{2018:JERuncert}. Unfolding is repeated with a modified response matrix for each JER uncertainty component, where the jet \PT\ has been smeared by 
\begin{equation}
    p_{\text{T}}^{*,\text{reco}} = p_{\text{T}}^{\text{reco}}\times\cal{N}\big(\text{1}, \sigma_{\text{JER}}^{\text{eff}}\big), 
\end{equation}
where $\cal{N}\big(\text{1}, \sigma_{\text{JER}}^{\text{eff}}\big)$ is the normal distribution with the effective resolution
\begin{equation}
\sigma_{\text{JER}}^{\text{eff}}
= \sqrt{ \big( \sigma_{\text{JER}} + \sigma_{\text{JER}}^{\text{syst}} \big)^2 - \sigma_{\text{JER}}^2  }.
\end{equation}

In addition to the 8 JER nuissance parameters produced using the SimpleJER configuration, a heavy ion specific JER uncertainty corresponding to HI jet collections and the cross-calibration was added to the uncertainties in \pxPb\,collisions.  

The relative JER systematic uncertainties broken down by individual components are displayed in Figures ~\ref{fig:zdc_jer_comps}~and~\ref{fig:fcal_jer_comps} for \AZDC\ and \EATP\ respectively. The overall systematic contribution due to the JER is within 1\% in most of the kinematic bins, except from the high \xp\ region, where it raises up to 3\%, and the low \xp\ bins, where it reaches 3-5\%.  

\begin{figure}[ht]
    \centering
    \includegraphics[width=1\linewidth]{sections//SupportMaterial//fig//syst_comps/systematic[JER]_zdc_comps.pdf}
    \caption{The components of the JER systematic uncertainty on the \AZDC\ 
a function of \xp\ for \SNN~=~8.16~\TeV\ collisions. The black dashed line shows the total JER systematic uncertainty.}
    \label{fig:zdc_jer_comps}
\end{figure}
\begin{figure}[ht]
    \centering
    \includegraphics[width=1\linewidth]{sections//SupportMaterial//fig//syst_comps/systematic[JER]_fcal_comps.pdf}
    \caption{The components of the JER systematic uncertainty on the \EATPB\ 
a function of \xp\ for \SNN~=~8.16~\TeV\ collisions. The black dashed line shows the total JER systematic uncertainty.}
    \label{fig:fcal_jer_comps}
\end{figure}


\subsection{Systematic Uncertainty Due to the Unfolding Procedure}
\label{sec:prior}
The systematic uncertainty on the unfolding is related to the sensitivity of the unfolding procedure to the choice of the input distributions, referred to as ``prior''. To determine the uncertainty on the re-weighting, the unfolding was performed using no data/MC weight. 
The analysis procedure was then repeated using the un-weighted response matrices, with the difference from the nominal result being taken as a systematic uncertainty, and will account for both jet and calorimeteric part of the uncertainty. 
Regarding the second, even though the data distribution are mimicked in the MC and passed through the unfolding with a diagonal response matrix, residual discrepancies between the distributions in the MC and the data have to be accounted for at level of systematics. In the case of the ZDC, residual discrepancies in each iteration of the unfolding derives from the random sampling of the real data \EZDC\ distribution in each kinematic bin convoluted with the JZ weights of the dijet MC sample, while for the FCal, residual differences from the data distribution may remain after the re-weighting the FCal distribution in MC to match the data. 


\subsection{Systematic Uncertainty due to Disabled HEC Quadrant}
Reconstructed level jets that fall within the region covered by the disabled $+\eta$ HEC have been removed from the the analysis (see Section~\ref{sec:hec_issue}). A systematic related to this removal is included in the analysis, and is evaluated by increasing the exclusion region by 0.1 in all directions in azimuth and pseudorapidity, and repeating the analysis procedure. The difference between the nominal exclusion and increased exclusion is taken as a systematic uncertainty and is symmetrized.

\subsection{Systematic studies related to residual pileup removal}
The effect of the residual pileup subtraction on \AZDC\ in evaluated in Section~\ref{aux:zdc_ref}.  In the extreme case of no residual pileup subtraction it is found that the greatest deviation is less than 0.25\%.  Since this contribution is minimal and subdominant compared to the other systematic sources analyzed and summarized in Figure~\ref{fig:zdc_jet_syst_uncertainties}, it was decided to do not assign an additional systematic uncertainty related to this particular source. Additionally, when considering the \EZDC\ distributions themselves, taking the extreme case of no residual pileup subtraction as a source of systematic uncertainty would lead to significant systematic uncertainties on the high \EZDC\ bins by the inclusion of clear pileup events in this region.

\subsection{Total systematic uncertainty}
The total systematic uncertainty on \AZDC\ and \EATPB\ is displayed in Figures~\ref{fig:zdc_jet_syst_uncertainties} and \ref{fig:fcal_jet_syst_uncertainties}, respectively. The \AZDC\ and \EATPB\ uncertainty is below 1\% in most of the phase space, with the exception of the lowest and highest \xp\ regions, where the relative systematic uncertainty reaches 1.5\% and 5\% for \AZDC\ and \EATPB, respectively. 


\begin{figure}[ht]
    \centering
    \includegraphics[width=0.75\linewidth]{sections//Systematics//fig/systematics_zdc.pdf}
    \caption{Total systematic uncertainty on \AZDC\ (black dashed line). The different components, are shown by different colored solid lines. }
    \label{fig:zdc_jet_syst_uncertainties}
\end{figure}

\begin{figure}[ht]
    \centering
    \includegraphics[width=0.75\linewidth]{sections//Systematics//fig/systematics_fcal.pdf}
    \caption{Total systematic uncertainties on \EATPB. (black dashed line). The different components, are shown by different colored solid lines. }
    \label{fig:fcal_jet_syst_uncertainties}
\end{figure}














\clearpage{}
\ifinternalnote \include{sections/ppReference/ppReference} \else \fi 
\clearpage{}\section{Unfolded Results}
\label{sec:unf_results}
\EZDC, and \ETPB were scrutinized in events characterized by a hard scattering yielding a dijet in the final state. The detailed selection for these dijet events is discussed in Section~\ref{sec:event_selec}.




Selected results for 0-90\% \pxPb\ collisions involving \xp\ and \EZDC\ have been unfolded using the strategy described in Section~\ref{ssec:unf_strat}. In order to average out fluctuations introduced by reproduction of data-like ZDC distributions in MC, and  evaluate the statistical uncertainty on $\EAZDC$ in each \xp\ distribution, 50 independent unfolding replicas were carried out to obtain the final distribution of the unfolded \EZDC\ in each kinematic bin. The number of iterations was chosen by evaluating the convergence of the averaged unfolding matrices, as described in Section~\ref{sec:convergence}. The results for \EZDC\ are displayed in  Figure~\ref{fig:unfresult_zdc_xpbin}, where a comparison between unfolded and not unfolded distributions is also provided, to demonstrate the impact of unfolding of experimental effects on the jets on the ZDC distributions. Please note that, for the final results, only unfolded distributions will be reported. 


\begin{figure}[ht]
    \centering
    \includegraphics[width=0.99\linewidth,trim={5mm 0cm 12mm 0mm},clip]{sections/Results/fig/unfolded/zdc_Pb_unf_comp_neutronDist_xp_v12.pdf}
    \caption{Unfolded Pb-going ZDC spectrum for the \pxPb\ orientation. Each panel represents an $x_p$ bin in dijet events.  In each bin of \xp\, the \AZDC\ with uncertainties  is displayed in blue text. These values directly correspond to Figure~\ref{fig:unfresult_avgzdc}. The event kinematic is corrected for experimental effects by the 2D unfolding procedure. The horizontal and vertical shaded areas represent the bin width and systematic uncertainty, respectively.  A breakdown of the systematic uncertainties for \EZDC\ can be found in Section~\ref{sec:zdc_panel_sys}. The vertical bars represent the statistical uncertainty. Energy below the ZDC detection threshold (e.g. formally assigned as zero by ZDC reconstruction) is reported in a dedicated bin. Please note that results in each bin are divided by the bin width. This is relevant when looking at the zero energy bin compared to the rest of the distribution. 
}
    \label{fig:unfresult_zdc_xpbin}
\end{figure}


A clear shift toward lower values of the mean of the distribution can be observed with increasing \xp. The distribution of \ETPB\ is presented, in the same fashion, in Figure~\ref{fig:unfresult_fcal_xpbin}. A stronger dependence on \xp\ can immediately be appreciated, as one could expect from the results presented in \cite{ATLAS:2023zfx}. As for the ZDC case, a comparison between unfolded and not-unfolded results is provided, showing minimal changes between the two versions of the results, mostly concentrated at the edges of the \xp\ phase-space. The lowest \xp\ bin appears to be the one most affected by fluctuations. This behavior can be understood in terms of statistics related to the missing HEC sector - that is heavily impacting the population of this particular bin, introducing instabilities in the unfolding procedure. Figure~\ref{fig:unfresult_avgfcal} reports the dependence of the mean of the \xp-unfolded distributions as a function of \xp. As for the ZDC case, unfolded results tracks the not-unfolded distribution with minimal deviations at the edge of the \xp\ phase space. To better understand the correlations between \xp\ and the energy accumulated in the calorimeters, the averages of the distributions are presented in Figure~\ref{fig:unfresult_avgzdc} and Figure~\ref{fig:unfresult_avgfcal} for the ZDC and FCal, respectively. 
One can notice how both distributions seem to be sensitive to \xp, in particular in the region \xp $\gtrsim 5\cdot10^{-2}$, corresponding to the onset of effects attributed to small proton size configurations in \cite{ATLAS:2023zfx}. The energy accumulated in the Pb-going side of the ZDC decreases up to $\sim$5\% in this region, while up to 40\% decrease is observed for the FCal. 
When considering the statistical and systematic uncertainties on the highest \xp\ bin, the measurement of \AZDC\ is more than 3 standard deviations away from the nominal value of \AZDC\ in the bin centered around \xp~$\approx$~0.02.
In order to test the effectiveness of the residual pile-up subtraction described in Section~\ref{ssec:pup_template}, these results were also produced in different classes of average interaction per bunch crossing. You can find the results of these additional x-check studies in the Auxiliary Material, Section~\ref{aux:mu_dep_averages}.

Systematic uncertainties on these distributions are well contained except for the outer-most \xp\ bins, where the analysis approaches the edge of the phase space. The means themselves are plotted in Figure~\ref{fig:unfresult_avgzdc}, where the statistical uncertainty is taken as the uncertainty on the mean of the 50 averaged unfolding distribution in the given \xp\ bin. In both the cases just discussed, results obtained using the unfolding are compared to not unfolded results. As one could expect from propagating the ZDC distributions through the 2D unfolding with a diagonal response matrix, the changes in the results appear to be very tame for the unfolded distributions and minimal for the mean. 
In the latter case, the estimator, as expected, is pretty resilient to fluctuations, with the changes in the results being well below 1\%.

\begin{figure}[ht]
    \centering \includegraphics[width=0.99\linewidth,trim={5mm 0cm 12mm 0mm},clip]{sections/Results/fig/unfolded/fcal_Pb_unf_neutronDist_xp_v11.pdf}
    \caption{Unfolded \ETPB\ spectrum for the \pxPb\ orientation. Each panel represents an $x_p$ bin in dijet events. In each bin of \xp\, the \EATPB\ with uncertainties  is displayed in red text.  These values directly correspond to Figure~\ref{fig:unfresult_avgfcal}. The event kinematic is corrected for experimental effects by the 2D unfolding procedure. The horizontal and vertical shaded areas represent the bin width and systematic uncertainty, respectively. A breakdown of the systematic uncertainties for \ETPB\ can be found in Section~\ref{sec:fcal_panel_sys}. The vertical bars represent the statistical uncertainty. }
    \label{fig:unfresult_fcal_xpbin}
\end{figure}

\begin{figure}[ht]
    \centering
    \includegraphics[width=0.8\linewidth,trim={2mm 1mm 2mm 5mm},clip]{sections/Results/fig/unfolded/zdc_Pb_unf_comp_xp_v12.pdf}
    \caption{Unfolded \AZDC\ as a function of \xp\ for the \pxPb\ orientation. The non-unfolded result is shifted for easier comparison.}

        \label{fig:unfresult_avgzdc}
\end{figure}

\begin{figure}[ht]
    \centering
    \includegraphics[width=0.8\linewidth,trim={2mm 1mm 2mm 2mm},clip]{sections/Results/fig/unfolded/fcal_Pb_unf_comp_xp_v11.pdf}
    \caption{Unfolded \EATPB\ as a function of \xp\ for the \pxPb\ orientation. The non-unfolded result is shifted for easier comparison.}
    \label{fig:unfresult_avgfcal}
\end{figure}


\clearpage

Figure~\ref{fig:unf_color_effect_zdc_fcal_xp} shows the comparison between \EZDC\ and \ETPB, both normalized to their maximum. The lower panel shows the ratio of the relative change in \EATPB\ over the relative change in \AZDC\ as a function of \xp.  For \xp\ bins ~$\lesssim 2\cdot10^{-2}$ the relative change in \AZDC\ is $\ll 1$ and the ratio becomes poorly defined, as one can observe by the large error bars.  However, for the region \xp~$\gtrsim 2\cdot10^{-2}$ the ratio is rather stable around a constant value, as demonstrated by the fit displayed in the Figure. This result shows how, compared to the transverse energy at forward rapidities, the energy deposited in the ZDC in \pxPb\ collisions is only $\sim$1/6 as sensitive to the hard process kinematics.  

\begin{figure}[ht]
        \centering
        \includegraphics[width=0.8\linewidth,trim={2mm 1mm 2mm 5mm}]{sections/Results/fig/unfolded/ratio_rel_effect_comb_inv.pdf}
        \caption{Unfolded \AZDC\ as a function of \xp\ (blue) compared to Unfolded \EATPB\ as a function of \xp\ (red). Both plots are normalized to their maximum value, in order to display both quantities on the same scale.
        The bottom plot shows the ratio of the relative change in \EATPB\ over the relative change in \AZDC. 
The horizontal and vertical shaded areas represent the bin width and systematic uncertainty, respectively.}
        \label{fig:unf_color_effect_zdc_fcal_xp}
    \end{figure}


\clearpage

\begin{figure}[ht]
        \centering
            \begin{subfigure}[c]{0.95\linewidth}
            \includegraphics[width=\linewidth]{sections/Results/fig/unfolded/zdc_ratio_selectXslices_v12.pdf}
            \end{subfigure}
        \caption{Unfolded \EZDC\ distributions plotted for three selected bins of \xp: a low, moderate, and high selection. These distributions are self-normalized in order to display the change of the \EZDC\ distribution shape with \xp. The bottom plot displays the ratio between the the moderate and high \xp\ selection over the low \xp\ selection. Systematic uncertainties on the ratios are propagated treating the JER, JES and HEC contribution as correlated, and the prior as uncorrelated. The three \EZDC\ distributions are taken from Figure~\ref{fig:unfresult_zdc_xpbin}. Energy below the ZDC detection threshold (e.g. formally assigned as zero by ZDC reconstruction) is reported in a dedicated bin. Please note that results in each bin are divided by the bin width. This is relevant when looking at the zero energy bin compared to the rest of the distribution. }
        \label{fig:unf_select_xp_zdc}
    \end{figure}

To better observe the change in the shape of the of the \EZDC\ distributions with \xp\, Figure~\ref{fig:unf_select_xp_zdc} shows a comparison of \EZDC\ distributions in well separated bins of \xp. From this figure, one can clearly see a shift to lower values of the distribution obtained in the high \xp\ region, characterized also by a slight increase in the fraction of counts in the bin corresponding to zero \EZDC\ compared to intermediate and low-x regions. 

It is relevant to note that pulses under the detection threshold in each ZDC module are considered compatible with noise and assigned to zero. This translates into a threshold for minimal energy detected in the ZDC slightly below 1 TeV. For this reason, the zeroth bin in Figure~\ref{fig:unf_select_xp_zdc} is set conservatively to 0-1 TeV. Please note that each bin in the top panel of Figure~\ref{fig:unf_select_xp_zdc} is divided by its bin width, introducing the visual jump observed between the first bin and the rest of the distribution.  
Similarly, Figure~\ref{fig:unf_select_xp_fcal} shows a comparison of \ETPB\ distributions in well separated bins of \xp.  

There is a striking difference in the distribution obtained in the high \xp\ (valence) region, when compared to the moderate and low \xp\ selections.  
The systematic error bars in Figures~\ref{fig:unf_select_xp_zdc}~and~\ref{fig:unf_select_xp_fcal} contains both correlated and uncorrelated sources of systematic uncertainty.\footnote{For this exercise, the JES, JER and HEC are considered correlated, while the prior uncertainty is treated as uncorrelated.}  A ratio of the distributions at intermediate and high-\xp\ with respect to the low-\xp\ one is taken to better display the actual significance in the change of the distributions, partially canceling out correlated uncertainties.  
\begin{figure}[ht]
        \centering
        \begin{subfigure}[c]{0.95\linewidth}
        \includegraphics[width=\linewidth]{sections/Results/fig/unfolded/fcal_ratio_selectXslices_v11.pdf}
        \end{subfigure}
        \caption{Unfolded \ETPB\ distributions plotted for three selected bins of \xp: a low, moderate, and high selection. Top~figure displays the \ETPB\ distributions with the shaded error bars representing all sources of systematic uncertainty.  Bottom~figure displays the \ETPB\ distributions with the shaded error bars representing only the uncorrelated systematic uncertainties. These results are taken from Figure~\ref{fig:unfresult_fcal_xpbin}.}
        \label{fig:unf_select_xp_fcal}
    \end{figure}

\clearpage

\section{Results without unfolding}
\label{sec:results}
Below, we present results for this analysis obtained without applying the 2D unfolding procedure described in Section~\ref{ssec:unf_strat}. A few of these results (e.g. the correlation between the energy in the calorimeters) cannot be unfolded with the proposed framework, while others are reported from completeness and may be unfolded for the final paper, if deemed of interest. Please note that not unfolded distributions as a function of other event kinematics, e.g.~\xf\ and \xPb, are reported in the Auxiliary Material, Sections~\ref{supp:xf_res} and \ref{supp:xPb_results}, respectively. 

\subsection{Correlation between \texorpdfstring{\EZDC}{EZDC} and \texorpdfstring{\ETPB}{FCAlET}}
A further characterization of the correlation between the forward energy from the participants (\ETPB\ in the FCal) and from the neutrals breakup neutrons (\EZDC\ in the ZDC) can be obtained by analyzing the correlations between the two detector energy distributions in each \xp\ bin.  Given their nature, these results, presented in Figure~\ref{fig:corr_panel_plot}, are not unfolded for experimental effects and are presented without the subtraction of residual pile-up in the ZDC. To give the reader an idea for how this distribution changes in each \xp\ bin, the correlation is first reported integrated over all bins of \xp\ in Figure~\ref{fig:corr_int}.
A comparison between two kinematic bins, at low and high-\xp\ values, is reported in Figure~\ref{fig:result_corr_xp}. The ratio between the two distributions highlight a clear gradient in the correlation as a function of the proton kinematics. To further investigate this behavior, 
Figure~\ref{fig:fit_zdc_fcal} reports the \EATPB\ analyzed this time as a function of \EZDC\ in different \xp\ bins. Each distribution is fitted to a linear form to capture the underlying relation between the two energies. A striking persistent linearity of the \EATPB-\EZDC\ correlation can be noticed. Figure~\ref{fig:fit_select_zdc_fcal} overlays a the plots in Figure~\ref{fig:fit_zdc_fcal} for a few selections of \xp\, where it is clearly evident that at higher \xp\ there is a significant difference in the dependence of \EATPB\ on \EZDC.  The evolution of the linear fit slope is reported in Figure~\ref{fig:fit_slope_zdc_fcal}. One can observe a consistent reduction of the `per neutron FCal energy' in the same ranges of \xp\ where the log-linear \xp-driven dijet \RCP\ suppression was observed in \cite{ATLAS:2023zfx}. The offset term in the linear fits are shown in Figure~\ref{fig:fit_offset_zdc_fcal}.  Once can observe a slight reduction of the `FCal Energy at Zero Neutrons' in high \xp\ events, however, this is biased by the centrality cut on \ETPB.  The dependence of the linear fit slope and offset on the fit range was studied by systematically varying the high limit of the to cover $\pm$ 1 bin in \EZDC\ with respect to the nominal choice of 155~\TeV.  Systematic uncertainties associated with the fit range are reported in Figures~\ref{fig:fit_slope_zdc_fcal}~and~\ref{fig:fit_offset_zdc_fcal}, both the slope and offset are weakly dependent on this systematic. The Pearson correlation coefficient between \EZDC\ and \ETPB\ as a function of \xp\ is reported in Figure~\ref{fig:cov_xp_corr}.  A small decrease in the correlation coefficient is observed with \xp, supporting the above observation that the `per neutron FCal energy' decreases with \xp. This new data will aid in further understanding the QCD-driven dynamics at work in \pxPb\ collisions related to the proton configuration. 

\begin{figure}[ht]
    \centering
    \includegraphics[width=0.5\linewidth,trim={0mm 0cm 0mm 0mm},clip]{sections/Results/fig/corr/corr_Dijets.pdf}
    \caption{Correlation plot between \EZDC\ and \ETPB, integrated over all bins of \xp\ in this measurement}
    \label{fig:corr_int}
\end{figure}


\begin{figure}[ht]
    \centering
    \includegraphics[width=0.99\linewidth,trim={5mm 0cm 9mm 0mm},clip]{sections/Results/fig/corr/zdc_fcal_panel.pdf}
    \caption{Correlation plot between \EZDC\ and \ETPB. Each panel represents an $x_p$ bin, where \xp\ is determined starting from the kinematic of the dijet events.}
    \label{fig:corr_panel_plot}
\end{figure}

    \begin{figure}[ht]\centering
\begin{subfigure}[c]{.49\linewidth}
\includegraphics[width=\linewidth]{sections/Results/fig/corr/corr_xp_low_v5.pdf}
\end{subfigure}
\begin{subfigure}[c]{.49\linewidth}
\includegraphics[width=\linewidth]{sections/Results/fig/corr/corr_xp_high_v5.pdf}
\end{subfigure}
\begin{subfigure}[c]{.6\linewidth}
\includegraphics[width=\linewidth]{sections/Results/fig/corr/corr_xp_ratio_v8.pdf}
\end{subfigure}
\caption{Correlation between the Pb going neutral forward energy (\EZDC) and the transverse energy accumulated in the Pb-going FCal arm (\ETPB) for the \pxPb\ orientation for events with a relatively low \xp\ (top left) a relatively high \xp\ (top right) and the ratio of the two correlation plots (bottom).  The ratio plot indicates that the high \xp\ selection has significantly more events with less than 0.2 TeV of energy in the Pb-going FCal and greater than 50 TeV of energy in the Pb-going ZDC. Since these are correlation plots, it was not possible to clean the \EZDC\ energy distribution from residual pile-up. Still - that contribution is present in both the bins and therefore should not bias the comparison. }
\label{fig:result_corr_xp}
\end{figure}

\begin{figure}[ht]
    \centering
    \includegraphics[width=0.99\linewidth,trim={5mm 0cm 12mm 0mm},clip]{sections/Results/fig/corr/zdc_xp_prof_fcal_v3.pdf}
    \caption{\EATPB\ plotted as a function of \EZDC. \EATPB\ was produced by taking a profile of the correlation plots in Figure~\ref{fig:corr_panel_plot}. Each panel represents an \xp\ bin in dijet events.  The  relationships between \EATPB\ and \EZDC\ are fit to a linear functional form.  The slope and intercept of each fit are plotted in their respective panels.}
    \label{fig:fit_zdc_fcal}
\end{figure}


\begin{figure}[ht]
        \centering
        \includegraphics[width=0.8\linewidth,trim={2mm 1mm 2mm 5mm}]{sections/Results/fig/corr/zdc_fcal_corr_m_xp_xp_v3.pdf}
        \caption{Slope of linear fits shown in Figure~\ref{fig:fit_zdc_fcal} as a function of \xp.  Error bars represent the uncertainties on the fit parameters (smaller than the marker size). The shaded boxes represent the systematic dependence on the fit range.}
        \label{fig:fit_slope_zdc_fcal}
    \end{figure}
\begin{figure}[ht]
        \centering
        \includegraphics[width=0.8\linewidth,trim={2mm 1mm 2mm 5mm}]{sections/Results/fig/corr/zdc_fcal_corr_c_xp_xp_v3.pdf}
        \caption{Offset term in  linear fits shown in Figure~\ref{fig:fit_zdc_fcal} as a function of \xp. Error bars represent the uncertainties on the fit parameters. The shaded boxes represent the systematic dependence on the fit range. }
        \label{fig:fit_offset_zdc_fcal}
    \end{figure}
\begin{figure}[ht]
        \centering
        \includegraphics[width=0.8\linewidth,trim={2mm 1mm 2mm 5mm}]{sections/Results/fig/corr/zdc_fcal_corr_cov_xp_xp_v3.pdf}
        \caption{Sample Pearson correlation coefficient between \EZDC\ and \ETPB, binned as a function of \xp. Error bars represent the standard error associated to the correlation.}
        \label{fig:cov_xp_corr}
    \end{figure}

\begin{figure}[ht]
    \centering
    \includegraphics[width=0.99\linewidth,trim={0mm 0cm 6mm 0mm},clip]{sections/Results/fig/corr/zdc_xp_prof_fcal_fit_selectXslices_v3.pdf}
        \caption{\EATPB\ plotted as a function of \EZDC for three selected bins of \xp: a low, moderate, and high selection. Relation is fit to a linear form with fit parameters (and fit uncertainties) listed on the figure. These results are taken from Figure~\ref{fig:fit_zdc_fcal}.}
    \label{fig:fit_select_zdc_fcal}
\end{figure}

\clearpage









\subsection{ZDC Energy and Neutron Distributions in MB events}
Since this measurement is the first of its kind to feature the ATLAS ZDC so prominently, we decided to report not only ZDC distributions in events characterized by the presence of a hard-scattering featuring a dijet, but also ZDC distributions in MB events. In particular, all the plots in this section were obtained from analyzed events where the \texttt{HLT\_mb\_sptrk\_L1MBTS\_1} trigger was fired, without any dijet requirement. Figure \ref{fig:result_zdcall_ppb} demonstrates the distribution of forward neutral energy measured in the Pb-going arm of the ZDC during the data taking with the \pxPb\ orientation. The inset show clear peaks visible for 1, 2 and 3 neutrons. A prominent change in slope in correspondence with a high number of neutrons, likely due to events characterized by residual pile-up contribution, can be seen. The corresponding distribution in case of dijet events is reported in Figure~\ref{fig:result_zdcSpectra_pPb} as a term of comparison for Figure~\ref{fig:result_zdcall_ppb}. Please note that the distribution in dijet events is prescale corrected (since it is constructed from a mix of different jet triggers) and, therefore, affected by larger uncertainties in the residual pile-up tail where statistics is relatively poor. 
\begin{figure}[ht]\centering
\includegraphics[width=0.75\linewidth]{sections/Results/fig/minbias/ZdcSpectrum_HLT_mb_sptrk_L1MBTS_1.pdf}
\caption{Distribution of \EZDC\ with a minimum bias selection.  The inset shows a zoomed view of the single neutron peak. Please note that residual pile-up is not subtracted from this distribution, as it can be noted by the tail of events at high ZDC energies. To access information about the magnitude of this contribution in this distribution, please see Section~\ref{ssec:pup_template}.}
\label{fig:result_zdcall_ppb}
\end{figure}
\begin{figure}[ht]
    \centering
    \includegraphics[width=0.75\linewidth]{sections/Results/fig/dijet/ZdcSpectrum_DJSample.pdf}
    \caption{Distribution of \EZDC\ with a dijet selection. The Inset shows a zoomed view of the single neutron peak. Please note that residual pile-up is not subtracted from this distribution. Refer to Section~\ref{ssec:pup_template} for information about the contribution of residual pile-up. 
}
    \label{fig:result_zdcSpectra_pPb}
\end{figure}

The correlation between the energy accumulated in the Pb-going ZDC and the same arm of the FCal is shown in Figure~\ref{fig:result_corr_mb_ppb}. The same plot in the bins of the analysis can be seen in Figure~\ref{fig:mb_corr}. Interestingly, one can notice a group of events characterized by $\leq$ 1 neutron in the ZDC and also associated to low energy deposit in the FCal. These events may have a specific physics contribution, for instance from diffractive processes, and will be further scrutinized later in the Note. 

Figure~\ref{fig:result_zdcCent_ppb} (\ref{fig:result_zdcNColl_ppb}) reports the \AZDC\ as a function of event centrality (average number of binary collisions), where the latter are calculated using standard Glauber MC techniques as described in \cite{Hill:2301540}. The evaluation of centrality-related quantities in \pxPb\ based on event activity has been found to not be robust due to experimental effects heavily dependent on the initial state kinematics \cite{ATLAS:2023zfx} and the analysis of the results in such fashion is merely shown to provide further information to the global modeling of these effects.

Figure~\ref{fig:mb_comp} reports normalized distributions of \EZDC\ and \ETPB\ in minimum bias events.  The axes and binning was chosen to directly compare to Figures~\ref{fig:unf_select_xp_zdc}~and~\ref{fig:unf_select_xp_fcal}. One can see significant differences in these distributions, as is expected due to the absence of a hard scattering selection in minimum bias events.

\begin{figure}[ht]\centering
\begin{subfigure}{.75\linewidth}
\includegraphics[width=\linewidth]{sections/Results/fig/minbias/cent_zdc_pb_bin_nom_HLT_mb_sbtrk_L1MBTS_1_templated.pdf}
\end{subfigure}
\caption{\AZDC\ as a function of centrality for MB Events in the \pxPb\ orientation data period. Centrality is classified using the event activity  selection described in \cite{Hill:2301540}. }
\label{fig:result_zdcCent_ppb}
\end{figure}

\begin{figure}[ht]\centering
\begin{subfigure}{.5\linewidth}
\includegraphics[width=\linewidth]{sections/Results/fig/minbias/npart_zdc_pb_bin_nom_HLT_mb_sbtrk_L1MBTS_1_templated.pdf}
\end{subfigure}
\caption{\AZDC\ as a function of \NPART\ for the \pxPb\ orientation.   \NPART\ is calculated using Table \ref{tab:cent_bin} for a fine selection from \cite{Hill:2301540}.}
\label{fig:result_zdcNColl_ppb}
\end{figure}



\begin{figure}[ht]\centering
\begin{subfigure}{.75\linewidth}
\includegraphics[width=\linewidth]{sections/Results/fig/minbias/corr_MinBias.pdf}
\end{subfigure}
\caption{Correlation plot between \EZDC\ and \ETPB\ for a minimum bias selection.}
\label{fig:mb_corr}
\end{figure}

\begin{figure}[ht]\centering
\begin{subfigure}{.75\linewidth}
\includegraphics[width=\linewidth,trim={0mm 3cm 0cm 0mm},clip]{sections/Results/fig/minbias/zdc_mb.pdf}
\end{subfigure}
\begin{subfigure}{.75\linewidth}
\includegraphics[width=\linewidth,trim={0mm 3cm 0cm 0mm},clip]{sections/Results/fig/minbias/fcal_mb.pdf}
\end{subfigure}
\caption{Normalized distributions of \EZDC\ (top) and \ETPB\ (bottom) in MB events.}
\label{fig:mb_comp}
\end{figure}

\clearpage





























































\subsection{Analysis of \ETP($x_p$)}
\label{sec:pgoing_fcal}
The analysis of the energy deposited in the FCal facing the proton beam provides significant information for the global modeling of the data, aiding also to further discriminate between different explanations proposed for the \xp-dependent variation of event activity in the Pb-going direction. A number of models based on energy conservation were proposed to explain the ATLAS inclusive jet data, see for instance \cite{Armesto:2015kwa,Bzdak:2014rca}. Despite strongly disfavored by dedicated measurements of event activity dependence on the hard-scattering carried out by ATLAS \cite{HION-2015-02}, the energy-conservation hypothesis is still considered valid. A recent example is offered by the JETSCAPE collaboration, that using the new hybrid X-SCAPE framework \cite{JETSCAPE:2024dgu} attempted to describe the ATLAS inclusive jet data and boil down the explanation of the event activity dependence on the hard scattering kinematics to soft-hard correlations originating from energy-momentum conservation.

A dedicated analysis on both \ETPB\ and \ETP\ energies as a function of \xp\ can provide additional input to test the validity of different hypothesis. A energy conservation-based picture would require the depletion in \ETPB\ with increasing \xp\ to be paired with a corresponding increase in \ETP\ in the same kinematic domain.

While the analysis of \ETPB(\xp) has been presented in the previous sections, with unfolding of experimental effects related to the jet detection, the analysis of \ETP(\xp) presents different implications. These can be already grasped by analyzing the distribution of \EATP(\xp) in dijet events selected for this analysis, presented in Figure~\ref{fig:raw_etp_xp}. 
\begin{figure}
    \centering
    \includegraphics[width=0.5\linewidth]{sections/Results/fig/FCal_p/NC_fcal_p_xp_DJSample.pdf}
    \caption{Analysis of \EATP\ as a function of the parton level proton initial state kinematics, \xp. Please note that no unfolding was applied to obtain theses results.   }
    \label{fig:raw_etp_xp}
\end{figure}

The average energy accumulated in the p-going forward calorimeter remains constant for all \xp\ $ < $ 0.1, while it starts to steeply rise above this value. This trend can be trivially explained by looking at the pseudorapidity coverage of dijets in this kinematic range, demonstrated in Figure~\ref{fig:eta1eta2_mc}. 

\begin{figure}
    \centering
    \includegraphics[width=0.85\linewidth]{sections/Results/fig/xp/truth_xp_grid.pdf}
    \caption{\xp\ distribution of dijet events as a function of the pseudorapidity of the leading ($\eta_1$) and sub-leading ($\eta_2$) jets, obtained from analysis of MC truth.}
    \label{fig:eta1eta2_mc}
\end{figure}

It appears clear how the increase of average energy accumulated in the p-going FCal coincides with one or both the analyzed jets falling into the acceptance of the calorimeter. A dedicated analysis of \ETP\ was implemented to remove the bias associated to the presence of the jets.  First, events were identified with leading and/or sub-leading R=0.4 jet(s) detected at $\eta \geq 2.8$, e.g., with at least one jet in the dijet that could have constituents within the acceptance of the p-going FCal. The calorimeter towers associated with these jets and residing in the p-going FCal were identified, tagged, and subtracted at level of the \ETP\ computation. The results of this exercise, carried out using $E_\mathrm{T}$ measured the EM scale for all the FCal cells, are shown in Figure~\ref{fig:fcal_wo_jets}. The jet un-subtracted distribution is also shown for direct assessment of the impact of the jet(s) in the balance of energy deposited in the FCal.

\begin{figure}
    \centering
    \includegraphics[width=0.5\linewidth]{sections/Results/fig/FCal_p/ratio_pFCalSubtraction.pdf}
    \caption{\EATP\ measured as a function of \xp\ without (azure) and with (red) subtraction of the jet contribution. A log-linear fit is applied to the region with \xp\ $>$ 0.1 to assess the impact of the jet contribution subtraction in the region where a steep rise is observed. The very modest residual contribution observed can be explained with out-of-cone radiation not accounted for when subtracting jets of R = 0.4. }
    \label{fig:fcal_wo_jets}
\end{figure}

The results of this exercise clearly demonstrate how the strong rise in \EATP\ in the \xp\ $>$ 0.1 region is a direct manifestation of the jet contribution in the p-going FCal in this kinematic domain. A log-linear fit is applied to both subtracted and un-subtracted distributions, and shows how the rise in the energy in the subtracted distribution largely disappear. The residual slope captured by the log-linear fit can be interpreted in terms of out-of-cone radiation for the R = 0.4 jets, which is expected to get more relevant the more the dijet pair is boosted towards forward rapidity, e.g., the higher the \xp. 

The results shown in Figure~\ref{fig:fcal_wo_jets} strongly oppose an interpretation of the results in Figure~\ref{fig:unfresult_avgfcal} in terms of energy-momentum conservation. The increase in \EATP\ can be attributed mainly to the presence of the jets, and once this is subtracted, the bulk of energy in the calorimeter appears to be compatible across different \xp\ regions. More importantly then, the decrease for \EATPB\ as a function \xp\ starts well before \xp $\sim$ 0.1, while there is no evidence of \xp\ dependence of \EATP\ in the same kinematic region. This observation is at odds with an explanation in terms of energy-momentum conservation, that would expect the energy in the two halves of the calorimeter to show an anti-correlated \xp-dependence. 

{\color{red}ToDo: The analysis team is working on a comparison with Pythia8 MC - the processing time required has not made this possible in this first iteration.}

\clearpage



\clearpage{}
\clearpage{}\section{Conclusion}
\label{sec:conclusion}



This note presents the first characterization of forward energy deposited in the ATLAS Forward Calorimeter and Zero-Degree Calorimeter in \pxPb\ collisions at \SNN~=~8.16~\TeV\ as a function of the estimated Bjorken-x of the proton, \xp, accessed in dijet events. 
Both \EZDC\ and \ETPB\ are shown to be sensitive to \xp, especially in the region \xp~$\gtrsim 5\cdot10^{-2}$, where previous ATLAS results from the same dataset \cite{ATLAS:2023zfx} have shown strong evidence of event activity biases that can be linked to color fluctuation effects related to the proton configuration~\cite{Alvioli:2017wou}. 
Further studies on the correlation between the two physics quantities as a function of \xp\ are also reported. In this case, the correlation is studied without unfolding of the experimental effects. The analysis of the `per neutron FCal energy' and the `FCal energy at zero neutrons’ unveiled correlations that change with \xp, which can be used to advance the modeling of event activity and nuclear breakup in \pxPb\ collisions characterized by the presence of a hard-scattering. 


This new data will aid in further understanding the proton configuration-related QCD-driven dynamics at work in \pxPb\ collisions. This novel type of analysis can have significant impact on the modeling of the nuclear breakup in \pxA\ collisions. Recently, the idea that color fluctuations can play a role in the evaporation of neutrons in the nuclear breakup has been proposed resolved photon collisions in UPCs \cite{Alvioli:2024cmd}. Such effects, related to the number of wounded nucleons, can be used to study the \xp-dependence of shadowing effects in UPCs and can affect the understanding of ZDC energy measured in future $e^-+A$ collisions at the future Electron-Ion Collider (EIC) \cite{Zheng:2014cha}. This measurement can provide a first look at the impact of color fluctuations in hadron-nucleus collisions, laying the ground for new theoretical advancements in the understanding of the nuclear breakup that will benefit both LHC and EIC physics programs. 
\clearpage{}

\printbibliography


\clearpage



\clearpage{}\section{Supporting Material} 


\ifinternalnote
\else \fi 





\subsection{Intrinsic resolution on event-level kinematics when using dijets}
\label{sec:xp_res}
In this analysis, the Bjorken-x of the initial state partons are reconstructed using the kinematics of the leading and sub-leading jet in the dijet pair, using the relation 
\begin{equation}
x_p = 
\frac{ p_{\TN{T},1}e^{y^\TN{CM}_{1}} 
+ 
p_{\TN{T},2}e^{y^\TN{CM}_{2}}}
{\sqrt{s_\mathrm{NN}}}
\label{eq:xp_aux}
\end{equation}
Accessing the parton level kinematics using jets involves approximations built-in to the definition of the probe, which has a finite precision due to the radius characterizing it. If, in a given event, particles are emitted outside of the cone considered in the analysis, set by the jet radius $R$ (= 0.4 in this analysis), the corresponding jet momentum will be intrinsically lower than the one of the parton initiating the jet. Such a discrepancy can be evaluated by using Monte Carlo generators (in this case, Pythia8) to study the accuracy of the initial state kinematics description achieved by using a certain jet radius. 

We have analyzed the accuracy of Equation~\ref{eq:xp_aux} by using the MC Pythia8 dijet simulation described in Section~\ref{sec:data}. Since the measurement presented in this note is unfolded for experimental effects introduced by the jet detection, the comparison carried out is at truth MC level, between the parton-level \xp\ variable (e.g. what is directly provided as \xp\ from Pythia8), indicated in the following as $x_p^\mathrm{pythia}$, and the truth jet-reconstructed \xp, referred to in the following as $x_p^\mathrm{nominal}$. 

Figure~\ref{fig:xp_res} demonstrates the distribution of the ratio $x_p^\mathrm{nominal}$/$x_p^\mathrm{pythia}$ in each of the \xp\ kinematic bins defined for the analysis. Each distribution has been fit with a Gaussian to $\pm1.2\sigma$ of the mean, found with a preliminary fit to the distribution. 

\begin{figure}[ht]
    \centering
    \includegraphics[width=0.95\linewidth]{sections/SupportMaterial/fig/xp_reso/xP_Resolution_Studies_Panel.pdf}
    \caption{Relative scale on $x_p^\mathrm{pythia}$ when the kinematic of the initial state is accessed via unfolded dijets using Equation~\ref{eq:xp_aux}.}
    \label{fig:xp_res}
\end{figure}

The dependence on $x_p^\mathrm{nominal}$ (e.g. the variable used in this analysis to study the forward energy production) of both mean and standard deviation of the ratio $x_p^\mathrm{nominal}$/$x_p^\mathrm{pythia}$ is shown in Figure~\ref{fig:xp_mean} and ~\ref{fig:xp_sig}, respectively. The $x_p^\mathrm{nominal}$ scale is in average 6-8\% below the parton-level value from the generator. This behavior can be explained in terms of radiation escaping the radius set for the jet reconstruction that, in average, leads to a lower than expected $x_p^\mathrm{nominal}$ value reconstructed. The relative resolution on $x_p^\mathrm{nominal}$ ranges from 10\% at the lowest values of $x_p^\mathrm{nominal}$ to $\sim$3\% in the highest $x_p^\mathrm{nominal}$ bin defined for this analysis. 

\begin{figure}[ht]
    \centering
    \includegraphics[width=0.75\linewidth]{sections/SupportMaterial/fig/pythia/xp_mu.pdf}
    \caption{Relative scale on $x_p^\mathrm{pythia}$ when the kinematic of the initial state is accessed via unfolded dijets using Equation~\ref{eq:xp_aux}.}
    \label{fig:xp_mean}
\end{figure}

\begin{figure}[ht]
    \centering
    \includegraphics[width=0.75\linewidth]{sections/SupportMaterial/fig/pythia/xp_sigma.pdf}
    \caption{Relative resolution on $x_p^\mathrm{pythia}$ when the kinematic of the initial state is accessed via unfolded dijets using Equation~\ref{eq:xp_aux}.}
    \label{fig:xp_sig}
\end{figure}









\clearpage
\subsection{Dependence of calorimetric energy on average interactions per bunch crossing}
\label{aux:mu_dep_averages}
The 2016 \pxPb\ dataset was taken at a relatively large $\mu$ for heavy ion collisions, around 0.1.  However, two of the first runs in the \pxPb\ orientation were taken at lower $\mu$ to slowly ramp the LHC up to peak intensities.  In order to determine if this measurement has any sensitivity to changing the number of interactions per bunch crossing, non-unfolded results of \AZDC\ and \EATPB\ are presented for several bins of $\mu$, see Figure~\ref{fig:support_mu_study}. Please note here that $\mu$ denotes a per lumiBlock quantity indicating the average interaction per bunch crossing measured in a given lumiBlock. The \AZDC\ plotted in this study includes residual pileup subtraction.

\begin{figure}[ht]\centering
\begin{subfigure}[c]{.49\linewidth}
\includegraphics[width=\linewidth]{sections/SupportMaterial/fig/muTemplateStudy/zdc_xp_mu_dependence.pdf}
\end{subfigure}
\begin{subfigure}[c]{.49\linewidth}
\includegraphics[width=\linewidth]{sections/SupportMaterial/fig/muTemplateStudy/fcal_xp_mu_dependence.pdf}
\end{subfigure}
\caption{\AZDC\ (left) and \EATPB\ (right) as a function of \xp\ for three different selections of average interactions per bunch crossing ($\mu$). The \EATPB\ is very resilient to $\mu$ selections.  The variation in \AZDC\ is sightly larger, but a striking dependence is not observed. The highest \xp\ point of \AZDC\ for the lowest $\mu$ is an outlier.  The amount of sampled luminosity in the low  $\mu$ runs is relatively low, and therefore the amount of jets sampled with a very large \xp\ is also low.  This leads to the large statistical variation observed here.}
\label{fig:support_mu_study}
\end{figure}
\clearpage


\subsection{Impact of residual pileup subtraction on \texorpdfstring{\AZDC}{<EZDC>}}
\label{aux:zdc_ref}
In Figure~\ref{fig:template_impact}, non-unfolded results for \AZDC\ are plotted before and after the residual pileup subtraction is performed.  The \AZDC\ values are shown to be very insensitive to the residual pileup subtraction.

\begin{figure}[ht]
    \centering
    \includegraphics[width=0.75\linewidth]{sections//SupportMaterial//fig//muTemplateStudy/zdc_xp_ratio_template_test.pdf}
    \caption{\AZDC as a function of \xp\ with and without the residual pileup subtraction via template fit.  The \AZDC\ is found to be very insensitive to the residual pileup subtraction, as the ratio of the unsubtracted over the subtracted is a fraction of a percent over unity. }
    \label{fig:template_impact}
\end{figure}
\clearpage

\subsection{Systematic Uncertainties on \texorpdfstring{\EZDC}{ezdc}}
\label{sec:zdc_panel_sys}

\begin{figure}[ht]\centering
\begin{subfigure}[c]{.49\linewidth}
\includegraphics[width=\linewidth]{sections/SupportMaterial/fig/zdc_syst/zdc_syst_xp[22].pdf}
\end{subfigure}
\begin{subfigure}[c]{.49\linewidth}
\includegraphics[width=\linewidth]{sections/SupportMaterial/fig/zdc_syst/zdc_syst_xp[21].pdf}
\end{subfigure}
\begin{subfigure}[c]{.49\linewidth}
\includegraphics[width=\linewidth]{sections/SupportMaterial/fig/zdc_syst/zdc_syst_xp[20].pdf}
\end{subfigure}
\begin{subfigure}[c]{.49\linewidth}
\includegraphics[width=\linewidth]{sections/SupportMaterial/fig/zdc_syst/zdc_syst_xp[19].pdf}
\end{subfigure}
\begin{subfigure}[c]{.49\linewidth}
\includegraphics[width=\linewidth]{sections/SupportMaterial/fig/zdc_syst/zdc_syst_xp[18].pdf}
\end{subfigure}
\begin{subfigure}[c]{.49\linewidth}
\includegraphics[width=\linewidth]{sections/SupportMaterial/fig/zdc_syst/zdc_syst_xp[17].pdf}
\end{subfigure}

\caption{Relative systematic uncertainties on \EZDC. Total Uncertainty (black dashed line) shown along with uncertainties for each systematic.  Panels correspond to measurements in different bins of \xp.  Cont. on next page.}
\label{fig:syst_zdc_panel_1}
\end{figure}

\begin{figure}[ht]\centering
\begin{subfigure}[c]{.49\linewidth}
\includegraphics[width=\linewidth]{sections/SupportMaterial/fig/zdc_syst/zdc_syst_xp[16].pdf}
\end{subfigure}
\begin{subfigure}[c]{.49\linewidth}
\includegraphics[width=\linewidth]{sections/SupportMaterial/fig/zdc_syst/zdc_syst_xp[15].pdf}
\end{subfigure}
\begin{subfigure}[c]{.49\linewidth}
\includegraphics[width=\linewidth]{sections/SupportMaterial/fig/zdc_syst/zdc_syst_xp[14].pdf}
\end{subfigure}
\begin{subfigure}[c]{.49\linewidth}
\includegraphics[width=\linewidth]{sections/SupportMaterial/fig/zdc_syst/zdc_syst_xp[13].pdf}
\end{subfigure}
\begin{subfigure}[c]{.49\linewidth}
\includegraphics[width=\linewidth]{sections/SupportMaterial/fig/zdc_syst/zdc_syst_xp[12].pdf}
\end{subfigure}
\begin{subfigure}[c]{.49\linewidth}
\includegraphics[width=\linewidth]{sections/SupportMaterial/fig/zdc_syst/zdc_syst_xp[11].pdf}
\end{subfigure}
\begin{subfigure}[c]{.49\linewidth}
\includegraphics[width=\linewidth]{sections/SupportMaterial/fig/zdc_syst/zdc_syst_xp[10].pdf}
\end{subfigure}
\begin{subfigure}[c]{.49\linewidth}
\includegraphics[width=\linewidth]{sections/SupportMaterial/fig/zdc_syst/zdc_syst_xp[9].pdf}
\end{subfigure}

\caption{Relative systematic uncertainties on \EZDC. Total Uncertainty (black dashed line) shown along with uncertainties for each systematic.  Panels correspond to measurements in different bins of \xp.  Cont. on next page.}
\label{fig:syst_zdc_panel_2}
\end{figure}

\begin{figure}[ht]\centering
\begin{subfigure}[c]{.49\linewidth}
\includegraphics[width=\linewidth]{sections/SupportMaterial/fig/zdc_syst/zdc_syst_xp[8].pdf}
\end{subfigure}
\begin{subfigure}[c]{.49\linewidth}
\includegraphics[width=\linewidth]{sections/SupportMaterial/fig/zdc_syst/zdc_syst_xp[7].pdf}
\end{subfigure}
\begin{subfigure}[c]{.49\linewidth}
\includegraphics[width=\linewidth]{sections/SupportMaterial/fig/zdc_syst/zdc_syst_xp[6].pdf}
\end{subfigure}
\begin{subfigure}[c]{.49\linewidth}
\includegraphics[width=\linewidth]{sections/SupportMaterial/fig/zdc_syst/zdc_syst_xp[5].pdf}
\end{subfigure}
\caption{Relative systematic uncertainties on \EZDC. Total Uncertainty (black dashed line) shown along with uncertainties for each systematic.  Panels correspond to measurements in different bins of \xp. }
\label{fig:syst_zdc_panel_3}
\end{figure}

\clearpage


\subsection{Systematic Uncertainties on \texorpdfstring{\ETPB}{etpb}}
\label{sec:fcal_panel_sys}

\begin{figure}[ht]\centering
\begin{subfigure}[c]{.49\linewidth}
\includegraphics[width=\linewidth]{sections/SupportMaterial/fig/fcal_syst/fcal_syst_xp[22].pdf}
\end{subfigure}
\begin{subfigure}[c]{.49\linewidth}
\includegraphics[width=\linewidth]{sections/SupportMaterial/fig/fcal_syst/fcal_syst_xp[21].pdf}
\end{subfigure}
\begin{subfigure}[c]{.49\linewidth}
\includegraphics[width=\linewidth]{sections/SupportMaterial/fig/fcal_syst/fcal_syst_xp[20].pdf}
\end{subfigure}
\begin{subfigure}[c]{.49\linewidth}
\includegraphics[width=\linewidth]{sections/SupportMaterial/fig/fcal_syst/fcal_syst_xp[19].pdf}
\end{subfigure}
\begin{subfigure}[c]{.49\linewidth}
\includegraphics[width=\linewidth]{sections/SupportMaterial/fig/fcal_syst/fcal_syst_xp[18].pdf}
\end{subfigure}
\begin{subfigure}[c]{.49\linewidth}
\includegraphics[width=\linewidth]{sections/SupportMaterial/fig/fcal_syst/fcal_syst_xp[17].pdf}
\end{subfigure}

\caption{Relative systematic uncertainties on \ETPB. Total Uncertainty (black dashed line) shown along with uncertainties for each systematic.  Panels correspond to measurements in different bins of \xp.  Cont. on next page.}
\label{fig:syst_fcal_panel_1}
\end{figure}

\begin{figure}[ht]\centering
\begin{subfigure}[c]{.49\linewidth}
\includegraphics[width=\linewidth]{sections/SupportMaterial/fig/fcal_syst/fcal_syst_xp[16].pdf}
\end{subfigure}
\begin{subfigure}[c]{.49\linewidth}
\includegraphics[width=\linewidth]{sections/SupportMaterial/fig/fcal_syst/fcal_syst_xp[15].pdf}
\end{subfigure}
\begin{subfigure}[c]{.49\linewidth}
\includegraphics[width=\linewidth]{sections/SupportMaterial/fig/fcal_syst/fcal_syst_xp[14].pdf}
\end{subfigure}
\begin{subfigure}[c]{.49\linewidth}
\includegraphics[width=\linewidth]{sections/SupportMaterial/fig/fcal_syst/fcal_syst_xp[13].pdf}
\end{subfigure}
\begin{subfigure}[c]{.49\linewidth}
\includegraphics[width=\linewidth]{sections/SupportMaterial/fig/fcal_syst/fcal_syst_xp[12].pdf}
\end{subfigure}
\begin{subfigure}[c]{.49\linewidth}
\includegraphics[width=\linewidth]{sections/SupportMaterial/fig/fcal_syst/fcal_syst_xp[11].pdf}
\end{subfigure}
\begin{subfigure}[c]{.49\linewidth}
\includegraphics[width=\linewidth]{sections/SupportMaterial/fig/fcal_syst/fcal_syst_xp[10].pdf}
\end{subfigure}
\begin{subfigure}[c]{.49\linewidth}
\includegraphics[width=\linewidth]{sections/SupportMaterial/fig/fcal_syst/fcal_syst_xp[9].pdf}
\end{subfigure}

\caption{Relative systematic uncertainties on \ETPB. Total Uncertainty (black dashed line) shown along with uncertainties for each systematic.  Panels correspond to measurements in different bins of \xp.  Cont. on next page.}
\label{fig:syst_fcal_panel_2}
\end{figure}

\begin{figure}[ht]\centering
\begin{subfigure}[c]{.49\linewidth}
\includegraphics[width=\linewidth]{sections/SupportMaterial/fig/fcal_syst/fcal_syst_xp[8].pdf}
\end{subfigure}
\begin{subfigure}[c]{.49\linewidth}
\includegraphics[width=\linewidth]{sections/SupportMaterial/fig/fcal_syst/fcal_syst_xp[7].pdf}
\end{subfigure}
\begin{subfigure}[c]{.49\linewidth}
\includegraphics[width=\linewidth]{sections/SupportMaterial/fig/fcal_syst/fcal_syst_xp[6].pdf}
\end{subfigure}
\begin{subfigure}[c]{.49\linewidth}
\includegraphics[width=\linewidth]{sections/SupportMaterial/fig/fcal_syst/fcal_syst_xp[5].pdf}
\end{subfigure}
\caption{Relative systematic uncertainties on \ETPB. Total Uncertainty (black dashed line) shown along with uncertainties for each systematic.  Panels correspond to measurements in different bins of \xp. }
\label{fig:syst_fcal_panel_3}
\end{figure}

\clearpage


\subsection{Results - proton going ZDC}
To better understand the correlations between \xp\ and the energy accumulated in the ZDC, the averages of the distributions presented in Figure~\ref{fig:support_panelplot_pPb_zdc_pside_xp} are presented in Figure~\ref{fig:result_zdc_xp_p}. One can notice how the distribution seems to be sensitive to \xp, in particular in the region \xp $\gtrsim 5\cdot10^{-2}$, corresponding to small proton configurations. The energy accumulated in the $p$-going side demonstrates an opposite trend, with a deferred turn-on around \xp$\gtrsim 3\cdot 10^{-1}$.  
Given the more peculiar physics involved in interpreting neutral energy on the $p$-going side, the timescale to finalize these studies has been set to match the final publication, and not the CONF note.
\begin{figure}[ht]
    \centering
    \includegraphics[width=1\linewidth]{sections/SupportMaterial/fig/pside/zdc_proton_xp_panel_DJSample.pdf}
    \caption{\pxPb\ orientation p-going ZDC spectrum.  Panels indicate which $x_p$ bin the associated dijet falls into. The $p$-going ZDC was calibrated using the first run after the reversal of the beam orientation, which was characterized by the same HV settings as the \pxPb\ orientation, as described in \cite{Hoppesch:2900504}. }
    \label{fig:support_panelplot_pPb_zdc_pside_xp}
\end{figure}



\begin{figure}[ht]\centering
\begin{subfigure}[c]{.49\linewidth}
\includegraphics[width=\linewidth]{sections/SupportMaterial/fig/pside/zdc_p_xp_DJSample.pdf}
\end{subfigure}
\begin{subfigure}[c]{.49\linewidth}
\includegraphics[width=\linewidth]{sections/SupportMaterial/fig/pside/zdc_p_xp_2d_DJSample.pdf}
\end{subfigure}
\caption{Average forward neutral energy in the proton-going ZDC as a function of \xp\ for the \pxPb\ orientation (left). Correlation between \xp\ and forward neutral proton-going ZDC energy, for the \pxPb\ orientation (right).  In the correlation plot, slices of \xp\ were weighted by the inverse of the sum of bins in each slice.}
\label{fig:result_zdc_xp_p}
\end{figure}




\clearpage




\subsection{Results - Both periods}
\label{auxsec:both_periods}
In this section, we report and compare results obtained for both periods. As a reminder, period C, also indicated as \Pbxp\ orientation, was not considered in the analysis due to the absence of the ZDC EM module on the Pb going side in favor of the LHCf detector, which was moved into garage position for this data taking, see Section~\ref{sec:mc_samples} for more details.   In the results comparing \AZDC\ for both running periods, the units of energy are divided by the nominal beam energy for neutrons, 2.56 TeV, and reported as number of neutrons.
The results reported in this section are not corrected for residual pile-up, and not unfolded for experimental effects. 
The 2d correlations, if not stated otherwise, display data only from the \pxPb\ orientation.

\begin{figure}[ht]
    \centering
    \includegraphics[width=0.75\linewidth]{sections/Results/fig/dijet/ZdcSpectrumComp_pPb_v_Pbp.pdf}
    \caption{Distribution of \EZDC\ with a Dijet Selection. The Inset shows a zoomed view of the single neutron peak. Plotted for both beam orientations \pxPb\ (blue) and \Pbxp\ (red). The \Pbxp\ distribution has a lower reach compared to the \pxPb\ distribution.  The one neutron peak sets the energy scale to 1n for both run periods, but the fraction of energy lost to punch through events is different between the sides.  This difference leads to the compression of the \Pbxp\ ZDC energy spectrum.}
    \label{fig:support_zdcSpectra_pPb_v_Pbp}
\end{figure}








\begin{figure}[ht]\centering
\begin{subfigure}[c]{.49\linewidth}
\includegraphics[width=\linewidth]{sections/Results/fig/dijet/zdc_comb_prof_xp_ratio_DJSample_comb_20Bins.pdf}
\end{subfigure}
\begin{subfigure}[c]{.49\linewidth}
\includegraphics[width=\linewidth]{sections/Results/fig/dijet/zdc_xp_2d_DJSample_comb_20Bins.pdf}
\end{subfigure}
\caption{\AZDC\ as a function of \xp\ for both periods (left). Correlation between \xp\ and forward neutral Pb-going ZDC energy, for the \pxPb\ orientation (right) .  In the correlation plot, slices of \xp\ were weighted by the inverse of the sum of bins in each slice.  Jet prescales applied when filling these histograms.}
\label{fig:support_zdc_xp}
\end{figure}


\begin{figure}[ht]
    \centering
    \includegraphics[width=0.75\linewidth]{sections/Results/fig/dijet/zdc_comb_prof_xp_ratio_p_DJSample_comb_20Bins.pdf}
    \caption{Average forward neutral energy in the proton-going ZDC as a function of \xp\ for both periods. The proton going ZDC has no energy calibration in the \Pbxp\ orientation, so the units are arbitrary; however, the energy was normalized to the \AZDCp\ in minimum bias \pxPb\ events.}
    \label{fig:support_zdc_xp_p_include_zeros}
\end{figure}


\begin{figure}[ht]\centering
\begin{subfigure}[c]{.49\linewidth}
\includegraphics[width=\linewidth]{sections/Results/fig/dijet/zdc_comb_prof_ex0_xp_ratio_p_DJSample_comb_20Bins.pdf}
\end{subfigure}
\begin{subfigure}[c]{.49\linewidth}
\includegraphics[width=\linewidth]{sections/SupportMaterial/fig/pside/zdc_p_xp_2d_DJSample_comb_20Bins.pdf}
\end{subfigure}
\caption{\textbf{Zero bin not included in the making of this plot. see Figure~\ref{fig:result_zdc_xp_p}.} Average forward neutral energy in the proton-going ZDC as a function of \xp\ for both periods (left). Correlation between \xp\ and forward neutral proton-going ZDC energy, for the \pxPb\ orientation (right).  In the correlation plot, slices of \xp\ were weighted by the inverse of the sum of bins in each slice. Jet prescales applied when filling these histograms. The proton going ZDC has no energy calibration so the units are arbitrary; however, the energy was normalized to the \AZDCp\ in minimum bias events. }
\label{fig:support_zdc_xp_p}
\end{figure}



\begin{figure}[ht]
    \centering
    \includegraphics[width=0.75\linewidth]{sections/Results/fig/dijet/zdc_comb_prof_XF_xp_ratio_DJSample_comb_20Bins.pdf}
    \caption{Average forward neutral energy in the Pb-going ZDC as a function of $x_F$ for both periods. Jet prescales applied when filling this histogram.}
    \label{fig:support_zdc_xf}
\end{figure}


\begin{figure}[ht]\centering
\begin{subfigure}[c]{.49\linewidth}
\includegraphics[width=\linewidth]{sections/SupportMaterial/fig/bothperiods/fcal_comb_prof_xp_ratio_DJSample_comb_20Bins.pdf}
\end{subfigure}
\begin{subfigure}[c]{.49\linewidth}
\includegraphics[width=\linewidth]{sections/SupportMaterial/fig/bothperiods/fcal_xp_2d_DJSample_comb_20Bins.pdf}
\end{subfigure}
\caption{Average forward energy in the Pb-going FCal as a function of $x_p$ for both periods (left). 
Correlation between \xp\ and forward Pb-going FCal energy, for the \pxPb\ orientation (right). The dijet system from which \xp\ is calculated is required to not deposit any energy in the Pb-going FCal. In the correlation plot, slices of \xp\ were weighted by the inverse of the sum of bins in each slice.
}
\label{fig:support_fcal_xp}
\end{figure}

\begin{figure}[ht]\centering
\begin{subfigure}[c]{.49\linewidth}
 \includegraphics[width=\linewidth]{sections/SupportMaterial/fig/bothperiods/fcal_comb_prof_xp_ratio_p_DJSample_comb_20Bins.pdf}
\end{subfigure}
\begin{subfigure}[c]{.49\linewidth}
\includegraphics[width=\linewidth]{sections/SupportMaterial/fig/bothperiods/fcal_p_xp_2d_DJSample_comb_20Bins.pdf}
\end{subfigure}
\caption{Average forward energy in the proton-going FCal as a function of \xp\ for both periods (left). Correlation between \xp\ and forward proton-going FCal energy, for the \pxPb\ orientation (right). The dijet system from which \xp\ is calculated can deposit energy in the proton-going FCal; therefore, the sharp rise in \EATP\ after $\sim$0.1 in \xp\ comes from a significant fraction of events with a large forward boost to the dijet system. In the correlation plot, slices of \xp\ were weighted by the inverse of the sum of bins in each slice. }
\label{fig:support_fcal_xp_p}
\end{figure}







\begin{figure}[ht]\centering
\begin{subfigure}[c]{.8\linewidth}
\includegraphics[width=\linewidth]{sections/Results/fig/dijet/zdc_sel_pPb_xp_DJSample_comb_20Bins.pdf}
\end{subfigure}
\begin{subfigure}[c]{.8\linewidth}
\includegraphics[width=\linewidth]{sections/Results/fig/dijet/zdc_sel_Pbp_xp_DJSample_comb_20Bins.pdf}
\end{subfigure}
\caption{Spectrum of forward neutral energy in the Pb-going ZDC (\EZDC) for selected bins of \xp.  Plotted for both the \pxPb\ (top) and \Pbxp\ (bottom) beam orientations.  Normalized to the weighted number of events in each slice of \xp.
}
\label{fig:support_zdc_spectrum_xp}
\end{figure}



















\clearpage
\subsection{ZDC distributions in different centrality selections}
\label{ssec:zdc_in_cent}
This appendix section includes the ZDC distributions (without removal of residual pile-up) in different FCal-based centrality selections. 
Figure~\ref{fig:zdc_dijet_A} (\ref{fig:zdc_dijet_C}) shows the comparison between the ZDC distributions in the 0-10\% and 60-90\% centrality selections for \pxPb\ (\Pbxp) data taking. From these distributions, it is evident the source of the trends observed in the \RCP\ results reported in Section~\ref{sec:results}.





\begin{figure}[ht]
    \centering
    \includegraphics[width=1\linewidth]{sections//Results//fig/dijet/zdc_cent0.pdf}
    \caption{\pxPb\ orientation Pb-going ZDC Spectrum for both central and peripheral centrality selections.  Panels indicate which kinematic bin the associated dijet falls into. }
    \label{fig:zdc_dijet_A}
\end{figure}

\begin{figure}[ht]
    \centering
    \includegraphics[width=1\linewidth]{sections//Results//fig/dijet/zdc_cent0_sidec.pdf}
    \caption{\Pbxp\ orientation Pb-going ZDC Spectrum for both central and peripheral centrality selections.  Panels indicate which kinematic bin the associated dijet falls into. }
    \label{fig:zdc_dijet_C}
\end{figure}

\clearpage

\subsection{Results as a function of \texorpdfstring{\xPb}{xPb}}

\textbf{These non-unfolded results are kept for documentation of the evolution of this analysis, but are not part of the final result}
\label{supp:xPb_results}
This section includes not unfolded results as a function of \xPb. Figure~\ref{fig:support_zdc_xPb} demonstrates the results for $\EAZDC$ as a function of \xPb. The excursion of the average of the ZDC energy is notably lower ($\sim$1/5) compared to what observed as a function of \xp. Figures~\ref{fig:support_zdc_xPb_p_include_zeros} and \ref{fig:support_zdc_xPb_p_exclude_zeros} show the results for $\EAZDCp$ for all dijet events and for those characterized by a non-null energy deposited in the $p$-going ZDC, respectively. Interestingly, also in this case requiring a non-null energy in the calorimeter sets the scale of the average energy to $\sim$ 1-1.2 TeV. Please note that these results were obtained using the ZDC distributions prior to the subtraction of residual pile-up and will be updated in the future. 

For completeness, Figure~\ref{fig:support_zdc_xPb} and \ref{fig:support_zdc_xPb_p_include_zeros} reports the average energy in the Pb and $p$-going FCal, respectively. The interpretation of the observed correlation between the results and \xPb\ is straightforward in both cases.

\begin{figure}[ht]
    \centering
    \includegraphics[width=0.75\linewidth]{sections/SupportMaterial/fig/xPb/zdc_Pb_xPb_DJSample.pdf}
    \caption{Average forward neutral energy in the Pb-going ZDC as a function of $x_{\mathrm{Pb}}$ for the \pxPb\ orientation.}
    \label{fig:support_zdc_xPb}
\end{figure}
\begin{figure}[ht]
    \centering
    \includegraphics[width=0.75\linewidth]{sections/SupportMaterial/fig/xPb/zdc_p_xPb_DJSample.pdf}
    \caption{Average forward neutral energy in the proton-going ZDC as a function of \xPb\ for the \pxPb\ orientation.}
    \label{fig:support_zdc_xPb_p_include_zeros}
\end{figure}
\begin{figure}[ht]
    \centering
    \includegraphics[width=0.75\linewidth]{sections/SupportMaterial/fig/xPb/zdc_p_ex0_xPb_DJSample.pdf}
    \caption{\textbf{Zero bin not included in the making of this plot. see Figure~\ref{fig:support_zdc_xPb_p_include_zeros}.} Average forward neutral energy in the proton-going ZDC as a function of \xPb\ for the \pxPb\ orientation.}
    \label{fig:support_zdc_xPb_p_exclude_zeros}
\end{figure}

\begin{figure}[ht]
    \centering
    \includegraphics[width=0.75\linewidth]{sections/SupportMaterial/fig/xPb/fcal_Pb_xPb_DJSample.pdf}
    \caption{Average forward energy in the Pb-going FCal as a function of \xPb\ for the \pxPb\ orientation. The dijet system from which \xPb\ is calculated is required to not deposit any Energy in the Pb-going FCal.}
    \label{fig:support_fcal_xPb}
\end{figure}

\begin{figure}[ht]
    \centering
    \includegraphics[width=0.75\linewidth]{sections/SupportMaterial/fig/xPb/fcal_p_xPb_DJSample.pdf}
    \caption{Average forward energy in the p-going FCal as a function of \xPb\ for the \pxPb\ orientation.}
    \label{fig:support_fcal_p_xPb}
\end{figure}
\clearpage



\subsection{Results as a function of \texorpdfstring{\xf}{xF}}
\label{supp:xf_res}

\textbf{These non-unfolded results are kept for documentation of the evolution of this analysis, but are not part of the final result}

This section includes not unfolded results as a function of \xf. 
Figure~\ref{fig:result_zdc_xf} reports the dependence of \EZDC\ on the Feynman variable, \xf. The trends observed in \xp\ are reflected also on the \xf\ dependence, even though at the extreme of the distribution the results are dominated by statistical uncertainties. Figure~\ref{fig:color_effect_zdc_fcal_xf} shows the comparison between \EZDC\ and \ETPB, both normalized to their maximum. The ratio between the two distributions clearly shows how the FCal (e.g. the products of the collisions) is much more sensitive to fluctuations in the proton configurations compared to the ZDC (e.g. the forward neutral energy, mostly originated from neutrons released from the nucleus in the nuclear breakup).

\begin{figure}[ht]
    \centering
    \includegraphics[width=0.75\linewidth]{sections/Results/fig/dijet/zdc_Pb_xf_temp_DJSample.pdf}
    \caption{Average forward neutral energy in the Pb-going ZDC as a function of $x_F$ for the \pxPb\ orientation. The ZDC residual pile-up contribution was removed to produce this plot. }
    \label{fig:result_zdc_xf}
\end{figure}
\begin{figure}[ht]
    \centering
    \includegraphics[width=.85\linewidth]{sections/Results/fig/dijet/coloreffect_percent_prof_XF_xp_ratio_DJSample_comb_20Bins.pdf}
    \caption{Average forward neutral energy in the Pb-going ZDC (after ZDC residual pile-up subtraction) as a function of \xf\ (blue). Average forward energy in the Pb-going FCal as a function of \xf\ (red). Both plots are normalized to their maximum value, in order to display both quantities on the same scale. The bottom plot shows the ratio of the average values of the above quantities. Plotted for runs in the \pxPb\ orientation.  }
    \label{fig:color_effect_zdc_fcal_xf}
\end{figure}
\clearpage



\subsection{Additional studies carried out for CONF note}

\subsubsection{\texorpdfstring{\RCP}{Rcp} Results}

In order to provide further information for comprehensive modeling of the data published in \cite{ATLAS:2023zfx}, an \RCP\ in the same centrality classes used in the former analysis was constructed using average ZDC energies. 
Please note that this part of the analysis is not yet updated with residual pile-up correction.
The mean \EZDC distributions for central (0--10\%) and peripheral (60--90\%) events, plotted as a function of \xp, are presented in the left and right portion of Figure~\ref{fig:result_zdc_comb_cent_peri}, respectively. The ZDC distributions in different centrality classes can be found in the Supporting Material, see Section~\ref{ssec:zdc_in_cent}. The \RCP\ result is presented, as a function of \xp, in Figure~\ref{fig:result_zdc_comb_rcp}. A suppression of about 20\% of this \RCPZDC\ is observed from low to high \xp. Interestingly, the suppression turn-on as a function of \xp\ is compatible to the standard \RCP\ suppression observed in \cite{ATLAS:2023zfx}. The same results, as a function of \xf, are reported in Figures~\ref{fig:result_zdc_comb_cent_peri_xf} and ~\ref{fig:result_zdc_comb_rcp_xf}. The limited statistical precision at the edges of the kinematic distribution limits the observations that could be made for \xf\ out of the range -0.4 $\leq$ \xf\ $\leq$ 0.5.
\begin{figure}[ht]\centering
\begin{subfigure}[c]{.49\linewidth}
\includegraphics[width=\linewidth]{sections/Results/fig/rcp/zdc_Pb_cent_xp_DJSample.pdf}
\end{subfigure}
\begin{subfigure}[c]{.49\linewidth}
\includegraphics[width=\linewidth]{sections/Results/fig/rcp/zdc_Pb_peri_xp_DJSample.pdf}
\end{subfigure}
\caption{Average forward neutral energy in the Pb-going ZDC as a function of \xp\ for a central selection (left) and a peripheral selection (right). Jet prescales applied when filling these histograms.
}
\label{fig:result_zdc_comb_cent_peri}
\end{figure}
\begin{figure}[ht]
    \centering
    \includegraphics[width=0.75\linewidth]{sections/Results/fig/rcp/zdc_Pb_rcp_xp_DJSample.pdf}
    \caption{Ratio central to peripheral of the average forward neutral energy in the Pb-going ZDC as a function of \xp.  The downwards scaling with \xp\ is similar to the effect seen in \cite{Longo:2860813}. Jet prescales applied when filling this histogram.}
    \label{fig:result_zdc_comb_rcp}
\end{figure}

\begin{figure}[ht]\centering
\begin{subfigure}[c]{.49\linewidth}
\includegraphics[width=\linewidth]{sections/Results/fig/rcp/zdc_Pb_cent_xf_DJSample.pdf}
\end{subfigure}
\begin{subfigure}[c]{.49\linewidth}
\includegraphics[width=\linewidth]{sections/Results/fig/rcp/zdc_Pb_peri_xf_DJSample.pdf}
\end{subfigure}
\caption{Average forward neutral energy in the Pb-going ZDC as a function of \xf\ for a central selection (left) and a peripheral selection (right). Jet prescales applied when filling these histograms.
}
\label{fig:result_zdc_comb_cent_peri_xf}
\end{figure}
\begin{figure}[ht]
    \centering
    \includegraphics[width=0.75\linewidth]{sections/Results/fig/rcp/zdc_Pb_rcp_xf_DJSample.pdf}
    \caption{Ratio central to peripheral of the average forward neutral energy in the Pb-going ZDC as a function of \xf.  Jet prescales applied when filling this histogram.}
    \label{fig:result_zdc_comb_rcp_xf}
\end{figure}
\clearpage{}


\end{document}
